% Copyright 2026 Christophe Jorssen <christophe.jorssen@gmail.com>.
% This file is part of luapourbaix.
% Distributed under the LPPL 1.3c or, at your option, any later version.
% Maintained by Christophe Jorssen. See COPYING.md, LICENSE and MANIFEST.md.

\documentclass[a4paper,11pt]{ltxdoc}
\usepackage{fontspec}
\usepackage[english]{babel}
\usepackage[margin=25mm]{geometry}
\usepackage{makeidx,calc,amsmath,longtable,array}
\usepackage{pgfplots}
\pgfplotsset{compat=1.18}
\usepgfplotslibrary{pourbaix}
\usepackage[hyperindex=false]{hyperref}
\usepackage{pgfmanual}
\input{pgfmanual-en-macros.tex}
\makeindex
\tracinglostchars=3
\hypersetup{pdftitle={The luapourbaix library},pdfauthor={Christophe Jorssen}}
\title{The \texttt{luapourbaix} Library\\
  \large Limiting Potential--pH Diagrams with PGFPlots}
\author{Christophe Jorssen\\
 \href{mailto:christophe.jorssen@gmail.com}{\texttt{christophe.jorssen@gmail.com}}}
\date{Version 1.0 --- 2026-10-08}

\begin{document}
\maketitle
\begin{abstract}
\texttt{luapourbaix} constructs limiting potential--pH (Pourbaix) diagrams
using ordinary PGFPlots commands. Declare a chemical system, check its
thermodynamic data, calculate candidate boundaries and finalize the retained
domains. Customize final plots, diagnostic overlays, water stability lines,
TikZ species maps and automatically composed reports. The generic library runs
with plain LuaTeX, LuaLaTeX and ConTeXt MkIV. It uses the ideal dilute limiting
model described in this manual; it is not a full aqueous speciation solver.
\end{abstract}
\begin{center}
\fbox{\begin{minipage}{.9\linewidth}
\textbf{AI assistance.} This package was ``vibe coded'' with the
assistance of \textbf{ChatGPT 6 Astra}, under the direction of Christophe Jorssen.
The author and current maintainer is Christophe Jorssen. Source provenance and
the limitations of the thermodynamic model are documented in this manual.
\end{minipage}}
\end{center}
\noindent Copyright 2026 Christophe Jorssen. This work is distributed under the
\href{https://www.latex-project.org/lppl/}{LaTeX Project Public License},
version 1.3c or (at your option) any later version. Maintenance status:
\texttt{maintained}. See \texttt{COPYING.md} and the unmodified license text in
\texttt{LICENSE}. The work's files are listed in \texttt{MANIFEST.md}.

\noindent Project repository:
\url{https://github.com/cjorssen/luapourbaix}.


\tableofcontents
\clearpage


\part{Getting started}
\section{Installation and formats}
\subsection{Replacement of the unsupported pourbaix package}
The earlier \href{https://github.com/cjorssen/pourbaix}{\texttt{pourbaix} package}
is no longer maintained or supported. \textbf{Use luapourbaix for new work and
migrate existing documents to the interface described in this manual.}
The original \texttt{pourbaix} was published only on GitHub; it was never
distributed through CTAN. Its repository and old manuals document the unsupported
implementation, not this library.

Migration requires adapting declarations and plotting commands, not merely
changing a package name. Load this implementation with
\verb|\usepgfplotslibrary{pourbaix}| after PGFPlots, rather than
\verb|\usepackage{pourbaix}|. The library identifier \texttt{pourbaix}, its key
family and the \texttt{pourbaix} axis environment are part of luapourbaix;
they do not imply continued support for the legacy package.

\subsection{Files, dependencies and loading}
The runtime files must be on the engine's search path. The repository uses
\path{tex/generic/luapourbaix/}; the CTAN source archive places the same files
in \texttt{tex/}. Copy them to \path{tex/generic/luapourbaix/} in a personal
TDS tree. Refresh the
TeX filename database if required by that tree. The generic Lua modules must
remain alongside the PGFPlots library loader.

Only LuaTeX, PGF/TikZ and PGFPlots are runtime dependencies. No network access,
external executable or shell escape is used by the library. Building and testing
the sources additionally uses Python~3.8 or later, Make, MakeIndex and the TeX
distribution's PGF documentation tools. The build runner sets local search paths
and caches without altering the global TeX installation.

For LuaLaTeX, load the library after PGFPlots:
\begin{verbatim}
\usepackage{pgfplots}
\pgfplotsset{compat=1.18}
\usepgfplotslibrary{pourbaix}
\end{verbatim}
For plain LuaTeX:
\begin{verbatim}
\input pgfplots.tex
\pgfplotsset{compat=1.18}
\usepgfplotslibrary{pourbaix}
\end{verbatim}
For ConTeXt MkIV, compiled using LuaTeX:
\begin{verbatim}
\usemodule[pgfplots]
\pgfplotsset{compat=1.18}
\usepgfplotslibrary[pourbaix]
\end{verbatim}
The declaration syntax is identical in all three formats. The library does not
select a global compatibility level, font, axis range or coordinate transformation.
Engines without \verb|\directlua| are rejected. Other Lua-capable engines have not
been qualified merely by possessing a Lua interpreter.

\section{Tutorial: An Iron Exercise for Maxime}
\label{sec:tutorial-iron}
Maxime is preparing a worksheet on aqueous corrosion. He wants his students to
identify the domains of five iron species, derive the boundary equations, and
explain whether metallic iron can coexist with water at equilibrium. His plotting
concentration is $c_t=0.01\,\mathrm{mol\,L^{-1}}$. Drawing the lines by hand would
be a useful exercise for the students, but it need not be another evening's work
for their teacher.

\subsection{A diagram before a table of constants}
Maxime chooses Fe(s), Fe$^{2+}$, Fe$^{3+}$, Fe(OH)$_2$(s) and Fe(OH)$_3$(s), at
298.15 K. These species are covered by the offline catalogue. He does not need
to look up and type their standard potentials and solubility products: an empty
\texttt{constants} list lets the library supply the missing independent data.
Nothing is downloaded while his document compiles.

He first loads PGFPlots and the library, as shown in the installation section.
Then he puts this complete picture in his document. The numbers use volts,
kelvins and mol/L; \texttt{.01} is the requested concentration. Maxime also
wants characteristic tick values to help his students read the diagram and
answer the boundary-equation questions. He enables extra ticks and places
their labels above and to the right, away from the regular graduations.

\begin{codeexample}[width=10cm]
\begin{tikzpicture}
\begin{pourbaix}[
 pH min=0,pH max=14,
 build={name=MaximeFe,element=Fe,
   temperature=298.15,concentration=.01,
   species={Fe(s),Fe^2+,Fe^3+,Fe(OH)2(s),Fe(OH)3(s)},
   boundary convention=equal element concentrations,
   constants={}},
 xmin=0,xmax=14,ymin=-1.2,ymax=1.5,width=10cm,height=7cm,
 extra x tick style={xtick pos=top,xticklabel pos=upper,
   xticklabel={\pgfmathprintnumber[fixed,precision=3]{\tick}}},
 extra y tick style={ytick pos=right,yticklabel pos=right,
   yticklabel={\pgfmathprintnumber[fixed,precision=3]{\tick}}}]
 \addpourbaix[domain palette=tab10 light,extra ticks=all,
   domain label mode=letters,
   final domain label style={fill=white,inner sep=1pt}]{MaximeFe};
 \addpourbaixwater[prepare={diagram=MaximeFe},
   domain fills=true,domain labels=false,
   domain/H2={style={fill=none}},
   domain/O2={style={fill=none}},
   domain/H2O={style={fill=none,pattern=north east lines,
     pattern color=black,opacity=.35}}];
\end{pourbaix}
\end{tikzpicture}
\end{codeexample}

The \texttt{build} key declares the chemistry, checks and completes the
thermodynamic data, calculates the provisional equations and constructs the
retained domains. It performs the four actions \texttt{declare}, \texttt{check},
\texttt{calculate} and \texttt{finalize}; errors such as missing constants are
reported normally. Maxime places the pH bounds before \texttt{build}, since keys
execute in written order. The two drawing commands
then produce the iron diagram and the water boundaries. Each drawing command
ends with a semicolon, as does an ordinary \cs{addplot} command. Maxime has written no
boundary equation and no polygon coordinates.

The concentration is a total concentration in iron, not the concentration of
every ion at every point. On each limiting boundary, species outside the pair
are neglected. Maxime explicitly chooses equal element concentrations for an
aqueous--aqueous boundary; here every species contains one iron atom, so equal
species concentrations would give the same result. The declaration does not
need the plot's potential limits; those belong to the view.

The plot option \texttt{extra ticks=all} adds the pH values of the retained
vertical boundaries and the intercepts of retained redox boundaries that reach
pH zero. Here the extra x ticks are about 2.297 and 7.393; the extra y ticks are
about $-0.499$ and $0.771$ V. These numbers come from the calculated diagram,
not from values Maxime typed into the axis. Their three-decimal labels do not
round the underlying calculation. Oblique boundaries that start at positive
pH contribute no extrapolated intercepts.

He can instead request a subset with \texttt{extra x ticks=\{a\}} and
\texttt{extra y ticks=\{1,3\}}, or remove these hints with
\texttt{extra ticks=none}. The unprefixed \texttt{extra x tick style} and
\texttt{extra y tick style} in the environment options are ordinary PGFPlots
keys, forwarded to the underlying axis. They change presentation only.

\subsection{Giving the students questions, not the answers}
\texttt{domain label mode=letters} replaces the species names with mathematical
uppercase letters. For this declaration, the answer key is A: Fe(s), B:
Fe$^{2+}$, C: Fe$^{3+}$, D: Fe(OH)$_2$(s), E: Fe(OH)$_3$(s). Boundary identifiers
remain numbers for redox boundaries and lowercase letters for vertical ones.
To print a corrected diagram, Maxime changes the label mode to \texttt{species}.
For a black-and-white worksheet, he changes the palette to \texttt{monochrome};
adjacent domains then have different black TikZ patterns.

Here is the worksheet he can put beside the figure:
\begin{enumerate}
\item Assign the five species to A--E. Distinguish predominance of a dissolved
  species from existence of a solid.
\item Write the balanced reaction for each numbered or lettered boundary.
  Use the limiting element balance and the appropriate activities to obtain
  its equation. Explain why the two vertical boundaries do not depend on $E$.
\item Identify the hydrogen and oxygen lines. Does the Fe(s) domain intersect
  the water-stability region? What does this say about oxidation of iron, and
  what does it not say about the corrosion rate?
\end{enumerate}
For numerical work he supplies the data actually used by this catalogue:
$E^\circ(\mathrm{Fe^{3+}/Fe^{2+}})=0.771$ V,
$E^\circ(\mathrm{Fe^{2+}/Fe})=-0.440$ V,
$pK_s(\mathrm{Fe(OH)_2})\simeq15.2131$,
$pK_s(\mathrm{Fe(OH)_3})=37.109$, and $pK_w=14$.
The phase and source conventions are those documented in the data reference;
these rounded teaching values are not a replacement for that provenance.

\subsection{Checking the answer sheet}
\label{sec:maxime-answer-sheet}
Let $s=RT\ln(10)/F=0.05915935$ V at 298.15 K. For example, on the Fe$^{2+}$/Fe
boundary the incipient solid contributes negligibly to the element balance,
so $[\mathrm{Fe^{2+}}]=c_t$ and
\[
 E=-0.440+\frac{s}{2}\log_{10}(c_t/c^\circ)
   \simeq-0.499159\ \mathrm V.
\]
On Fe$^{3+}$/Fe$^{2+}$ both concentrations equal $c_t/2$; their ratio is one and
$E=0.771$ V. For Fe(OH)$_m$(s), precipitation gives
$K_s=(c_t/c^\circ)a(\mathrm{OH^-})^m$, hence
\[
 \mathrm{pH}=pK_w-\frac{pK_s+\log_{10}(c_t/c^\circ)}{m}.
\]
The complete numerical answer sheet, rounded only for display, is:
\begin{center}
\begin{tabular}{cll}
Boundary & Equation ($E$ in V) & pH support\\\hline
1 & $E=-0.499159$ & $0\leq\mathrm{pH}\leq7.39343$\\
2 & $E=-0.0591593\,\mathrm{pH}-0.0617691$ & $7.39343\leq\mathrm{pH}\leq14$\\
3 & $E=0.771$ & $0\leq\mathrm{pH}\leq2.297$\\
4 & $E=-0.177478\,\mathrm{pH}+1.17867$ & $2.297\leq\mathrm{pH}\leq7.39343$\\
5 & $E=-0.0591593\,\mathrm{pH}+0.303887$ & $7.39343\leq\mathrm{pH}\leq14$\\
a & $\mathrm{pH}=7.39343$ & retained vertical segment\\
b & $\mathrm{pH}=2.297$ & retained vertical segment
\end{tabular}
\end{center}
The data keys and the report command in the reference part provide these same
coefficients and the constants used by a diagram, without transcribing a table.

\subsection{Letting the package write the answer sheet}
\cs{pgfpourbaixset} reads keys relative to \texttt{/pgfplots/pourbaix},
so Maxime can omit that prefix in settings and data queries. It respects the
current TeX scope and executes keys in their written order.

Maxime need not copy the coefficients from a printed diagram. The native PGF
handler \texttt{.get} copies a stored value into a TeX macro. The following
self-contained example builds the same diagram and reads its final boundaries.
The loop produces the five redox equations; the last two nodes give the vertical
pH values. Rounding is applied only when printing.

\begin{codeexample}[width=0pt]
\pgfpourbaixset{pH min=0,pH max=14,
 build={name=MaximeAnswers,element=Fe,
  species={Fe(s),Fe^2+,Fe^3+,Fe(OH)2(s),Fe(OH)3(s)},
  temperature=298.15,concentration=.01,
  boundary convention=equal element concentrations,constants={}}}
\pgfplotsgetpourbaix{MaximeAnswers}\FeHandle
\begin{tikzpicture}
\foreach \id in {1,...,5}{
 \pgfpourbaixset{data/\FeHandle/final/\id/.cd,
   a/.get=\slope,b/.get=\intercept,expression/.get=\formula}
 \node[anchor=west] at (0,-.5*\id)
  {$\id:\quad E=(\pgfmathprintnumber[precision=5]{\slope})\,
    \mathrm{pH}+(\pgfmathprintnumber[precision=5]{\intercept})$};
}
\foreach \id/\row in {a/6,b/7}{
 \pgfpourbaixset{data/\FeHandle/final/\id/pH/.get=\value}
 \node[anchor=west] at (0,-.5*\row)
  {$\id:\quad\mathrm{pH}=\pgfmathprintnumber[precision=5]{\value}$};
}
\end{tikzpicture}
\end{codeexample}

\texttt{expression} stores the PGFMath formula in \texttt{x}. Use it inside
its scope with:
\begin{verbatim}
\addplot {\formula};
\end{verbatim}
Restrict this plot to the retained support; do not extrapolate it over the
whole axis.
For a given final identifier, read \texttt{segment 1} to obtain the first segment
index, then \texttt{selection/segment N/pH min} and \texttt{pH max} for its limits.
A pair can have multiple disjoint segments. The reference lists every data key.

The handle and values are TeX-scoped. To reuse the initial diagram, place the
queries inside its \texttt{pourbaix} environment before closing it. To generate
an answer sheet outside the picture, build outside the picture as above.
Reading a nonexistent key does not calculate a missing value: retrieve only
fields appropriate to the boundary type, after a successful build.

\subsection{What the water overlay tells Maxime}
Maxime wants the water-stability region to remain visible over the colored iron
regions. In the first example, he draws the iron diagram first, then enables
\texttt{domain fills=true} for the water overlay. The addressed key
\texttt{domain/H2O} assigns diagonal TikZ hatching to just the central region
between the hydrogen and oxygen lines. The styles of \texttt{domain/H2} and
\texttt{domain/O2} use \texttt{fill=none}, leaving the outer regions unfilled
by the water overlay.

The pattern has a transparent background: the iron colors remain visible between
its strokes. \texttt{opacity=.35} softens the hatch strokes; use 1 for
fully opaque strokes, still with a transparent background. The pattern is the
standard TikZ \texttt{north east lines}; the library already loads
\texttt{patterns}. These water styles do not change the iron palette.
Water-domain labels are hidden in this worksheet, while the two dashed water
boundaries remain visible. The hatching marks the equilibrium water region; it
does not identify an additional iron species.

Recall the Nernst factor introduced in Section~\ref{sec:maxime-answer-sheet}:
$s=RT\ln(10)/F\simeq0.05915935\,\mathrm V$ at 298.15 K.
With the overlay's independent default gas pressures of 1 bar and $E$ in volts,
$E_{\mathrm{H_2}}=-s\,\mathrm{pH}$ and
$E_{\mathrm{O_2}}=1.229-s\,\mathrm{pH}$.
Over the displayed interval, Fe(s) lies entirely below the hydrogen line.
There is therefore no common equilibrium domain of metallic iron and water in
this model. Oxidation of iron coupled to hydrogen formation is thermodynamically
favored: in acid the products can be Fe$^{2+}$ and H$_2$, while the alkaline
limiting construction allows Fe(OH)$_2$(s) and H$_2$. Dissolved oxygen provides
another oxidant, represented by the upper water line.

Maxime insists on the last sentence of the exercise. A thermodynamic driving
force is not a corrosion rate. Hydroxide existence domains do not by themselves
prove that a protective film forms, remains adherent, or passivates the metal.
The diagram supplies equilibrium comparisons, not kinetics or film morphology.

\subsection{Using the classroom approximation of 0.06 V}
For another worksheet, Maxime wants to use $s=0.06$ V at $25\,{}^\circ$C.
He adds \texttt{nernst factor=0.06} inside \texttt{build} (or \texttt{declare}).
Temperature remains 298.15 K: this is an explicit teaching approximation, not
a change of temperature. The preceding figures and answer sheet use the default
\texttt{auto}, namely $RT\ln(10)/F$.

\begin{codeexample}[width=0pt]
\pgfpourbaixset{pH min=0,pH max=14,
 build={name=RoundedFe,element=Fe,
  temperature=298.15,nernst factor=0.06,concentration=.01,
  species={Fe(s),Fe^2+,Fe^3+,Fe(OH)2(s),Fe(OH)3(s)},
  boundary convention=equal element concentrations,constants={}},
 water/prepare={diagram=RoundedFe}}
\pgfplotsgetpourbaix{RoundedFe}\RoundedHandle
\pgfpourbaixset{data/\RoundedHandle/.cd,
 nernst factor/.get=\factor,final/1/b/.get=\metalPotential}
\pgfpourbaixset{water/data/a/.get=\waterSlope}
\begin{tikzpicture}
\node[align=left] {
 $s=\pgfmathprintnumber[fixed,precision=2]{\factor}\,\mathrm V$\\
 Fe$^{2+}$/Fe: $E=\pgfmathprintnumber[fixed,precision=3,fixed zerofill]
   {\metalPotential}\,\mathrm V$\\
 Hydrogen line: $E=(\pgfmathprintnumber[fixed,precision=2]{\waterSlope})\,\mathrm{pH}$ (V)};
\end{tikzpicture}
\end{codeexample}

The metal boundary is now $E=-0.440+(0.06/2)\log_{10}(0.01)=-0.500$ V.
A water overlay prepared with \texttt{diagram=RoundedFe} inherits the same
factor, so its hydrogen line is $E=-0.06\,\mathrm{pH}$ at 1 bar. The usual
\verb|\addpourbaixwater[prepare={diagram=RoundedFe}];| uses this inheritance too.

This setting affects calculation, not just printed precision. The library uses
the chosen factor consistently in Nernst equations and in the conversion
$\log_{10}K=nE^\circ/s$. Consequently, derived constants, intersections and even
some derived precipitation values can change. Supplied $E^\circ$ and $K$ values
are not rewritten. A redundant dataset consistent with the exact factor may
be inconsistent with the rounded one; the normal consistency error still
applies. Rebuild the diagram to change this approximation; already calculated
snapshots retain their own factor. Use \texttt{nernst factor=auto} to return to
the temperature-derived value.

\section{Tutorial: Sophie Makes Iodine Disproportionate}
\label{sec:tutorial-iodine}
Sophie wants a different lesson. Her students know that a boundary separates
two domains, but they have not yet seen why an intermediate oxidation state can
lose its domain altogether. She chooses I$^-$, I$_2$(aq) and IO$_3^-$, with iodine
oxidation numbers $-1$, $0$ and $+5$. She will first keep all three pair equations,
then let the class see what survives.

\subsection{Supplying the two missing potentials}
The offline catalogue has no iodine entries. Unlike Maxime, Sophie supplies two
independent standard potentials. She uses the iodine data in
\href{https://sci.rmutp.ac.th/woravith/?page_id=1786}{Woravith's university teaching
compilation, \emph{Reduction Standard Potential}}: 0.620 V for dissolved
I$_2$/I$^-$ and 0.269 V for the alkaline IO$_3^-$/I$^-$ reaction. She works at
298.15 K with $pK_w=14$, potentials versus SHE and unit water activity.
The balanced half-reactions are
\begin{align*}
 \mathrm{I_2(aq)}+2e^-&=2\mathrm{I^-},\\
 \mathrm{IO_3^-}+3\mathrm{H_2O}+6e^-&=
   \mathrm{I^-}+6\mathrm{OH^-}.
\end{align*}
The water coefficient in the alkaline reaction is three, as required by atom
balance. The package's pair-potential input uses H$^+$ and H$_2$O, so Sophie
converts the alkaline value before entering it:
\[
 \mathrm{IO_3^-}+6\mathrm{H^+}+6e^-=
 \mathrm{I^-}+3\mathrm{H_2O},\qquad
 E^\circ_{\mathrm{acid}}=0.269+s\,pK_w=1.0972309\ \mathrm V.
\]
She must not substitute the approximately 0.535 V value for I$_2$(s)/I$^-$:
her diagram contains \emph{aqueous} iodine, whose activity is concentration based.
These are explicit teaching inputs, not new catalogue entries. The conversion
retains more digits than the source data warrant only to make the calculation
reproducible.

\subsection{A convention that now matters}
Sophie sets $c_t=10^{-3}\,\mathrm{mol\,L^{-1}}$ in iodine atoms:
\[
 c_t=[\mathrm{I^-}]+2[\mathrm{I_2(aq)}]+[\mathrm{IO_3^-}].
\]
At I$_2$/I$^-$, equal \emph{species} concentrations mean
$[\mathrm{I_2}]=[\mathrm{I^-}]=c_t/3$. Equal \emph{element} concentrations mean
$2[\mathrm{I_2}]=[\mathrm{I^-}]$, hence
$[\mathrm{I_2}]=c_t/4$ and $[\mathrm{I^-}]=c_t/2$.
The same distinction applies to I$_2$/IO$_3^-$; for IO$_3^-$/I$^-$ the conventions
coincide. The library requires Sophie to state her choice. She plots both choices
side by side, without changing the selected species or the input potentials.
She adds characteristic ticks and a grid so the class can read the intercepts
and the pH at which the iodine domain disappears.

\begin{codeexample}[width=14cm]
\foreach \convention in
 {equal element concentrations,equal concentrations}{
 \pgfpourbaixset{declare={name=SophieI,element=I,
   species={IO3^-,I^-,I2(aq)},temperature=298.15,concentration=.001,
   boundary convention/.expanded=\convention,
   thermodynamic database=none,
   constants={
     standard potential/I2(aq)|I^-={value=.620},
     standard potential/IO3^-|I^-={value=1.0972309}}},
   check=SophieI,calculate=SophieI,
   pH min=0,pH max=10,finalize=SophieI}
 \pgfplotsgetpourbaix{SophieI}\IHandle
 \pgfpourbaixset{data/\IHandle/selection/event 1/.cd,
   pH min/.get=\iodineThreshold}
 \begin{tikzpicture}
 \begin{axis}[pourbaix/axis,width=6cm,height=6cm,font=\small,
   xmin=0,xmax=10,ymin=.3,ymax=1.3,title={\convention},
   grid=both,minor tick num=1,
   major grid style={gray!45},minor grid style={gray!20},
   extra x ticks/.expanded=\iodineThreshold,
   extra x tick style={grid=major,xtick pos=top,xticklabel pos=upper,
     xticklabel={\pgfmathprintnumber[fixed,precision=3]{\tick}}},
   extra y tick style={grid=major,ytick pos=right,yticklabel pos=right,
     yticklabel={\pgfmathprintnumber[fixed,precision=3]{\tick}}}]
  \addpourbaix[domain palette=tab10 light,extra ticks=all,
    show provisional boundaries=true,
    provisional boundary style={red!70!black,dashed},
    boundary label style={pos=.25},
    final domain label style={fill=white,inner sep=1pt}]{SophieI};
 \end{axis}
 \end{tikzpicture}
}
\end{codeexample}

The grid uses native PGFPlots keys: \texttt{grid=both}, \texttt{minor tick num},
\texttt{major grid style} and \texttt{minor grid style}. The library's plots
participate in the ordinary axis machinery, so tick formatting, titles, grid
styles and other applicable PGFPlots options remain available. The linear
potential--pH coordinate requirements still apply.

\texttt{extra ticks=all} supplies the two retained intercepts at pH zero.
This iodine set has no vertical boundary, so it supplies no automatic x tick.
Sophie obtains the start of the first elimination event with \texttt{.get} and
passes it to the native \texttt{extra x ticks} key. This is the disappearance
threshold for this three-species example, not an additional vertical boundary.
The value is read separately for each convention inside the loop, yielding
approximately 6.567 and 6.328. The extra-tick styles also draw grid lines at these
special coordinates. Reading the event leaves the scientific snapshot unchanged.

\texttt{thermodynamic database=none} makes the supplied data the complete input
set. Checking obtains the missing iodate/iodine potential by adding reaction free
energies. In this normalization,
\[
 6E^\circ(\mathrm{IO_3^-/I^-})=
 5E^\circ(\mathrm{IO_3^-/I_2})+E^\circ(\mathrm{I_2/I^-}),
\]
so the missing potential is approximately 1.192677 V. An unweighted average of
potentials would not give the correct result.

\subsection{Watching the middle domain disappear}
The dashed red lines are the provisional pair equations, including their
extensions beyond retained supports. Solid black lines are the final boundaries.
For equal element concentrations, the lower and upper iodine limits are
\begin{align*}
 E_{\mathrm{I_2/I^-}}&=0.708739\ \mathrm V,\\
 E_{\mathrm{IO_3^-/I_2}}&=1.174929-0.0709912\,\mathrm{pH}\quad(\mathrm V).
\end{align*}
The I$_2$ domain must lie above the first and below the second. They meet at
$\mathrm{pH}\simeq6.567$. Beyond that point there is no common interval of
potential: iodine disproportionates according to
\[
 3\mathrm{I_2(aq)}+3\mathrm{H_2O}=
 5\mathrm{I^-}+\mathrm{IO_3^-}+6\mathrm{H^+}.
\]
Finalization replaces the incompatible limits with the direct IO$_3^-$/I$^-$
boundary, $E=1.097231-0.0591593\,\mathrm{pH}$ in volts. The existing identifiers
1 and 3 terminate; identifier 2 continues to higher pH. With equal species
concentrations the crossing is instead near pH 6.328. The change is a boundary
convention effect, not a change in the thermodynamic constants.

Sophie can also enable \texttt{show removed domains=true} to display the band
between incompatible bounds; it is labeled as an incompatible domain, not as a
stable iodine region. The three-species model deliberately omits triiodide,
hypoiodous acid and solid iodine. It illustrates limiting-domain construction,
not complete iodine speciation or kinetic persistence.

\subsection{Reading the disappearance point and the replacement line}
Sophie can query both provisional and final data from the same completed
snapshot. Here she chooses equal element concentrations. The provisional limits
retain their \texttt{-temp} names even after iodine loses its domain.

\begin{codeexample}[width=0pt]
\pgfpourbaixset{pH min=0,pH max=10,
 build={name=SophieReadout,element=I,
  species={IO3^-,I^-,I2(aq)},temperature=298.15,concentration=.001,
  boundary convention=equal element concentrations,
  thermodynamic database=none,constants={
   standard potential/I2(aq)|I^-={value=.620},
   standard potential/IO3^-|I^-={value=1.0972309}}}}
\pgfplotsgetpourbaix{SophieReadout}\IHandle
\pgfpourbaixset{data/\IHandle/.cd,
  1-temp/b/.get=\lowerIntercept,
  3-temp/a/.get=\upperSlope,
  3-temp/b/.get=\upperIntercept,
  final/2/a/.get=\newSlope,
  final/2/b/.get=\newIntercept,
  final/2/expression/.get=\replacement,
  selection/event 1/pH min/.get=\disappearance}
\pgfkeys{/pgf/fpu=true,/pgf/fpu/output format=fixed}
\pgfmathsetmacro\intersection
  {(\lowerIntercept-\upperIntercept)/\upperSlope}
\pgfkeys{/pgf/fpu=false}
\begin{tikzpicture}
\node[align=left] {
  Intersection: $\mathrm{pH}=\pgfmathprintnumber[precision=3]{\intersection}$\\
  Event starts: $\mathrm{pH}=\pgfmathprintnumber[precision=3]{\disappearance}$\\
  Replacement: $E=(\pgfmathprintnumber[precision=5]{\newSlope})\,
  \mathrm{pH}+(\pgfmathprintnumber[precision=5]{\newIntercept})$};
\end{tikzpicture}
\end{codeexample}

The intersection computed from the queried provisional coefficients agrees with
the start of the recorded elimination event to the three displayed decimals.
PGF's floating-point unit avoids
the limited precision of its dimension-based arithmetic in this calculation.
Changing the boundary convention
in this declaration reproduces the other threshold, without changing any
numerical coefficient in the extraction code. The identifiers used here apply
to this species inventory and pH window; they are not universal iodine labels.
In the side-by-side example, put analogous queries inside the \verb|\foreach|
body, where its named diagram and convention are still in scope.

\subsection{Three ways to organize the species}
Before discussing the sloping lines, Sophie asks the class to organize the
species. The oxidation-number picture places I$^-$, I$_2$(aq) and IO$_3^-$
in the order $-1,0,+5$. The acidity picture compares species \emph{within each
oxidation group}. Here each row contains just one species: there is no acid--base
pair to order. In particular, IO$_3^-$ is not being classified as more or less
acidic than I$^-$ or I$_2$. A hydroxide-containing system such as Maxime's iron
example has nontrivial acidity rows.

Finally, the combined map places oxidation number upwards and acidity
horizontally. Its links carry the same final boundary identifiers as the
potential--pH plot. They indicate retained pairs somewhere in the selected pH
window, not simultaneous coexistence of every linked species. I$_2$ remains
visible because it retains a domain over part of this window.

All three are ordinary TikZ pics, used outside an \texttt{axis}. Sophie chooses
\texttt{acidity direction=increasing} for left-to-right reading; the default is
\texttt{decreasing}. This only changes presentation, not the chemical data,
provisional inventory or final boundary identifiers. The declaration is repeated
here to make the example self-contained.

\begin{codeexample}[width=0pt]
\pgfpourbaixset{pH min=0,pH max=10,
 build={name=SophieMap,element=I,
  species={IO3^-,I^-,I2(aq)},temperature=298.15,concentration=.001,
  boundary convention=equal element concentrations,
  thermodynamic database=none,constants={
   standard potential/I2(aq)|I^-={value=.620},
   standard potential/IO3^-|I^-={value=1.0972309}}}}
\begin{tikzpicture}[
 pourbaix/species/.append style={fill=blue!8},
 pourbaix/acidity direction=increasing,
 pourbaix/caption/.append style={font=\small}]
 \pic (states) {pourbaix oxidation=SophieMap};
 \pic (acids) at (0,-2.5) {pourbaix acidity=SophieMap};
 \pic[pourbaix/map row sep=1.3cm] (combined)
   at (6,-5.4) {pourbaix map=SophieMap};
\end{tikzpicture}
\end{codeexample}

The oxidation and acidity pics need only a declaration: no thermodynamic
constants or boundary calculations are needed for these classifications.
The combined \texttt{pourbaix map} needs finalization to obtain its labeled
links; \texttt{build} has already performed it here. The three pic names keep
their node names distinct and allow additional TikZ annotations.

\section{Tutorial: Aya Brings Her Own Uranium Data}
\label{sec:tutorial-uranium}
Aya is preparing a presentation on uranium thermochemistry. A ready-made figure
will not do: she wants the diagram to use the particular numerical dataset in
her notes. Her chosen species are U(s), U$^{4+}$, U$^{3+}$, UO$_2^{2+}$,
U(OH)$_3$(s), U(OH)$_4$(s), and UO$_2$(OH)$_2$(s), and her total uranium
concentration is $1\,\mathrm{mol\,L^{-1}}$.

\subsection{Transcribing the data without changing them}
For this worked example Aya takes her constants to apply at 298.15 K and to the
SHE reference. The potentials and solubility products are:
\begin{center}
\begin{tabular}{lr}
Datum & Value\\\hline
$E^\circ(\mathrm{U^{4+}/U(s)})$ & $-1.51$ V\\
$E^\circ(\mathrm{U^{4+}/U^{3+}})$ & $-0.61$ V\\
$E^\circ(\mathrm{UO_2^{2+}/U^{4+}})$ & $0.33$ V\\
$pK_s(\mathrm{U(OH)_3(s)})$ & $18.5$\\
$pK_s(\mathrm{U(OH)_4(s)})$ & $49.0$\\
$pK_s(\mathrm{UO_2(OH)_2(s)})$ & $24.0$
\end{tabular}
\end{center}
These are Aya's supplied example data, not evaluated uranium entries from the
offline catalogue. No crystalline form is specified, so the example makes no
identification with a particular mineral phase. She disables numerical fallback
and supplies $\log_{10}K_w=-14$ explicitly. Every selected species contains one
uranium atom; equal species and equal element concentrations are equivalent here.
At this concentration the ideal activity model is an illustrative approximation,
not a quantitatively validated model of a concentrated uranium solution.

\subsection{Converting the uranyl solubility product by hand}
A \texttt{log value} is $\log_{10}K$, so the first two solubility products are
entered as $-18.5$ and $-49$, not as positive p$K_s$ values.
For the uranyl solid, the dissolution reaction is
\[
 \mathrm{UO_2(OH)_2(s)}=
 \mathrm{UO_2^{2+}}+2\mathrm{OH^-}.
\]
The \texttt{solubility product} shorthand infers only dissolution to a monatomic
cation. It must not interpret the two oxo oxygen atoms as hydroxide groups.
Aya first performs the conversion by hand, using the equal-oxidation-state \texttt{equilibrium constant}
record for the acid-balanced pair:
\[
 \mathrm{UO_2(OH)_2(s)}+2\mathrm{H^+}=
 \mathrm{UO_2^{2+}}+2\mathrm{H_2O},\qquad
 \log_{10}K=-24-2(-14)=4.
\]
This is precisely her third solubility product expressed in the normalization
of the pair interface. It is not an additional fitted datum.

\begin{codeexample}[width=0pt]
\pgfpourbaixset{pH min=0,pH max=14,
 build={name=AyaU,element=U,temperature=298.15,concentration=1,
   species={U(s),U^4+,U^3+,UO2^2+,U(OH)3(s),U(OH)4(s),UO2(OH)2(s)},
   boundary convention=equal element concentrations,
   thermodynamic database=none,
   constants={
     standard potential/U^4+|U(s)={value=-1.51},
     standard potential/U^4+|U^3+={value=-.61},
     standard potential/UO2^2+|U^4+={value=.33},
     solubility product/U(OH)3(s)={log value=-18.5},
     solubility product/U(OH)4(s)={log value=-49},
     equilibrium constant/UO2(OH)2(s)|UO2^2+={log value=4},
     water ion product={log value=-14}}}}
\foreach \palette/\caption in
 {tab10 light/Presentation,monochrome/Handout}{
 \begin{tikzpicture}
 \begin{axis}[pourbaix/axis,title={\caption},
   xmin=0,xmax=14,ymin=-2.5,ymax=1.2,width=7.4cm,height=7cm]
  \addpourbaix[domain palette=\palette,
    boundary/b={label style={right,xshift=2pt}},
    final domain label style={font=\scriptsize,fill=white,inner sep=1pt}]{AyaU};
 \end{axis}
 \end{tikzpicture}
}
\end{codeexample}

\subsection{One calculation, two audiences}
Aya uses \texttt{build} once, outside the pictures. The \verb|\foreach| loop
reuses the same immutable diagram: \texttt{tab10 light} gives a color version
for her presentation, and \texttt{monochrome} gives black TikZ patterns for her
handout. Adjacent retained domains receive different automatic fills in both
versions. The white backgrounds keep species labels legible over the patterns.
No water overlay is needed for this comparison. These are presentation choices;
they change neither the thermodynamic data nor the computed boundaries.

\subsection{Letting the library read the dissolution reaction}
Aya can also supply the dissolution equation directly with the optional
\texttt{reaction} key. She keeps the original $\log_{10}K_s=-24$; the library
checks atoms and charge, then combines this equation with the water ion product.
It obtains the same acid-dissolution constant as her hand calculation.
The following complete input gives the same diagram. Only the uranyl constant
record differs from the first version.

\begin{codeexample}[width=0pt]
\pgfpourbaixset{pH min=0,pH max=14,
 build={name=AyaReaction,element=U,
   temperature=298.15,concentration=1,
   species={U(s),U^4+,U^3+,UO2^2+,U(OH)3(s),U(OH)4(s),UO2(OH)2(s)},
   boundary convention=equal element concentrations,
   thermodynamic database=none,constants={
     standard potential/U^4+|U(s)={value=-1.51},
     standard potential/U^4+|U^3+={value=-.61},
     standard potential/UO2^2+|U^4+={value=.33},
     solubility product/U(OH)3(s)={log value=-18.5},
     solubility product/U(OH)4(s)={log value=-49},
     solubility product/UO2(OH)2(s)={
       reaction={UO2(OH)2(s) = UO2^2+ + 2 OH^-},log value=-24},
     water ion product={log value=-14}}}}
\pgfplotsgetpourbaix{AyaReaction}\AyaReactionHandle
\pgfpourbaixset{data/\AyaReactionHandle/.cd,
 input 6/log K/.get=\AyaSourceLogK,
 species 7/log solubility product/.get=\AyaLogKs}
\begin{minipage}{11cm}
Source $\log_{10}K_s=\pgfmathprintnumber{\AyaSourceLogK}$;
checked $\log_{10}K_s=\pgfmathprintnumber{\AyaLogKs}$.
\par
\begin{tikzpicture}
 \begin{axis}[pourbaix/axis,width=10cm,height=7cm,
   xmin=0,xmax=14,ymin=-2.5,ymax=1.2]
  \addpourbaix[domain palette=tab10 light,
    boundary/b={label style={right,xshift=2pt}},
    final domain label style={font=\scriptsize}]{AyaReaction};
 \end{axis}
\end{tikzpicture}
\end{minipage}
\end{codeexample}

Spaces around the separating \texttt{+} distinguish addition from ionic charges;
\verb|2 OH^-| has a coefficient of two, whereas \verb|UO2^2+| is one ion.
The explicit equation fixes both direction and normalization: doubling all
coefficients would require \texttt{log value=-48}, not $-24$.
The input reaction and its coefficients remain queryable, and report input cards
show the supplied reaction alongside its logarithmic constant. Neither hydroxide
ions nor water become additional diagram domains. Omitting \texttt{reaction}
keeps the documented automatic convention; it never guesses which convention
Aya's numerical source used.

\subsection{Organizing seven species before discussing the boundaries}
Aya also wants a slide that explains the structure of the diagram. Unlike
Sophie's iodine set, her uranium set has several species at the same oxidation
number. The oxidation-only picture separates four groups: U(s) at 0,
U$^{3+}$ and U(OH)$_3$(s) at $+3$, U$^{4+}$ and U(OH)$_4$(s) at $+4$, and
UO$_2^{2+}$ and UO$_2$(OH)$_2$(s) at $+6$.

The acidity picture then orders the species within each group. Aya chooses
increasing acidity, as in Sophie's example: each hydroxide solid is to the
left of its dissolved cation. This is the stoichiometric H$^+$/H$_2$O order, not an order
of numerical solubility products. For example, acid dissolves the uranyl solid
according to the reaction above. No acidity comparison between different
oxidation groups is implied, and positions are not numerical pH values.

Aya combines the two classifications in the third picture. Oxidation number
increases upwards and acidity increases rightwards. The letters and numbers
match the final uranium diagram. Here is a self-contained version using the same
input data; if \texttt{AyaU} is still in scope, she can instead reuse it in the
three pic commands without building another diagram. Aya sets
\texttt{pourbaix/oxidation number format=roman} on the picture to display
$0$, $+\mathrm{III}$, $+\mathrm{IV}$ and $+\mathrm{VI}$ in all three pics.
This changes only the labels; the default notation uses Arabic numerals.

\begin{codeexample}[width=0pt]
\pgfpourbaixset{pH min=0,pH max=14,
 build={name=AyaMap,element=U,temperature=298.15,concentration=1,
  species={U(s),U^4+,U^3+,UO2^2+,U(OH)3(s),U(OH)4(s),UO2(OH)2(s)},
  boundary convention=equal element concentrations,
  thermodynamic database=none,constants={
   standard potential/U^4+|U(s)={value=-1.51},
   standard potential/U^4+|U^3+={value=-.61},
   standard potential/UO2^2+|U^4+={value=.33},
   solubility product/U(OH)3(s)={log value=-18.5},
   solubility product/U(OH)4(s)={log value=-49},
   equilibrium constant/UO2(OH)2(s)|UO2^2+={log value=4},
   water ion product={log value=-14}}}}
\begin{tikzpicture}[
 pourbaix/species/.append style={fill=blue!8,font=\small},
 pourbaix/oxidation number format=roman,
 pourbaix/acidity direction=increasing,
 pourbaix/caption/.append style={font=\small}]
 \pic (states) {pourbaix oxidation=AyaMap};
 \pic (acids) at (0,-3) {pourbaix acidity=AyaMap};
 \pic[pourbaix/map row sep=1.5cm,
   pourbaix/map long edge/.style={bend right=75}] (combined)
   at (7.5,-7.5) {pourbaix map=AyaMap};
\end{tikzpicture}
\end{codeexample}

All seven species appear in the two classification pictures, which require
only a declaration. Classification alone does not decide whether a species has
a stable domain. In the combined map, U(OH)$_3$(s) is dimmed and dashed, with no
retained boundary: finalization eliminates it throughout the selected pH range.
U$^{3+}$ remains normally drawn because it has a domain over part of that range.
The map summarizes retained pairs over the entire window; its links do not claim
that all domains coexist at one pH. This distinction lets Aya introduce the
elimination reactions before returning to the quantitative E--pH plot.

\subsection{Explaining the result in the presentation}
The seven declared species give 21 provisional pairs. Aya did not need 21
independent constants: the three potentials, three solubility data and water
constant determine them by reaction algebra. For example,
\[
 E^\circ(\mathrm{U^{3+}/U})=
 \frac{4(-1.51)-(-0.61)}{3}=-1.81\ \mathrm V.
\]
At the existence limits of U(OH)$_4$(s) and UO$_2$(OH)$_2$(s), the dissolved
species concentration is $c_t$. Their vertical boundaries are respectively at
pH $14-49/4=1.75$ and $14-24/2=2.00$.
For U(OH)$_3$(s), the corresponding \emph{provisional} value is
$14-18.5/3=7.8333$; it is not a retained vertical boundary.

Aya notices that U(OH)$_3$(s) has no final domain. This is a result, not a missing
input. Its solid--solid upper potential limit lies below its metal--solid lower
limit for this dataset; disproportionation into U(s) and U(OH)$_4$(s) is favored:
\[
 4\mathrm{U(OH)_3(s)}=
 \mathrm{U(s)}+3\mathrm{U(OH)_4(s)}.
\]
The dissolved U$^{3+}$ domain also disappears above about pH 6.821, where its
limits against U(s) and U(OH)$_4$(s) meet. As in Sophie's example, provisional
line overlays can show why the final inventory is smaller than the declaration.

Aya keeps one qualification on her presentation slide: $c_t=1$ mol/L is used here
with the library's ideal-activity model. Activity corrections can be substantial
at that concentration. The chosen seven-species diagram and its unspecified
solid phases are a calculation with her supplied data, not a complete uranium
speciation or mineral-stability model. Standard potentials determine reaction
Gibbs energies at the stated temperature; they do not by themselves determine
enthalpies or temperature dependence.

\subsection{Preparing the speaker's notes}
Aya expects questions about the missing U(OH)$_3$ domain and the constants she
did not enter herself. Instead of copying numbers into a separate document,
she asks \cs{pgfplotspourbaixreport} to assemble her notes from the completed
diagram. The report below contains the final plot, the elimination events and
their reactions, the thermodynamic inputs, recovered constants, and the equations
of all 21 provisional boundaries with their retained or inactive status.

The classifications have already appeared on her introductory slide, so she
omits their report sections. She keeps the plot free of diagnostic overlays
with \texttt{show removed=false}; the event cards still explain the eliminated
domains. To include the three schematics as well, add \texttt{oxidation},
\texttt{acidity} and \texttt{map} to the section list. There is no water overlay.

This self-contained code repeats her input data. If \texttt{AyaU} is still in
scope, only the report command is needed, with \texttt{AyaU} as its argument.
A report reuses the selected snapshot; it does not change the diagram or run a
new chemical calculation. Use the command in ordinary document flow, outside
\texttt{tikzpicture} and \texttt{axis}, so pages can break between its cards.

\begin{verbatim}
\begingroup
\pgfpourbaixset{pH min=0,pH max=14,
 build={name=AyaNotes,element=U,temperature=298.15,concentration=1,
  species={U(s),U^4+,U^3+,UO2^2+,U(OH)3(s),U(OH)4(s),UO2(OH)2(s)},
  boundary convention=equal element concentrations,
  thermodynamic database=none,constants={
   standard potential/U^4+|U(s)={value=-1.51},
   standard potential/U^4+|U^3+={value=-.61},
   standard potential/UO2^2+|U^4+={value=.33},
   solubility product/U(OH)3(s)={log value=-18.5},
   solubility product/U(OH)4(s)={log value=-49},
   equilibrium constant/UO2(OH)2(s)|UO2^2+={log value=4},
   water ion product={log value=-14}}}}
\pgfplotspourbaixreport[
 sections={header,plot,events,inputs,constants,boundaries},
 text width=12cm,
 axis options={xmin=0,xmax=14,ymin=-2.5,ymax=1.2},
 plot options={show removed=false,domain palette=tab10 light,
   boundary/b={label style={right,xshift=2pt}},
   final domain label style={fill=white,inner sep=1pt}}
]{AyaNotes}
\endgroup
\end{verbatim}

Here are Aya's generated notes. Input log constants are normalized reaction
constants, not all solubility products. Each boundary card links its equation
to the inputs used to derive it, so she can explain, for example, the recovered
$E^\circ(\mathrm{U^{3+}/U})=-1.81$ V without treating potentials as unweighted
additive quantities. The report includes inactive candidates for reference;
only retained boundaries appear as solid lines in the plot.

\begingroup
\pgfpourbaixset{pH min=0,pH max=14,
 build={name=AyaNotes,element=U,temperature=298.15,concentration=1,
  species={U(s),U^4+,U^3+,UO2^2+,U(OH)3(s),U(OH)4(s),UO2(OH)2(s)},
  boundary convention=equal element concentrations,
  thermodynamic database=none,constants={
   standard potential/U^4+|U(s)={value=-1.51},
   standard potential/U^4+|U^3+={value=-.61},
   standard potential/UO2^2+|U^4+={value=.33},
   solubility product/U(OH)3(s)={log value=-18.5},
   solubility product/U(OH)4(s)={log value=-49},
   equilibrium constant/UO2(OH)2(s)|UO2^2+={log value=4},
   water ion product={log value=-14}}}}
\pgfplotspourbaixreport[
 sections={header,plot,events,inputs,constants,boundaries},
 text width=12cm,
 axis options={xmin=0,xmax=14,ymin=-2.5,ymax=1.2},
 plot options={show removed=false,domain palette=tab10 light,
   boundary/b={label style={right,xshift=2pt}},
   final domain label style={fill=white,inner sep=1pt}}
]{AyaNotes}
\endgroup

\section{Reading the Reference Parts}

\subsection{How to use this manual}
Part~\ref{part:reference} documents the input keys, named snapshots and plotting
commands by task. Part~\ref{part:science} explains the limiting chemical model,
thermodynamic checking and selection algorithms. Part~\ref{part:examples}
contains complete diagrams and composed reports. Part~\ref{part:data} describes
the offline catalogue, standard states, solid phases and provenance.
Part~\ref{part:technical} covers the TeX/Lua boundary, diagnostics and local builds.

\subsection{Scope and limitations}
A diagram describes the selected species only. It does not silently add missing
domains and does not solve full aqueous speciation. Source data can mix molality
and molarity conventions under the stated teaching approximation. No temperature
extrapolation or general reference-electrode conversion is performed. Nonideal
activities, arbitrary ligand equilibria and kinetic overpotentials are outside
the current model. A diagram outside water's stability window can still be drawn,
but this does not predict experimentally persistent aqueous phases.

\clearpage\part{User reference}\label{part:reference}
\section{Commands, key families and workflow}
The usual workflow is to build a named diagram, then draw it inside an axis.
For intermediate inspection, use the separate stages in this order:
\texttt{declare}, \texttt{check}, \texttt{calculate}, \texttt{finalize}.
A declaration already supplies classification pics; a check supplies completed
thermodynamic data; calculation supplies provisional equations; finalization
supplies the retained domains, final boundaries and report data.
The \texttt{build} action performs all four stages atomically.

Chemical declarations and plot appearance have distinct key families:
\begin{center}
\begin{tabular}{ll}
Configuration & Relative key family\\\hline
\cs{pgfpourbaixset} & \texttt{/pgfplots/pourbaix}\\
\cs{addpourbaix} options & \texttt{/pgfplots/pourbaix/plot}\\
\cs{addpourbaixwater} options & \texttt{/pgfplots/pourbaix/water}\\
Classification pic styles & \texttt{/tikz/pourbaix}\\
\end{tabular}
\end{center}
The \texttt{pourbaix} environment combines chemical keys with an axis: its
unknown top-level keys are searched in the PGFPlots family. Explicit absolute
paths disambiguate collisions. Neither the setters nor the environment draw a
diagram implicitly. Normal PGF styles, handlers and TeX grouping apply.

\begin{command}{\pgfpourbaixset\marg{keys}}
Processes a key list in the \texttt{/pgfplots/pourbaix} family.
The definition is exactly:
\begin{verbatim}
\def\pgfpourbaixset#1{\pgfkeys{/pgfplots/pourbaix/.cd,#1}}
\end{verbatim}
There is no extra group or automatic calculation. Relative paths start in this
family; absolute paths, styles, \texttt{.cd}, \texttt{.get} and other PGF handlers
keep their usual meaning. Unknown keys use normal PGF error handling.
Assignments and retrieved macros follow the caller's TeX scope. The command
is available in plain LuaTeX, LuaLaTeX and ConTeXt MkIV after loading the library.
For example, \verb|\pgfpourbaixset{pH min=0,pH max=14}| sets the sweep bounds.
Unlike the \texttt{pourbaix} environment, this shortcut does not forward unknown
keys to the axis family. Use \cs{pgfplotsset} for ordinary PGFPlots settings and
\cs{tikzset} for pic styles. The tutorials demonstrate declarations, actions
and data queries with this command.
\end{command}

\section{Chemical notation and numerical input}
\subsection{Species grammar}
Species are chemical data strings, not typesetting code. Supported forms are:
\begin{center}
\begin{tabular}{ll}
\textbf{Input} & \textbf{Meaning}\\
\texttt{Fe(s)} & Neutral solid\\
\texttt{Fe\string^2+} & Aqueous ion of charge $+2$\\
\texttt{Fe\string^\{2+\}(aq)} & The same ion, with an explicit aqueous suffix\\
\texttt{Fe(OH)2(s)} & Solid with one Fe, two O and two H atoms\\
\texttt{Fe2(SO4)3(s)} & Solid with two Fe, three S and twelve O atoms\\
\texttt{H2CO3(aq)} & Neutral aqueous species, usable for element C\\
\texttt{Cr2O7\string^2-} & Aqueous ion with two chromium atoms
\end{tabular}
\end{center}
Element symbols are case sensitive. Subscripts are positive integers after an
element symbol or a parenthesized group; groups may be nested. Ionic charges must
follow a caret, optionally inside braces. A unit charge can be written
\texttt{\string^+} or \texttt{\string^1+}. A charged species without a phase suffix
is aqueous. A neutral species requires either \texttt{(s)} or \texttt{(aq)}.
Solids must be electrically neutral. The selected element must occur in every
listed species.

Outer whitespace is ignored. Internal whitespace, TeX macros, math delimiters,
underscores, Unicode charge glyphs, gas/liquid suffixes, hydration dots and square
brackets are not supported by this grammar. Use plain formula data such
as \texttt{Fe\string^2+}, not \verb|\mathrm{Fe}^{2+}| or \verb|\ce{Fe^2+}|.
Unsupported notation produces an error rather than a guessed composition.

The parser normalizes optional aqueous suffixes on ions, braced charges and unit
charge magnitudes for identity checks. It preserves formula group structure;
it does not infer isomers, crystal phases or chemical equivalence between
alternative structural formulas. Distinct polymorphs with the same formula are
not supported by the chemical formula parser.

\subsection{Numbers and units}
Numbers are signed decimal literals, optionally with a decimal point and an
\texttt{e} or \texttt{E} exponent. Examples include \texttt{1E-2}, \texttt{.77} and
\texttt{-0.44}. Expressions, hexadecimal notation and attached unit text are
rejected. The concentration must be finite and strictly positive, in mol/L.
Potentials must be finite and are in volts; negative potentials are allowed.
Solubility products must be finite and strictly positive and are dimensionless
thermodynamic constants. No logarithm is implied: \texttt{value=15} means $K_s=15$,
not $pK_s=15$. Underflow to zero is rejected for positive quantities.

The working temperature is required as a finite positive number in kelvins.
No default temperature, reference electrode, standard concentration or activity
model is inferred. Constants must be appropriate to the declared temperature;
the library does not apply temperature corrections. Thermodynamic closure
requires compatible input standard states and a common potential reference.
The water default during checking is the exception described below.

\section{Declaration key reference}


\begin{key}{/pgfplots/pourbaix/build=\marg{options}}
Required option-list value, with no default. Accepts exactly the declaration
keys below, then executes declaration, checking, calculation and finalization
in that order. Set \texttt{pourbaix/pH min} and \texttt{pH max} before this key;
\texttt{pH step} defaults to 0.1. Bounds are finite pH values and are independent
of axis limits. No picture or plot is created.

The action commits one completed snapshot and exports all checked, provisional
and final data for \texttt{.get}. It updates \cs{pgfplotslastpourbaix} and the
optional name in the caller's TeX scope. Any failed stage reports its normal
error and commits nothing: existing aliases, handles and data remain intact.
Missing, contradictory or invalid thermodynamic inputs are errors, not warnings
that permit an unchecked plot. Available catalogue data and derived constants
are used under the same rules as \texttt{check}. Use the separate stage keys
when inspecting an incomplete declaration or retaining intermediate snapshots.

Both the \texttt{pourbaix} environment and \cs{pgfpourbaixset} accept
unprefixed keys such as \texttt{pH min} and \texttt{build}.
With \cs{pgfplotsset}, add the \texttt{pourbaix/} prefix.
\end{key}

\begin{key}{/pgfplots/pourbaix/declare=\marg{options}}
Action key, with a required value and no default. Parses \meta{options} relative
to \path{/pgfplots/pourbaix/diagram}. Use through \cs{pgfpourbaixset} before
plotting, outside or inside a picture. It validates all data before committing
a snapshot, then sets \verb|\pgfplotslastpourbaix| and the optional name alias.
It has no drawing or axis side effects. A failed action changes neither alias
nor last handle.
\end{key}

The following keys are relative to \texttt{/pgfplots/pourbaix/diagram}.
\begin{key}{/pgfplots/pourbaix/diagram/name=\meta{name}}
Optional. The default is an anonymous diagram. A nonempty name starts with an
ASCII letter and then contains letters, digits, underscores, dots, colons or
hyphens. Names are case sensitive. Redefinition replaces the alias in the current
TeX group; an outer definition reappears when that group ends. No uniqueness
constraint is imposed across scopes, and no old snapshot is modified.
\end{key}
\begin{key}{/pgfplots/pourbaix/diagram/element=\meta{symbol}}
Required, with no default. One chemical element symbol, for example \texttt{Fe}
or \texttt{Cr}. Element names such as \texttt{iron} are not accepted. Every species
must contain the selected element. Other elements may occur in its formula.
\end{key}
\begin{key}{/pgfplots/pourbaix/diagram/species=\marg{list}}
Required, nonempty comma-separated list of formulas in the grammar above.
Order is preserved. Duplicate normalized identities, empty entries and formulas
without the selected element are rejected. This is the entire selected species
set: the library does not add species while validating constants.
\end{key}
\begin{key}{/pgfplots/pourbaix/diagram/concentration=\meta{number}}
Required, finite and strictly positive, in mol/L of the selected element.
This is the total diagram concentration $c_{\mathrm t}$, including generalized
solid contributions in the balance. It is not an independently assigned
concentration for every species.
\end{key}
\begin{key}{/pgfplots/pourbaix/diagram/nernst factor=\meta{auto or volts}}
Default \texttt{auto}: compute $s=RT\ln(10)/F$ from the required temperature.
Alternatively give a strictly positive finite decimal value in volts, without
a unit suffix, for example \texttt{0.06}. A missing value, zero, a negative
value or a nonfinite value is an error. The value belongs to the declared
snapshot and follows the declaration's TeX scope; it is not a plotting style.

Checking uses the effective $s$ for all conversions between standard potentials
and log equilibrium constants, including fallback records. Equation calculation,
finalization and reports use the same value. Temperature, catalogue matching,
water-product defaults and the standard states remain unchanged. This is a
numerical teaching approximation when it differs from $RT\ln(10)/F$, not a
new physical temperature or a temperature extrapolation.

After checking or building, the data root exports \texttt{nernst factor} in volts
and \texttt{nernst factor source}, either \texttt{temperature} or
\texttt{specified}. Both support native \texttt{.get}. The report summary also
prints $s$. Invalid input preserves existing diagrams and aliases.
\end{key}
\begin{key}{/pgfplots/pourbaix/diagram/temperature=\meta{number}}
Required positive finite absolute temperature in K, without a default. It belongs
to the diagram snapshot. Celsius suffixes and unit expressions are not accepted.
Input constants are taken to apply at this temperature; no extrapolation is
performed. Declarations written for version 0.1 must add this key.
\end{key}
\begin{key}{/pgfplots/pourbaix/diagram/boundary convention=\meta{choice}}
Required; there is deliberately no default. The only values are
\texttt{equal concentrations} and \texttt{equal element concentrations}.
Their definitions are given above. The setting belongs to this diagram and does
not change another diagram's convention.
\end{key}
\begin{key}{/pgfplots/pourbaix/diagram/constants=\marg{records}}
Required; \texttt{constants=\{\}} explicitly supplies an empty list. Parses records
relative to \texttt{/pgfplots/pourbaix/constants}. The currently supported records
are \texttt{standard potential}, \texttt{solubility product},
\texttt{equilibrium constant} and \texttt{water ion product}. Multiple records
are allowed for distinct couples or solids; water has at most one input record. Record order is preserved.
Repeated constants for the same unordered couple or the same solid are errors,
even if their values agree. No sufficiency, thermodynamic cycle consistency or
reaction balancing check is claimed at declaration time.
\end{key}

\subsection{Thermodynamic constant records}
\begin{key}{/pgfplots/pourbaix/constants/standard potential=\marg{record}}
Adds one standard \emph{reduction} potential. The record has exactly three
supported fields, all required, relative to
\texttt{/pgfplots/pourbaix/standard potential}. Each invocation starts with empty
fields; values never leak from a previous record.
\end{key}
\begin{key}{/pgfplots/pourbaix/standard potential/oxidized=\meta{species}}
The oxidized member, using the species grammar. It must be among the declared
species. This explicit orientation avoids interpreting a slash-separated label.
\end{key}
\begin{key}{/pgfplots/pourbaix/standard potential/reduced=\meta{species}}
The reduced member. It must be declared and distinct from the oxidized member.
The assigned oxidation numbers are checked for the declared direction, but the
declaration stage does not balance half-reactions. The explicit check
subsequently verifies balance and determines electron counts.
\end{key}
\begin{key}{/pgfplots/pourbaix/standard potential/value=\meta{number}}
Finite signed reduction potential in V at the diagram's explicit temperature.
There is no implicit potential reference. All constants intended for one
calculation must use consistent physical conventions.
\end{key}

\begin{key}{/pgfplots/pourbaix/constants/solubility product=\marg{record}}
Adds one thermodynamic solubility product associated with a declared solid.
Fields are relative to \texttt{/pgfplots/pourbaix/solubility product}. This stores
an input datum. The check infers dissolution into a free monatomic ion and
hydroxide for hydroxide solids only. Use the optional \texttt{reaction} field
for an explicit dissolution equation, including polyatomic ions.
\end{key}
\begin{key}{/pgfplots/pourbaix/solubility product/solid=\meta{species}}
Required reference to one of the declared solid species. An aqueous reference
is rejected. No solid is added automatically to the diagram.
\end{key}
\begin{key}{/pgfplots/pourbaix/solubility product/value=\meta{number}}
Dimensionless, finite, positive $K_s^\circ$. Supply exactly one of this field
and \texttt{log value}. It is a direct constant, not its negative logarithm.
\end{key}

\begin{key}{/pgfplots/pourbaix/solubility product/log value=\meta{number}}
Finite $\log_{10}K_s^\circ$, as an alternative to \texttt{value}. For example,
\texttt{log value=-15} means $K_s^\circ=10^{-15}$. There is no default.
\end{key}
\begin{key}{/pgfplots/pourbaix/constants/equilibrium constant=\marg{record}}
Adds a non-redox pair constant. Fields are relative to
\texttt{/pgfplots/pourbaix/equilibrium constant}; every record starts empty.
\end{key}
\begin{key}{/pgfplots/pourbaix/equilibrium constant/first=\meta{species}}
Required declared species. Without \texttt{reaction}, it is the reactant:
the check balances the pair with H$^+$ and H$_2$O
using primitive integer coefficients of the two selected species.
\end{key}
\begin{key}{/pgfplots/pourbaix/equilibrium constant/second=\meta{species}}
Required distinct declared species of the same oxidation number. Without
\texttt{reaction}, it is the product. With an explicit equation, the pair must
appear on opposite sides, in either orientation. The full
reaction direction and normalization are essential to interpreting its constant.
\end{key}
\begin{key}{/pgfplots/pourbaix/equilibrium constant/value=\meta{number}}
Finite positive dimensionless $K^\circ$ for that reaction. Supply exactly one
of this key and \texttt{log value}; no default is provided.
\end{key}
\begin{key}{/pgfplots/pourbaix/equilibrium constant/log value=\meta{number}}
Finite signed $\log_{10}K^\circ$, alternatively. This supports constants outside
the range representable as positive finite floating-point numbers.
\end{key}
\begin{key}{/pgfplots/pourbaix/solubility product/reaction=\marg{equation}}
Optional; omitted by default. An explicitly empty equation is an error.
The dimensionless \texttt{value} or
\texttt{log value} applies to exactly this equation, with its written direction
and coefficients. It overrides automatic reaction inference for this record.
The addressed solid must be consumed; all other net participants must be aqueous
products, except liquid water, which may occur on either side. Acid dissolution
with proton reactants belongs in an \texttt{equilibrium constant} record.
See Aya's tutorial for an executed uranyl example.
\end{key}
\begin{key}{/pgfplots/pourbaix/equilibrium constant/reaction=\marg{equation}}
Optional; omitted by default. An explicitly empty equation is an error.
The same syntax specifies the reaction for a pair
constant, overriding the automatic H$^+$/H$_2$O balance. Both addressed species
must occur on opposite sides and have the same oxidation number. Additional
carriers of the selected element are rejected for pair constants.
These fields are local to their constant record and follow declaration scope;
neither changes a previously stored snapshot.

Write exactly one \texttt{=}, separate participants with a plus sign surrounded
by whitespace, and separate a coefficient from its species with whitespace.
Coefficients are positive integers or positive rational numbers, such as
\texttt{1/2}; omitted coefficients equal one. Decimal coefficients and TeX
commands are not accepted. Formulae and charges use the species grammar;
neutral aqueous species require \texttt{(aq)}, solids require \texttt{(s)}.
Liquid water is \texttt{H2O(l)} (alias \texttt{H2O}); \verb|HO^-| is a
specific alias of \verb|OH^-|. This does not identify arbitrary isomers.
Gas participants and electrons are outside this non-redox interface.

The parser checks atom and charge conservation on a common integer scale;
coefficients and intermediate integer sums/products must not exceed $2^{53}-1$
in absolute magnitude. Invalid syntax, unbalanced equations, empty net reactions
and incompatible references are errors. No missing coefficient is repaired.
An inverse equation needs $-\log_{10}K$, and an equation multiplied by $m$ needs
$m\log_{10}K$. Rational coefficients retain their scale in the stored reaction;
the thermodynamic solver subsequently uses floating-point coefficients.

Auxiliary participants are algebraic species, not new diagram domains. Supplying
an equation does not provide an unknown ligand activity: the check must still
resolve every canonical boundary reaction by eliminating auxiliary participants.
Unsupported boundary models or insufficient independent data remain errors.
Reaction text is read as unexpanded macro content through the bridge, never
executed as Lua or TeX source; the normal explicit \texttt{.expanded} handler
remains the caller's opt-in. The validated source is exported as detokenized text.
\end{key}
\begin{key}{/pgfplots/pourbaix/constants/water ion product=\marg{record}}
Optional single water constant at the diagram temperature, with fields relative
to \texttt{/pgfplots/pourbaix/water ion product}. No selected-domain species are
added for H$^+$, OH$^-$ or H$_2$O. See Section~\ref{sec:thermodynamic-check}
for inference and the fallback at 298.15 K.
\end{key}
\begin{key}{/pgfplots/pourbaix/water ion product/value=\meta{number}}
Finite positive dimensionless $K_w^\circ$ for water dissociation. Supply exactly
one of this field and \texttt{log value} in an explicit record.
\end{key}
\begin{key}{/pgfplots/pourbaix/water ion product/log value=\meta{number}}
Finite signed $\log_{10}K_w^\circ$. A value of $-14$ represents $K_w^\circ=10^{-14}$,
not $pK_w=-14$.
\end{key}

\subsection{Classification overrides and boundary names}
\begin{key}{/pgfplots/pourbaix/diagram/species data=\marg{record}}
Optional repeatable record, with fields relative to
\path{/pgfplots/pourbaix/species data}. Each species may have at most one record.
References to undeclared species are errors. Fields start empty on each use.
\end{key}
\begin{key}{/pgfplots/pourbaix/species data/species=\meta{species}}
Required species reference using the usual formula grammar and normalization.
\end{key}
\begin{key}{/pgfplots/pourbaix/species data/oxidation number=\meta{value}}
Optional signed finite decimal or integer fraction with positive denominator.
Omission requests automatic assignment; an explicit value overrides it. No
restriction to integral values is imposed on an explicit effective average.
\end{key}
\begin{key}{/pgfplots/pourbaix/species data/acidity rank=\meta{number}}
Optional positive finite ordering number; smaller means earlier, more acidic.
If supplied for one species, it must be supplied for all members of that oxidation
group, with no ties. Required when automatic H$^+$/H$_2$O comparison is inapplicable.
Ranks belong to a single group and may be reused in another group.
\end{key}
\begin{key}{/pgfplots/pourbaix/diagram/boundary labels=\marg{options}}
Optional naming configuration relative to \path{/pgfplots/pourbaix/boundary labels}.
Repeated uses update the specified fields. Defaults are reset for each declaration.
\end{key}
\begin{key}{/pgfplots/pourbaix/boundary labels/vertical prefix=\meta{text}}
Default empty. Prepended to the alphabetic counter for vertical candidates.
\end{key}
\begin{key}{/pgfplots/pourbaix/boundary labels/redox prefix=\meta{text}}
Default empty. Prepended to the decimal counter for redox candidates.
\end{key}
\begin{key}{/pgfplots/pourbaix/boundary labels/suffix=\meta{text}}
Default \texttt{-temp}, appended to both counters. All naming fragments may
contain ASCII letters, digits, underscores, dots, colons and hyphens, or be empty.
They are identifiers, not TeX code.
\end{key}
\begin{key}{/pgfplots/pourbaix/diagram/boundary name=\marg{record}}
Optional repeatable override for one unordered pair. Its fields are relative to
\path{/pgfplots/pourbaix/boundary name}; an explicit name replaces the generated
name without changing either automatic counter.
\end{key}
\begin{key}{/pgfplots/pourbaix/boundary name/first=\meta{species}}
Required declared species reference, distinct from the second member.
\end{key}
\begin{key}{/pgfplots/pourbaix/boundary name/second=\meta{species}}
Required declared species reference. Endpoint order does not affect matching.
\end{key}
\begin{key}{/pgfplots/pourbaix/boundary name/name=\meta{name}}
Required nonempty identifier. Starts with an ASCII letter or digit, followed by
letters, digits, underscores, dots, colons or hyphens. It must be unique among
the final candidate names within this diagram.
\end{key}

\subsection{Reuse through PGF styles}
Declare reusable input under the diagram key family. Each call starts with empty
fields, applies the given keys in order and validates the completed object.
The following example uses only three species and two potentials.
\begin{codeexample}[width=3.2cm]
\pgfplotsset{
  pourbaix/diagram/simple iron/.style={
    element=Fe,temperature=298.15,
    species={Fe(s),Fe^2+,Fe^3+},
    constants={
      standard potential={
        oxidized=Fe^3+,reduced=Fe^2+,value=0.77},
      standard potential={
        oxidized=Fe^2+,reduced=Fe(s),value=-0.44}
    },
    boundary convention=equal element concentrations
  },
  pourbaix/declare={
    simple iron,name=FeDilute,concentration=1e-3}
}
\texttt{FeDilute declared.}
\end{codeexample}
A style stores a recipe; a declaration stores its validated result. Changing the
style later does not change a diagram already declared from it. A later scalar
key overrides an earlier one, as usual in PGF. A later \texttt{constants} or
\texttt{species} key replaces the entire earlier list.

\section{Addressed species, constants and boundaries}
\label{sec:addressed}
An addressed key puts an object's identity in the path and its properties in the
value. These forms complement the original record syntax and share its validation,
units, scoping and failure behavior. They do not add species or change the reaction
model. The preferred forms are summarized below.

\begin{center}
\begin{tabular}{ll}
Context & Relative key path\\\hline
Declaration & \texttt{species/}\meta{species}\\
Constants & \texttt{solubility product/}\meta{solid}\\
Constants & \texttt{standard potential/}\meta{oxidized}\texttt{\char124}\meta{reduced}\\
Constants & \texttt{equilibrium constant/}\meta{first}\texttt{\char124}\meta{second}\\
Plot options & \texttt{boundary/}\meta{final identifier}
\end{tabular}
\end{center}
Every addressed key requires an option-list value; empty lists are syntactically
allowed but cannot satisfy a required numerical constant. A bare key is an error.
A path identity cannot be replaced by a nested identity field. Unknown properties,
unknown species and malformed numerical data retain their usual diagnostics.
Existing valid snapshots and aliases survive failed declarations.

\subsection{Species properties}
\begin{key}{/pgfplots/pourbaix/diagram/species/\meta{species}=\meta{options}}
Apply properties to a species already present in the declaration's \texttt{species}
list. The value accepts \texttt{oxidation number}, \texttt{acidity rank} and
\texttt{thermodynamic phase}, with exactly the defaults and restrictions of the
older \texttt{species data} record. Unspecified fields are unchanged; repeated
property records retain the existing update behavior. No species is added
implicitly. Do not supply a nested \texttt{species} field in the addressed form.
\end{key}

\begin{codeexample}[width=6cm]
\pgfplotsset{pourbaix/declare={
 name=Magnetite,element=Fe,
 species={Fe3O4(s)},
 species/Fe3O4(s)={oxidation number=8/3},
 temperature=298.15,concentration=.01,
 boundary convention=equal concentrations,
 constants={}}}
\begin{tikzpicture}
 \pic {pourbaix oxidation=Magnetite};
\end{tikzpicture}
\end{codeexample}
The fractional oxidation number is an average for the selected element; the
classification does not resolve distinct iron sites. This example requires only
preparation and therefore supplies no numerical equilibrium data.

The older form remains equivalent:
\begin{verbatim}
species data={species=Fe3O4(s),oxidation number=8/3}
\end{verbatim}
For a solid phase, write for example
\texttt{species/Al(OH)3(s)=\{thermodynamic phase=amorphous\}} inside a declaration
that includes Al(OH)3(s). The usual database phase availability rules still apply.

\subsection{Thermodynamic constants and the pipe separator}
\begin{key}{/pgfplots/pourbaix/constants/solubility product/\meta{solid}=\meta{options}}
Equivalent to \texttt{solubility product=\{solid=}\meta{solid}\texttt{,...\}}.
The options require exactly one of \texttt{value} (positive dimensionless $K_s$)
or \texttt{log value} ($\log_{10}K_s$). The solid must be selected and supported
by the existing hydroxide-dissolution model. A nested \texttt{solid} field is
forbidden because the path already identifies it.
\end{key}
\begin{key}{/pgfplots/pourbaix/constants/standard potential/\meta{pair}=\meta{options}}
The path contains exactly two species, in \emph{oxidized--reduced} order. Supply
\texttt{value=}\meta{volts}, the standard reduction potential. The library balances
the reduction with H$^+$, water and electrons using its existing reaction model.
No temperature or reference-electrode conversion is introduced. The nested
\texttt{oxidized} and \texttt{reduced} fields are forbidden. Incompatible species,
reversed oxidation order and inconsistent constants remain errors.
\end{key}
\begin{key}{/pgfplots/pourbaix/constants/equilibrium constant/\meta{pair}=\meta{options}}
The two species have equal oxidation number. The constant refers to the existing
balanced reaction from \meta{first} to \meta{second}, including H$^+$ and water
as required; it is not simply their concentration ratio. Reversing the species
reverses the reaction and therefore inverts $K$. Supply exactly one of
\texttt{value} or \texttt{log value}, both dimensionless. Nested \texttt{first} and
\texttt{second} fields are forbidden. Different spectator-ligand compositions
remain outside the supported balancing model.
\end{key}

Use a literal \texttt{\char124} between species, not a slash. The slash is reserved for
PGF key path components. Zero separators, repeated separators and empty endpoints
are rejected. Surrounding endpoint spaces are trimmed by the chemical parser.
Do not write a mathematical \cs{vert}, \cs{mid}, or TeX typeset chemical formula
in a key address.

In plain LuaTeX and LuaLaTeX the ordinary pipe has character category~12. In
ConTeXt MkIV it can be active. The declaration locally assigns that active token
an expandable ordinary-pipe meaning before PGF expands key names, then restores
its meaning on leaving the declaration group. No global catcode or active-character
change is made. This normalization concerns declarative chemical data, not node
text. Endpoint splitting occurs in the bridge as inert string processing; species
text never becomes Lua code. Custom host active-pipe meanings are restored too.

\begin{codeexample}[width=6cm]
\pgfplotsset{pourbaix/declare={
 name=Hydrolysis,element=Fe,
 species={Fe^2+,FeOH^+},
 temperature=298.15,concentration=.01,
 boundary convention=equal concentrations,
 thermodynamic database=none,
 constants={
  equilibrium constant/Fe^2+|FeOH^+={
    log value=-9}}},
 pourbaix/check=Hydrolysis,
 pourbaix/calculate=Hydrolysis}
\pgfplotsgetpourbaix{Hydrolysis}\acidHandle
\edef\acidRoot{/pgfplots/pourbaix/data/\acidHandle}
\pgfkeysgetvalue{\acidRoot/a-temp/pH}\acidPH
$pH=\pgfmathprintnumber{\acidPH}$
\end{codeexample}
The illustrative constant is for
$\mathrm{Fe^{2+}+H_2O= FeOH^++H^+}$.
With equal aqueous concentrations, this gives the independently checkable
boundary pH~9. It is an interface example, not a new catalogue reference value.

\subsection{A complete inline example}
This example combines species, solubility, potential, boundary and domain
addresses. Numerical values are illustrative, as in the introductory iron example;
assigning a solid phase does not certify these numbers for that phase.

\begin{codeexample}[width=8cm]
\pgfplotsset{pourbaix/declare={
 name=AddressedFe,element=Fe,
 species={Fe(s),Fe^2+,Fe^3+,
          Fe(OH)2(s),Fe(OH)3(s)},
 species/Fe(OH)3(s)={
   thermodynamic phase=amorphous},
 temperature=298.15,concentration=.01,
 boundary convention=equal concentrations,
 constants={
  standard potential/Fe^3+|Fe^2+={value=.77},
  standard potential/Fe^2+|Fe(s)={value=-.44},
  solubility product/Fe(OH)2(s)={log value=-15},
  solubility product/Fe(OH)3(s)={log value=-38}}},
 pourbaix/check=AddressedFe,
 pourbaix/calculate=AddressedFe,
 pourbaix/pH min=0,pourbaix/pH max=14,
 pourbaix/finalize=AddressedFe}
\begin{tikzpicture}
\begin{axis}[pourbaix/axis,
 width=8cm,height=6cm,
 xmin=0,xmax=14,ymin=-1.2,ymax=1.2]
 \addpourbaix[
  domain label mode=letters,
  boundary/1={style={red,thick},
    label=$r_1$,label style={above,pos=.65}},
  boundary/a={style={dashed},
    label style={right}},
  domain/Fe(s)={label=$M$}
 ]{AddressedFe};
\end{axis}
\end{tikzpicture}
\end{codeexample}
The result is still queried under \texttt{final/1/}, even though its drawn label
is $r_1$. The older constant forms remain accepted in the same declaration:
\begin{verbatim}
standard potential={oxidized=Fe^3+,reduced=Fe^2+,value=.77}
solubility product={solid=Fe(OH)2(s),log value=-15}
\end{verbatim}
Do not provide both versions of the same constant: normal duplicate-constant
validation still applies. The addressed keys introduce no new scientific defaults.

\section{Names, handles and TeX scope}
\begin{command}{\pgfplotsgetpourbaix\marg{name}\meta{macro}}
Looks up a literal declared name visible in the current group and defines
\meta{macro} to its opaque handle. The assignment is local to the caller's group.
An unknown or invalid name is an error; the target macro remains unchanged.
Copy the returned handle rather than deferring a lookup of the name.
\end{command}
\begin{command}{\pgfplotslastpourbaix}
Expands to the handle of the most recent successful declaration in the current
scope, including anonymous declarations. Initially empty. Use
\verb|\let\saved\pgfplotslastpourbaix| to capture a handle. Using
\verb|\def\saved{\pgfplotslastpourbaix}| would instead follow future changes.
\end{command}

A name is an alias; a handle identifies an independent snapshot. An inner group
may redeclare a name with a different concentration. Within that group the alias
points to the new snapshot; afterwards the outer alias is restored. A previously
captured handle remains valid and retains its original data. Snapshots are kept
for the engine session, including those whose local aliases have expired. This
prevents deferred plots from losing their input before visualization.
Applications declaring very large numbers of diagrams should account for that
session-lifetime storage; garbage collection of diagram records is not exposed.

\section{The pourbaix axis environment}
\label{sec:pourbaix-environment}
The \texttt{pourbaix} environment is an ordinary PGFPlots axis inside an existing
\texttt{tikzpicture}, with its optional keys read first in
\texttt{/pgfplots/pourbaix}. It creates neither a picture nor an automatic plot.
Use \cs{addpourbaix}, \cs{addpourbaixwater} and ordinary \cs{addplot} commands in
its body. The \texttt{pourbaix/axis} presentation style is applied before your
options: axes are on top, with pH and potential labels. No axis limits or global
compatibility settings are imposed.

\begin{command}{\pourbaix\oarg{options}}
Begin the generic axis wrapper. An omitted optional argument means an empty
option list. Close it with \cs{endpourbaix}. The commands are installed inside
TikZ pictures, following the \cs{axis}/\cs{endaxis} convention.
In LaTeX, use \verb|\begin{pourbaix}[options]| and \verb|\end{pourbaix}|.
In ConTeXt MkIV, \cs{startpourbaix}\oarg{options} and \cs{stoppourbaix} are also
provided. All spellings use the same underlying axis and key processing.
\end{command}

\subsection{Key lookup and execution order}
Relative keys are looked up in \texttt{/pgfplots/pourbaix} first. PGFKeys'
\texttt{.search also} handler forwards unresolved keys to \texttt{/pgfplots},
including ordinary axis styles and their normal TikZ fallback. For example,
\texttt{build}, \texttt{declare}, \texttt{check}, \texttt{calculate}, \texttt{pH min},
\texttt{pH max} and \texttt{finalize} address the chemical family;
\texttt{xmin}, \texttt{width}, \texttt{title} and \texttt{axis background/.style}
address the underlying axis. Use \texttt{plot/domain palette=monochrome} for a
plot-family setting. A same-named chemical key takes precedence over an axis key;
use an absolute path such as \texttt{/pgfplots/title} to disambiguate. Unknown
keys remain errors. Nested declaration and plot families retain their own strict
validation; an unknown chemical record field is not reinterpreted as an axis key.

Options execute once, in their written order, during axis initialization.
No scientific stage is inferred or reordered. In particular, \texttt{check=Fe}
requires an existing declaration; \texttt{calculate=Fe} requires successful
checking; \texttt{finalize=Fe} requires calculation and explicit \texttt{pH min}
and \texttt{pH max} already in scope. Selection belongs here when a final diagram
is wanted. A calculated-only axis can instead plot provisional equations, and
an axis using an already selected diagram needs no preparation keys.
The selection interval is independent of \texttt{xmin}/\texttt{xmax}; the latter
only specify the view. Existing handle-based forms remain available.

Scientific aliases, result keys, styles and the search fallback obey the axis
scope. Declarations made here are available in the body and while the axis
finishes rendering; they do not replace outer aliases after the environment.
To reuse a diagram in several axes, declare and process it in an enclosing scope.
Options remain TeX tokens, forwarded without an extra expansion step. Ordinary
PGF handlers such as \texttt{/.expanded} and reusable \texttt{/.style} definitions
work as usual. Deferred rendering and immutable plot handles are unchanged.

\begin{codeexample}[width=9cm]
\begin{tikzpicture}
\begin{pourbaix}[
 declare={name=AxisFe,element=Fe,temperature=298.15,
   concentration=.01,
   species={Fe(s),Fe^2+,Fe^3+,Fe(OH)2(s),Fe(OH)3(s)},
   boundary convention=equal element concentrations,constants={}},
 check=AxisFe,calculate=AxisFe,
 pH min=0,pH max=14,finalize=AxisFe,
 plot/domain palette=tab10 light,
 xmin=0,xmax=14,ymin=-1.2,ymax=1.5,width=9cm,height=6cm]
 \addpourbaix{AxisFe};
 \addpourbaixwater[prepare={diagram=AxisFe},domain labels=false];
\end{pourbaix}
\end{tikzpicture}
\end{codeexample}
For plain LuaTeX, begin with \cs{tikzpicture} and \cs{pourbaix}; close with
\cs{endpourbaix} and \cs{endtikzpicture}.
In ConTeXt, begin with \cs{starttikzpicture} and \cs{startpourbaix}; close with
\cs{stoppourbaix} and \cs{stoptikzpicture}. The option list and body stay identical.

\section{Final diagrams, overlays and boundary labels}
\label{sec:rendering}
Selection assigns definitive pair identifiers and reconstructs domain strips.
The renderer clips that geometry to a finite viewport and submits ordinary plots.
``Final'' means final within the selected limiting construction, species list and
pH range; it does not add water stability limits or solve a full speciation problem.

\subsection{Definitive names and result lookup}
Retained redox pairs receive \texttt{1}, \texttt{2}, \ldots; retained vertical pairs
receive \texttt{a}, \texttt{b}, \ldots, \texttt{z}, \texttt{aa}, \ldots.
The two counters are independent and follow the original candidate inventory
order, skipping pairs with no retained segment. Names have no \texttt{-temp}
suffix. Several disconnected pieces of the same pair share one identifier;
their individual segment indices and pH intervals remain distinct. Changing the
selected pH range can change the set of retained pairs and therefore its numbering.

Every selected handle exports a new family below its existing data root:
\begin{quote}
\texttt{final/}\meta{identifier}\texttt{/}
\end{quote}
Its keys are \texttt{source}, \texttt{first species}, \texttt{second species},
\texttt{segment count}, \texttt{segment N} (an index into the segment array),
\texttt{equation type}, \texttt{equation}, \texttt{expression},
\texttt{a}, \texttt{b}, \texttt{pH}, \texttt{log K},
\texttt{standard potential}, and \texttt{equilibrium constant}.
Only applicable coefficients and constants are present. Units and equation
syntax are unchanged from the provisional readouts. \texttt{source} links to
the complete provisional record and its thermodynamic provenance.

Two selection-level counts report the number of final pairs and domain strips:
\texttt{final boundary count} and \texttt{domain count}.

Each \texttt{selection/domain N/} record contains
\texttt{species} (declaration index), \texttt{kind} (\texttt{existence} or
\texttt{predominance}), \texttt{pH min}, \texttt{pH max}, and optional
\texttt{lower source} and \texttt{upper source}. An absent bound is unbounded
in that direction. These are strips bounded by fixed equations, coalesced across
adjacent cells. One species can have several strips as its neighbors change.
They are geometric pieces, not extra chemical species or distinct phases.

\subsection{The plotting command}
Both \cs{addpourbaix} and \cs{addpourbaixwater} require a terminating semicolon.
Whitespace before it is allowed. A missing terminator is diagnosed before any
plot submission or water preparation, without consuming the following command.
The terminator ends the complete diagram request, which can generate several
ordinary PGFPlots plots. Appending path operations such as \texttt{node} is not
supported; use the domain and boundary label keys.
The long names \cs{pgfplotsaddpourbaix} and \cs{pgfplotsaddpourbaixwater} are
aliases with the same syntax. Calls without a semicolon must be updated.
The standalone picture and report commands do not require a semicolon.


\begin{command}{\addpourbaix\oarg{options}\marg{name};}
Inside a two-dimensional linear PGFPlots axis, plot an already selected diagram.
The optional key list belongs to \texttt{/pgfplots/pourbaix/plot/} and is local to
this invocation. The literal name is resolved immediately. No checking,
calculation or selection is performed implicitly; an incomplete diagram is an
error. The command changes neither name aliases nor \cs{pgfplotslastpourbaix}.
\end{command}

Supply explicit finite decimal \texttt{xmin}, \texttt{xmax}, \texttt{ymin} and
\texttt{ymax}, with increasing bounds. The renderer needs these bounds to clip
unbounded domains and vertical rays before submitting finite coordinates.
It does not choose a potential range on the user's behalf. Axis bounds given
through macros need explicit expansion, for example \texttt{xmin/.expanded}.
Logarithmic or three-dimensional axes are rejected. Reversed axis directions,
ordinary scaling and user-chosen physical dimensions are compatible with the
linear coordinate geometry. The selected pH interval limits the plotted data;
an axis extending beyond it can have blank margins.

The default output colors adjacent retained species differently and draws black
definitive boundaries, with species and boundary labels. The older phase-based
coloring remains available as an explicit option. Aqueous plot labels omit the redundant \texttt{(aq)}; solid labels retain
\texttt{(s)}. Domain labels use an interior area centroid; see Section~\ref{sec:appearance}.
Their text can still overlap another label in a narrow domain: placement is a
simple geometric choice, not a collision-avoidance algorithm. Local shifts,
anchors and fonts can be adjusted with the styles below.

\begin{codeexample}[width=8cm]
\pgfplotsset{pourbaix/declare={
 name=FinalFe,element=Fe,temperature=298.15,
 concentration=.01,
 species={Fe(s),Fe^2+,Fe^3+,Fe(OH)2(s),Fe(OH)3(s)},
 boundary convention=equal concentrations,
 constants={}},
 pourbaix/check=FinalFe,pourbaix/calculate=FinalFe,
 pourbaix/pH min=0,pourbaix/pH max=14,
 pourbaix/finalize=FinalFe}
\begin{tikzpicture}[
 pourbaix/boundary b label/.style={right}]
\begin{axis}[pourbaix/axis,width=8cm,height=6cm,
 xmin=0,xmax=14,ymin=-1,ymax=1]
 \addpourbaix{FinalFe};
\end{axis}
\end{tikzpicture}
\end{codeexample}
The empty input list in this example explicitly requests thermodynamic closure
with the default catalogue. Its constants can differ slightly from manually
supplied teaching values.

\subsection{Plot switches}
Every Boolean key accepts \texttt{true} or \texttt{false}; writing it without a
value means \texttt{true}. The following names are relative to the plot family.
They can also be set in an enclosing scope with \cs{pgfplotsset} using the full
\texttt{pourbaix/plot/} prefix.
\begin{center}
\begin{tabular}{>{\raggedright\arraybackslash}p{5cm}p{1.3cm}>{\raggedright\arraybackslash}p{7.5cm}}
Key & Initial & Meaning\\\hline
\texttt{boundaries} & true & Draw retained boundary pieces.\\
\texttt{boundary labels} & true & Append definitive identifiers to those plots.\\
\texttt{domain fills} & true & Fill final existence/predominance strips.\\
\texttt{domain labels} & true & Label visible final species.\\
\texttt{show removed boundaries} & false & Overlay the inactive portions of candidates.\\
\texttt{show removed domains} & false & Overlay diagnostic disproportionation bands and their labels.\\
\texttt{show provisional boundaries} & false & Overlay complete provisional lines within the selected pH range and viewport.\\
\texttt{provisional labels} & false & Add provisional names to either boundary overlay.\\
\texttt{show removed} & --- & Set both removed-boundary and removed-domain switches together.
\end{tabular}
\end{center}
The full-provisional and removed-only boundary overlays are independent; enabling
both can draw some portions twice. A degenerate equation has no line to draw.

\begin{key}{/pgfplots/pourbaix/plot/handle=\meta{handle}}
An optional captured numeric handle, initially empty. When nonempty it takes
precedence over the command's name argument, which may then be empty. For example,
\texttt{handle/.expanded=\string\saved} supports an anonymous diagram or a
snapshot whose name has since been rebound. The handle must identify a selected
snapshot. Styles and viewport settings are still taken from the current invocation.
\end{key}

\subsection{Eliminated domains: a diagnostic overlay}
A species removed by disproportionation has no common stable potential interval.
The optional red band shows the interval \emph{between its incompatible bounds},
not a stable domain. Its label explicitly says \texttt{incompatible domain}.
The diagnostic representation makes the inconsistency visible without
inventing an existence or predominance region. The label is placed three quarters
of the way from the lower visible band limit to the upper one at an interior pH,
to reduce coincidence with the replacement line.

Bands come from explicit \texttt{disproportionation} events, including transient
steps of a cascade; they may overlap. Zero-width events have no band to fill.
Candidates excluded solely by acid/base selection or nonadjacency have no
invented deleted-species region. Their inactive boundary pieces remain available
through the removed-boundary overlay and result inventory.

\begin{codeexample}[width=9cm]
\pgfplotsset{pourbaix/declare={
 name=OverlayCu,element=Cu,temperature=298.15,
 concentration=.01,species={Cu(s),Cu^+,Cu^2+},
 boundary convention=equal concentrations,
 constants={}},
 pourbaix/check=OverlayCu,
 pourbaix/calculate=OverlayCu,
 pourbaix/pH min=0,pourbaix/pH max=14,
 pourbaix/finalize=OverlayCu}
\begin{tikzpicture}
\begin{axis}[pourbaix/axis,width=9cm,height=6cm,
 xmin=0,xmax=14,ymin=0,ymax=.6]
 \addpourbaix[show removed,
  provisional labels,boundary label style={pos=.25},
  removed boundary style={red,dashed},
  removed domain style={fill=red!20,fill opacity=.5},
  removed domain label style={text=red,font=\small}
 ]{OverlayCu};
\end{axis}
\end{tikzpicture}
\end{codeexample}
The labels with \texttt{-temp} refer to retired candidates. The single definitive
boundary is labeled \texttt{1}. This is the same pair identifier used by the
standalone map below, although diagram handles are independent.

\subsection{Styles and appended label nodes}
These keys store TeX option lists, with the initial values listed below.
Assign an option list directly to the key. It is copied into each plot or node
before deferred visualization, without Lua interpretation.
\begin{center}
\begin{tabular}{p{6.5cm}p{7.5cm}}
Key & Initial option list\\\hline
\texttt{final boundary style} & \texttt{black,semithick}\\
\texttt{removed boundary style} & \texttt{red!70!black,dashed,thin}\\
\texttt{provisional boundary style} & \texttt{gray,densely dotted,thin}\\
\texttt{existence domain style} & \texttt{fill=blue!10}\\
\texttt{predominance domain style} & \texttt{fill=green!8}\\
\texttt{removed domain style} & \texttt{fill=red!18,fill opacity=.55}\\
\texttt{boundary label style} & empty\\
\texttt{final domain label style} & empty\\
\texttt{removed domain label style} & \texttt{text=red!70!black}
\end{tabular}
\end{center}
Fills start with \texttt{draw=none,mark=none}; boundary plots start with
\texttt{mark=none}. All generated plots use \texttt{forget plot}, so they do not
create legend entries or consume the user's plot-style cycle. The usual
\cs{addlegendimage} and \cs{addlegendentry} remain available for a custom legend.
For superpositions, fills may be disabled or given opacity and pattern options;
see Section~\ref{sec:appearance}. Base phase fills are overridden by the adjacency
palette unless \texttt{domain coloring=phase} is requested.

\begin{stylekey}{/tikz/pourbaix/boundary label}
Style for the actual node appended after a boundary's \cs{addplot} coordinates.
Initially \texttt{midway,fill=white,inner sep=1.5pt}. It can also be used in a
hand-written plot. Override it using ordinary \cs{tikzset}; a call's
\texttt{boundary label style} option list is applied afterwards.
\end{stylekey}
\begin{stylekey}{/tikz/pourbaix/domain label}
Base style for species-name nodes, initially \texttt{inner sep=2pt}. The relevant
final or removed domain label option list follows it.
\end{stylekey}

Optional styles provide individual overrides. Define them with \texttt{.style}
or \texttt{.append style}; they have no effect until defined:
\begin{description}
\item[\texttt{/pgfplots/pourbaix/boundary }\meta{id}]
  Applied after the boundary option list. Use a definitive identifier for a final
  line or a provisional name for an overlay. It also applies with labels disabled.
\item[\texttt{/tikz/pourbaix/boundary }\meta{id}\texttt{ label}]
  Applied last to that line's appended label node; useful for \texttt{pos},
  \texttt{above}, \texttt{sloped} or a local shift.
\item[\texttt{/pgfplots/pourbaix/species }\meta{index}]
  Applied last to that species' region fill, including its diagnostic band.
\item[\texttt{/tikz/pourbaix/species }\meta{index}\texttt{ label}]
  Applied last to the species-name node, including its diagnostic label.
\end{description}
Species indices always refer to declaration order, not oxidation order. Set
named node styles in a scope that lasts until axis visualization; per-call
option lists are captured, but macros inside them retain normal TeX semantics.
The library does not globally freeze or redefine user macros.

\begin{codeexample}[width=7cm]
\begin{tikzpicture}
\begin{axis}[pourbaix/axis,width=7cm,height=4cm,
 xmin=0,xmax=14,ymin=-1,ymax=1]
 \addplot[blue,domain=0:14,samples=2] {-.05*x+.4}
   node[pourbaix/boundary label,pos=.65,above] {1};
 \addplot[red] coordinates {(4,-1) (4,1)}
   node[pourbaix/boundary label,pos=.3,right] {a};
\end{axis}
\end{tikzpicture}
\end{codeexample}
These two lines illustrate the reusable node style; they are not asserted to be
a chemically complete diagram. Generated boundary plots attach their nodes in
the same manner, rather than adding unrelated labels after the axis has rendered.

\subsection{Geometry and deferred rendering}
Final domain strips are clipped polygons between their lower and upper affine
bounds, with the viewport supplying missing limits. Touching grid cells with
identical bounding equations were already merged during selection; the renderer
does not paint one rectangle per sample. Polygons close with the TikZ path
operation \texttt{--cycle}, not a fill-to-axis-baseline operation. A boundary
piece is clipped as a straight line segment. Vertical rays receive finite
endpoints from the viewport. No interpolation changes its chemical equation.

Plot submission order is final fills, incompatible bands, complete provisional
lines, inactive boundary pieces, final boundaries, final species names, and
incompatible-band names. This makes final boundaries legible above the overlays.
The convenience \texttt{pourbaix/axis} style sets \texttt{axis on top} so filled
regions do not cover the frame; the user can override that ordinary PGFPlots key.
No global compatibility level, font, axis limit or externalization is enabled.

Name resolution and clipping happen during survey. Geometry is captured before
any later declaration can rebind a name. The renderer uses ordinary \cs{addplot}
commands throughout, including invisible one-point plots carrying species-name
nodes. It does not mutate PGFPlots' private plot lists. PGFPlots performs final
drawing at axis end, alongside other user plots. Rendering does not create or
modify diagram snapshots. Invalid geometry/options are rejected before any
plot commands are submitted; invalid user-supplied TeX styles remain subject to
PGF's normal diagnostics during plot processing.

\subsection{PGFPlots presentation}
\begin{key}{/pgfplots/pourbaix/axis}
Optional style with no value. Sets the axis labels to pH and $E$/V. It neither
converts data nor chooses limits. Later PGFPlots options can override either
label, and \texttt{axis on top}. Applying the style alone does not finalize or
render a declared diagram.
\end{key}
This example checks presentation only; the horizontal line is an arbitrary
coordinate plot, not a calculated chemical boundary.
\begin{codeexample}[width=6cm]
\begin{tikzpicture}
  \begin{axis}[pourbaix/axis,
    width=5.8cm,height=5cm,
    xmin=0,xmax=14,ymin=-1,ymax=1]
    \addplot coordinates {(0,0) (14,0)};
  \end{axis}
\end{tikzpicture}
\end{codeexample}

\subsection{Boundary presentation by final identifier}
\begin{key}{/pgfplots/pourbaix/plot/boundary/\meta{identifier}=\meta{options}}
Customize a retained pair using its final number or lowercase-letter identifier.
The options are \texttt{style}, \texttt{label}, and \texttt{label style}, all with
required values. Omitted styles are empty and an omitted label retains the
identifier. Label text is arbitrary TeX; include math delimiters explicitly.
An empty label is allowed. Styles accept ordinary PGFPlots/TikZ line options or
TikZ node options, respectively. For placement along a segment use node options
such as \texttt{pos=.65,above}. The global \texttt{boundary labels=false} still
suppresses nodes, even when a custom label exists.

These overrides apply after the global boundary styles and older indexed styles.
They are captured as TeX tokens before deferred visualization and apply to every
retained piece of the pair. Unknown final identifiers and duplicate customization
records are errors, even when the boundary is outside the viewport. A final
identifier addresses only final pieces; diagnostic provisional pieces retain
separate styles. Scope options to each plotting call when overlaying diagrams.

Changing the label does not rename the boundary, change result-key paths or modify
the species map. The map continues to show chemical pair identifiers. The older
\texttt{pourbaix/boundary N} and \texttt{pourbaix/boundary N label} styles remain
supported. A named node is possible, but avoid reusing a single node name when
one boundary has several retained pieces.
\end{key}

\section{Final-only views and domain presentation}
\label{sec:appearance}
Presentation is independent of chemistry. Checking, equation calculation and
selection remain explicit operations; none of the options in this section changes
the selected snapshot. A final-only view shows retained boundaries and species,
without printing classifications, intermediate equations or a report.

\begin{command}{\pgfplotspourbaix\oarg{axis options}\marg{name}}
Compose an ordinary TikZ picture and PGFPlots axis.
The plot uses \texttt{final diagram}. Supply the same explicit
finite decimal axis limits required by the regular renderer. The named diagram
must already be selected. No fonts, physical dimensions or axis limits are imposed
by this shortcut. The generic syntax works in all three formats. Use
\texttt{pourbaix/plot/}\meta{key} within the axis options to configure domains.
For multiple diagrams or additional plots, use an ordinary axis instead.
\end{command}
\begin{stylekey}{/pgfplots/pourbaix/plot/final diagram}
Disable removed-domain, removed-boundary and full provisional-boundary overlays.
Other options (fills, labels and styles) are retained. The style also applies to
ordinary plotting calls inside an existing axis.
\end{stylekey}

The following setup is executed inline and shared by this section's examples.
It deliberately uses the sourced library rather than adding thermodynamic data.
\begin{verbatim}
\pgfplotsset{pourbaix/declare={name=StyledFe,element=Fe,
  species={Fe(s),Fe^2+,Fe^3+,Fe(OH)2(s),Fe(OH)3(s)},
  temperature=298.15,concentration=.01,
  boundary convention=equal concentrations,constants={}},
  pourbaix/check=StyledFe,pourbaix/calculate=StyledFe,
  pourbaix/pH min=0,pourbaix/pH max=14,pourbaix/finalize=StyledFe}
\end{verbatim}
\pgfplotsset{pourbaix/declare={name=StyledFe,element=Fe,
  species={Fe(s),Fe^2+,Fe^3+,Fe(OH)2(s),Fe(OH)3(s)},
  temperature=298.15,concentration=.01,
  boundary convention=equal concentrations,constants={}},
  pourbaix/check=StyledFe,pourbaix/calculate=StyledFe,
  pourbaix/pH min=0,pourbaix/pH max=14,pourbaix/finalize=StyledFe}

\begin{codeexample}[width=8cm]
\pgfplotspourbaix[
 width=8cm,height=6cm,
 xmin=0,xmax=14,ymin=-1.2,ymax=1.2,
 pourbaix/plot/domain label mode=letters
]{StyledFe}
\end{codeexample}
Letters follow retained species in declaration order, including species outside
the current viewport. Here $A$ is Fe(s), $B$ is Fe(II), $C$ is Fe(III), $D$ is the
ferrous hydroxide and $E$ is the ferric hydroxide. These mathematical uppercase
letters identify domains; lowercase letters and numbers still identify boundaries.

\subsection{Fills, adjacency and superposition}
All following plot keys belong to \texttt{/pgfplots/pourbaix/plot/} and are local
to a plotting call when supplied in its optional argument.
\begin{key}{/pgfplots/pourbaix/plot/domain coloring=\meta{mode} (initially adjacent)}
The required value is \texttt{adjacent}, \texttt{patterns}, \texttt{phase}, or \texttt{none}.
The first uses a deterministic saturation-degree graph coloring: each retained
species is a vertex and each retained chemical pair an edge. Ties use degree,
then original species order. Every adjacent pair receives different palette
indices, including pairs outside the viewport. This is a deterministic heuristic,
not a minimum-color solver. \texttt{patterns} applies the same assignment to
the pattern list instead of the color list. \texttt{phase} uses blue solid/green
aqueous base styles. \texttt{none} omits stable-domain fills. The existing
\texttt{domain fills=false} also disables them, regardless of this mode; bare
\texttt{domain fills} means true. Diagnostic bands have their own switch.
\end{key}
\begin{key}{/pgfplots/pourbaix/plot/domain colors=\meta{color list}}
A required comma-separated list of ordinary TikZ color expressions, initially
\texttt{blue!15}, \texttt{orange!20}, \texttt{green!15}, \texttt{violet!15},
\texttt{yellow!25}, \texttt{red!15}. Empty entries and identical text entries are
errors when adjacency fills are requested. Too few colors for the computed
coloring cause a clear error before plot submission; the palette never silently
cycles onto a neighbor. Supply more colors if needed. Distinct TeX color names
can denote the same visual color: ensuring a visually distinct custom palette
remains the user's responsibility. Use \texttt{/.expanded} for a macro containing
an entire list. Zooming alone never changes color indices.
\end{key}
\subsection{Named palettes and monochrome printing}
\begin{key}{/pgfplots/pourbaix/plot/domain palette=\meta{name}}
Select a named palette in the current plot or TeX scope. A value is required;
unknown names are errors. There is no implicit named palette: the original
six-color list remains the default. Each preset sets \texttt{domain coloring}
and the corresponding list. Later keys override earlier ones as usual.
The presets do not enable disabled fills, alter chemistry or change axis limits.
\begin{center}
\begin{tabular}{lll}
Name & Slots & Appearance\\\hline
\texttt{tab10} & 10 & Matplotlib/Tableau categorical RGB colors\\
\texttt{tab10 light} & 10 & The same colors, 25\% color and 75\% white\\
\texttt{grayscale} & 4 & Black at 8\%, 45\%, 22\%, 65\%\\
\texttt{patterns} & 6 & Diagonals, dots, horizontal/vertical lines, crosshatch\\
\texttt{hatching} & 5 & The same line patterns, without dots\\
\texttt{monochrome} & 6 & Alias for \texttt{patterns}
\end{tabular}
\end{center}
The RGB values and order of \texttt{tab10} follow Matplotlib v3.10.7's
\href{https://github.com/matplotlib/matplotlib/blob/v3.10.7/lib/matplotlib/_color_data.py}{Tableau color definitions}.
The light variant is a library tint, not a separate Matplotlib palette.
Grayscale and pattern selections are library choices. Patterns are the stock
\href{https://tikz.dev/library-patterns}{TikZ patterns}; the library loads
\texttt{patterns} automatically, in all three formats. No Matplotlib installation
or runtime download is involved.
\end{key}
\begin{key}{/pgfplots/pourbaix/plot/domain patterns=\meta{pattern list}}
A required comma-separated list of stock TikZ pattern names, initially
\texttt{north east lines}, \texttt{north west lines}, \texttt{dots},
\texttt{horizontal lines}, \texttt{vertical lines}, \texttt{crosshatch}.
Used when \texttt{domain coloring=patterns}; setting the list alone does not
switch modes. The \texttt{hatching} preset removes \texttt{dots} from this order.
Empty or repeated entries and insufficient capacity are errors before plot
submission. Pattern names are interpreted by TikZ; unknown names produce TikZ
errors. Parameterized pattern expressions containing commas are not supported
by this simple list interface.

The same deterministic adjacency assignment used for colors selects pattern
slots. Every retained adjacent pair gets different patterns, even outside the
viewport; zooming leaves assignments unchanged. Slots may be reused for
nonadjacent domains. Patterns have black ink and transparent backgrounds,
so they also work for overlays. Use \texttt{domain style} to change
\texttt{pattern color} or \texttt{fill opacity}. Styles remain TeX tokens and
are captured into each deferred addplot; no style becomes Lua source.
\end{key}
For black-and-white printing, prefer \texttt{monochrome} and opaque black ink.
A color palette can collapse to similar grays when printed. The grayscale
preset offers distinct tones, but thin domains and printer screening may make
patterns easier to identify. The tinted palette is useful behind black labels;
for dense patterns or saturated colors, give label nodes a white background.
Transparency and overlapping patterns can reduce contrast. Adjacency is checked
within each diagram, not between regions belonging to separate overlays.
As with colors, explicit common or per-domain fill/pattern overrides can defeat
the automatic distinction; the library does not interpret arbitrary TikZ styles.

\begin{codeexample}[width=14cm]
\pgfplotsset{pourbaix/declare={name=PaletteFe,element=Fe,
 species={Fe(s),Fe^2+,Fe^3+,Fe(OH)2(s),Fe(OH)3(s)},
 temperature=298.15,concentration=.01,
 boundary convention=equal concentrations,constants={}},
 pourbaix/check=PaletteFe,pourbaix/calculate=PaletteFe,
 pourbaix/pH min=0,pourbaix/pH max=14,pourbaix/finalize=PaletteFe}
\vbox{\hsize=14cm\noindent
\foreach \palette [count=\paletteindex] in
 {tab10,tab10 light,grayscale,patterns,hatching}{
 \begin{tikzpicture}
 \begin{axis}[pourbaix/axis,width=6cm,height=4.5cm,
   xmin=0,xmax=14,ymin=-1.2,ymax=1.5,title={\palette}]
  \addpourbaix[domain palette/.expanded=\palette,
    final domain label style={fill=white,inner sep=1pt}]{PaletteFe};
 \end{axis}
 \end{tikzpicture}%
 \ifodd\paletteindex\hskip1cm\else\par\noindent\fi
}
}
\end{codeexample}

\begin{key}{/pgfplots/pourbaix/plot/domain style=\meta{option list} (initially empty)}
Common PGFPlots/TikZ fill options, such as \texttt{fill opacity=.3},
\texttt{pattern=north east lines}, and \texttt{pattern color=red!60!black}.
They apply to final fills and enabled diagnostic bands, after the phase/removed
base style and the adjacency color (stable domains only). The existing indexed
species style follows, then the per-species \texttt{domain/style} override.
Explicit fill overrides can intentionally supersede adjacency coloring.
Transparency applies to fills without fading boundary or label nodes unless
requested separately. The stock TikZ \texttt{patterns} library is loaded automatically;
custom pattern libraries follow the normal TikZ loading rules.
\end{key}

The next example overlays iron with the simple three-state copper system. Copper
adds red hatching and a red boundary; its species nodes are suppressed here to
keep this superposition legible. Each diagram retains its own numbering, colors,
chemical data and immutable handle. Hatching is a rendering option, not an
additional chemical phase or a mixed-element equilibrium calculation.
\begin{codeexample}[width=8cm]
\usetikzlibrary{patterns}
\pgfplotsset{pourbaix/declare={
 name=OverlayCu,element=Cu,
 species={Cu(s),Cu^+,Cu^2+},
 temperature=298.15,concentration=.01,
 boundary convention=equal concentrations,
 constants={}},
 pourbaix/check=OverlayCu,
 pourbaix/calculate=OverlayCu,
 pourbaix/pH min=0,pourbaix/pH max=14,
 pourbaix/finalize=OverlayCu}
\begin{tikzpicture}
\begin{axis}[pourbaix/axis,
 width=8cm,height=6cm,
 xmin=0,xmax=14,ymin=-1.2,ymax=1.2]
 \addpourbaix[
  domain label mode=letters,
  domain style={fill opacity=.45}
 ]{StyledFe};
 \addpourbaix[
  final diagram,domain labels=false,
  boundary labels=false,
  domain style={pattern=north east lines,
    pattern color=red!20,
    fill opacity=.25},
  final boundary style={red!70!black,thick}
 ]{OverlayCu};
 % Repeat only iron's nodes above the overlay.
 \addpourbaix[
  domain fills=false,boundaries=false,
  domain label mode=letters
 ]{StyledFe};
\end{axis}
\end{tikzpicture}
\end{codeexample}
For an outline-only diagram use \texttt{domain coloring=none} or
\texttt{domain fills=false}. For a monochrome hatched overlay, use
\texttt{domain coloring=phase} with a common pattern style; the pattern then
carries the visual distinction between diagrams. Domain fills still require
positive visible area. Touching pieces have no stroked polygon edges.

\subsection{Letters, selected labels and custom TikZ nodes}
\begin{key}{/pgfplots/pourbaix/plot/domain label mode=\meta{mode} (initially species)}
Required value \texttt{species} or \texttt{letters}. The former uses chemical
formulas. The latter assigns $A,B,\ldots,Z,AA,AB,\ldots$ in mathematical mode to
retained species in declaration order, independently of clipping. All disconnected
parts of one species share its letter. Never-retained species keep their formula
in diagnostic overlays unless explicitly listed below. The existing Boolean
\texttt{domain labels} (default true, bare value true) controls final node emission.
\end{key}
\begin{key}{/pgfplots/pourbaix/plot/letter species=\meta{species list}}
Require an explicit list, enable letter mode, and assign $A,B,\ldots$ only to
those species, in the supplied order. All other species keep their formula.
Use the declaration's plain chemical notation, not TeX typesetting commands.
Unknown species and duplicate entries are errors. An empty list selects the
all-retained-species behavior. Internally this style sets \texttt{domain label
mode=letters} and \texttt{letter species list=}\meta{list}; the latter is also a
value key, initially empty, and does not independently enable letter mode.
\end{key}
\begin{key}{/pgfplots/pourbaix/plot/domain/\meta{species}=\meta{key list}}
Customize the declared species identified by the final path component, for example
\texttt{domain/Fe(s)=\{label=\$M\$\}}. This is the preferred syntax.
An option list is required; an empty list is allowed. The species uses the same
chemical notation as the declaration, including parentheses and caret charges.
Give at most one record per species per effective
plot configuration; duplicates are errors. Records inherit TeX grouping and are
best supplied directly to a plotting call when superposing different diagrams.
Nested keys belong to \texttt{/pgfplots/pourbaix/domain/}:
\begin{description}
\item[Chemical identifier] The species comes from the key path and is normalized
with the declaration parser. Unknown species are errors even if currently clipped.
A nested \texttt{species} key is forbidden in this form: it cannot redirect a label.
\item[\texttt{label=}\meta{TeX text}] Override the formula or generated letter.
There is no automatic math wrapper: write \texttt{label=\{\$X\$\}} for a mathematical
label, or supply arbitrary text and macros. An omitted key retains the automatic
label; empty text is allowed. A diagnostic label still appends ``incompatible
domain'' so the override cannot accidentally present it as a stable region.
\item[\texttt{label at=\{pH,E\}}] Two finite decimal data coordinates, without
parentheses; pH is dimensionless and $E$ is in volts. An empty value restores
automatic placement. Explicit positions may lie outside the chemical region;
axis clipping still applies. Use \texttt{/.expanded} for macro-supplied coordinates.
\item[\texttt{label style=\{options\}}] TikZ node options, applied last.\par
For example \texttt{anchor=west}, \texttt{xshift=3pt}, \texttt{font=\string\small},
\texttt{draw,fill=white}, \texttt{name=my-ion}, \texttt{text opacity=.8}, or
\texttt{rotate=15}. This is an actual node appended to an ordinary plot, not an
independent eagerly drawn object. Named nodes are available after axis completion.
\item[\texttt{style=\{options\}}] Per-species polygon options, applied last to
final and diagnostic fills of that species. They may override palette colors,
opacity or patterns. Omitted style values are empty.
\end{description}
All nested keys require values. Text and style tokens stay in TeX; they never
become Lua source. They are copied into the deferred plot command, so internal
scratch macros need not survive until visualization. User macros contained in
those token lists retain ordinary TeX expansion/lifetime semantics. Apply styles
in a scope that includes axis completion.
\end{key}

\begin{key}{/pgfplots/pourbaix/plot/domain=\meta{key list}}
Compatibility form: supply the required \texttt{species=}\meta{formula} field
inside the list, followed by the same label, position and style keys. For example,
\texttt{domain=\{species=Fe(s),label=\$M\$\}} is equivalent to the species-addressed
form above. Both forms share the same scoping, validation and duplicate checks.
\end{key}

\begin{codeexample}[width=8cm]
\begin{tikzpicture}
\begin{axis}[pourbaix/axis,
 width=8cm,height=6cm,
 xmin=0,xmax=14,ymin=-1.2,ymax=1.2]
 \addpourbaix[
  domain fills=false,
  letter species={Fe^2+,Fe^3+},
  domain/Fe(s)={label={$M$},
    label at={4,-.9},
    label style={draw,rounded corners,
      inner sep=3pt,name=metal}},
  domain/Fe^2+={
    label at={4,.15},
    label style={text=blue}},
  domain/Fe(OH)2(s)={
    label style={font=\small,xshift=-3pt}}
 ]{StyledFe};
\end{axis}
\draw[->] (metal) -- ++(0,-.4);
\end{tikzpicture}
\end{codeexample}
Only Fe(II) and Fe(III) become $A$ and $B$; the metal has the independent custom
label $M$. This figure also demonstrates that the label is a reusable named TikZ
node. Both formula and letter nodes accept the same styles.

\subsection{Automatic interior placement}
\begin{key}{/pgfplots/pourbaix/plot/domain label placement=\meta{mode} (initially centroid)}
Required value \texttt{centroid} or \texttt{largest piece}. The default takes the
area-weighted centroid of all visible polygon strips for one species. If that
point belongs to none of its strips (possible for concave or disconnected domains),
use the area centroid of the largest visible strip. The second mode always uses
that strip. Polygon moments and containment tests use normalized viewport
coordinates; the containment roundoff threshold is $64\,2^{-52}$ in normalized
cross products. A species with no positive-area visible region receives no node.
A per-species \texttt{label at} overrides this placement without creating a domain.
\end{key}
The automatic point is geometrically inside the visible domain, but a long text
box may cross a boundary or be clipped. There is no automatic text-fit or collision
solver. Adjust font, anchor, shifts, or explicit coordinates when necessary.
Only one final node is emitted per visible species; disconnected regions share
that node by default. A zoom changes the visible centroid, not the species' letter
or palette index. Diagnostic bands retain their dedicated placement near the
upper part of the band; they are not included in stable-domain centroid moments.

\section{The water stability overlay}

The water overlay adds the two equilibria
\[
 2\mathrm{H}^+ + 2\mathrm e^- = \mathrm{H}_2(\mathrm g),
 \qquad
 \mathrm{O}_2(\mathrm g)+4\mathrm{H}^++4\mathrm e^-
       = 2\mathrm{H}_2\mathrm O(\mathrm l).
\]
It has three redox regions: hydrogen below the lower line, water between the
lines, and oxygen above the upper line. Acidic and alkaline descriptions of
water do not introduce additional redox regions: protons, hydronium and
hydroxide are distinct acid/base species, but do not change the corresponding
oxidation levels of hydrogen and oxygen in water. This simplified overlay does
not include peroxide or other oxygen/hydrogen species.

\subsection{Gas activities and the Nernst equations}
For an ideal gas, fugacity equals partial pressure. Here
$a_i=p_i/p^\circ$, with $p^\circ=1\,\mathrm{bar}$ fixed. Each pressure is an
\emph{absolute partial pressure}, not a total pressure or a concentration.
The default $p_{\mathrm H_2}=p_{\mathrm O_2}=1\,\mathrm{bar}$ represents two
separate equilibria; it does not describe a two-component mixture at a total
pressure of one bar. For wet gases, supply the relevant partial pressure after
accounting for water vapour; the library does not perform that correction.

With $f(T)=RT\ln(10)/F$ by default (or the specified Nernst factor),
unit liquid-water activity and potentials versus the
standard hydrogen electrode (SHE), the implemented equations are
\begin{align*}
 E_{\mathrm H_2}&=-f(T)\,\mathrm{pH}
                  -\tfrac12 f(T)\log_{10}(p_{\mathrm H_2}/p^\circ),\\
 E_{\mathrm O_2}&=E^\circ_{\mathrm O_2/\mathrm H_2O}(T)-f(T)\,\mathrm{pH}
                  +\tfrac14 f(T)\log_{10}(p_{\mathrm O_2}/p^\circ).
\end{align*}
Thus both slopes are approximately $-0.05915935\,\mathrm{V}$ per pH unit at
298.15 K. Hydrogen's standard potential is zero by the SHE convention.
The oxygen default is the rounded value $1.229\,\mathrm V$ at 298.15 K only.
No temperature extrapolation is attempted. Nonideal gases, nonunit water
activity, reference-electrode conversions and kinetic overpotentials are
outside this model. Use the same SHE reference for any superposed metal diagram.

The default is sourced from Malkow, Pilenga and Blagoeva,
\emph{EU harmonised terminology for hydrogen generated by electrolysis}
(2021), JRC120120, DOI \href{https://doi.org/10.2760/732809}{10.2760/732809},
entry 720, printed page 105; entry 1069 defines SATP as 298.15 K and 1 bar.
The gas convention follows the
\href{https://goldbook.iupac.org/terms/view/I02935}{IUPAC definition of an ideal gas}.
These citations document equilibrium conventions, not electrolysis onset rates.

\subsection{Preparation, data and scope}
\begin{key}{/pgfplots/pourbaix/water/prepare=\marg{options}}
Validate and calculate an independent water data record. Options are parsed under
\texttt{/pgfplots/pourbaix/water/input}. Every call starts from the defaults below;
previous preparation options are not inherited. A successful call replaces the
water record in the current TeX group; failure preserves the previous record.
No element-diagram handle or species inventory is changed. Preparation is also
accepted inside the plotting command, where its data remain local to that call.
\end{key}
\begin{key}{/pgfplots/pourbaix/water/input/temperature=\meta{kelvins}}
Positive finite absolute temperature; no implicit default. Alternatively inherit
it with \texttt{diagram}. If both are supplied, they must agree numerically.
\end{key}
\begin{key}{/pgfplots/pourbaix/water/input/diagram=\meta{name}}
Read the temperature and effective Nernst factor from an existing diagram,
even before its selection stage. The name must exist in the current TeX scope.
No concentrations, species or pressure settings are inherited.
\end{key}
\begin{key}{/pgfplots/pourbaix/water/input/nernst factor=\meta{auto or volts}}
If omitted, inherit the associated diagram's effective factor; without a
\texttt{diagram}, compute it from temperature. Explicit \texttt{auto} computes
$RT\ln(10)/F$. A positive finite decimal overrides this calculation for a
standalone water overlay. With \texttt{diagram}, any explicit setting must agree
numerically with its factor, otherwise preparation fails and preserves the
previous water record. This prevents mixing approximations within an associated
element/water plot. Temperature is still required or inherited, and the oxygen
standard-potential default remains restricted to 298.15 K.
\end{key}
\begin{key}{/pgfplots/pourbaix/water/input/hydrogen pressure=\meta{bar} (initially 1)}
Positive finite hydrogen partial pressure in bar. Zero is an error, since the
Nernst logarithm is undefined there. Bare decimal numbers are required.
\end{key}
\begin{key}{/pgfplots/pourbaix/water/input/oxygen pressure=\meta{bar} (initially 1)}
Independent positive finite oxygen partial pressure in bar.
\end{key}
\begin{key}{/pgfplots/pourbaix/water/input/oxygen standard potential=\meta{volts}}
Finite standard reduction potential versus SHE at the requested temperature,
with gas standard pressure 1 bar and water activity 1. Defaults to 1.229 V only
at 298.15 K; missing data at any other temperature produce an error.
\end{key}

The read-only convention for the generated keys is
\texttt{/pgfplots/pourbaix/water/data/}. Suffixes are \texttt{temperature}, \texttt{nernst factor} (V),
\texttt{hydrogen pressure}, \texttt{oxygen pressure},
\texttt{oxygen standard potential}, \texttt{standard pressure} (1 bar),
\texttt{reference} (SHE), \texttt{a} (common slope), \texttt{hydrogen b} and
\texttt{oxygen b} (intercepts in volts), and \texttt{hydrogen expression} and
\texttt{oxygen expression} (PGF mathematical expressions in \texttt{x}).
Use \verb|\pgfkeysgetvalue| to retrieve them. Preparation errors include
nonfinite values and parameters for which the upper equilibrium is no longer
strictly above the lower one; the three-region model requires an open water
stability interval.

\subsection{Plotting and appearance}
\begin{command}{\addpourbaixwater\oarg{options};}
Submit clipped ordinary plots inside a two-dimensional linear axis with explicit
finite decimal \texttt{xmin}, \texttt{xmax}, \texttt{ymin}, \texttt{ymax}.
Prepare first, or supply \texttt{prepare=\{...\}} in these options.
The optional list uses \texttt{/pgfplots/pourbaix/water}. There is no mandatory
argument. Plot coordinates and all node/style tokens are captured before axis
visualization, so later preparation does not change an earlier overlay.
\end{command}
\begin{key}{/pgfplots/pourbaix/water/domain fills=\meta{boolean} (initially false)}
Enable the three clipped fills. The defaults use orange, blue and green,
with opacity 0.15. Draw fills before other opaque plots if they should be behind
those plots. Plot order remains ordinary PGFPlots order.
\end{key}
\begin{key}{/pgfplots/pourbaix/water/domain labels=\meta{boolean} (initially true)}
Enable one TikZ node per visible region. Its default position is the mean of the
vertices of the convex clipped polygon, hence inside that region. Styles can
adjust the anchor or shift the node; labels are not automatically collision-free.
\end{key}
\begin{key}{/pgfplots/pourbaix/water/boundary style=\marg{options}}
Common line options, initially \texttt{blue!70!black,dashed,thick}.
\end{key}
\begin{key}{/pgfplots/pourbaix/water/domain style=\marg{options}}
Common fill options, initially \texttt{fill opacity=.15}. Ordinary TikZ patterns
are supported after loading the corresponding TikZ library.
\end{key}
\begin{key}{/pgfplots/pourbaix/water/label style=\marg{options}}
Common domain-node options, initially empty.
\end{key}
\begin{key}{/pgfplots/pourbaix/water/domain/\meta{id}=\marg{options}}
The identifiers are \texttt{H2}, \texttt{H2O}, \texttt{O2}. Nested \texttt{style},
\texttt{label} and \texttt{label style} set the fill, text and node options.
Local styles follow common styles; a local opacity therefore overrides the common
one. Labels default to the chemical names and may be empty or arbitrary TeX.
\end{key}
\begin{key}{/pgfplots/pourbaix/water/boundary/\meta{id}=\marg{options}}
The identifiers \texttt{H2} and \texttt{O2} denote the respective equilibria.
Nested \texttt{style}, \texttt{label} and \texttt{label style} are initially empty.
A nonempty label produces a node appended to that boundary's \verb|\addplot|,
using \texttt{pourbaix/boundary label} followed by the local node options.
These identifiers never consume the element diagram's numerical labels.
\end{key}
All options obey TeX grouping. Text and appearance stay as TeX tokens, never Lua
source. Numeric inputs cross the bridge as unexpanded macro replacement text;
use PGF's \texttt{.expanded} handler explicitly when expansion is needed.

\subsection{An independent water diagram}
\begin{codeexample}[width=8cm]
\begin{tikzpicture}
\begin{axis}[pourbaix/axis,width=8cm,height=6cm,
 xmin=0,xmax=14,ymin=-1.2,ymax=1.5]
 \addpourbaixwater[
  prepare={temperature=298.15},
  domain fills=true,
  domain/H2O={label style={fill=white}},
  boundary/O2={label=$E_{\mathrm O_2}$,
    label style={above,pos=.75}}];
\end{axis}
\end{tikzpicture}
\end{codeexample}

\subsection{Superposition on an iron diagram}
The water layer does not remove metal regions outside water's stability window.
Such regions can be useful for discussing corrosion, metastability and kinetics.
\begin{codeexample}[width=8cm]
\pgfplotsset{pourbaix/declare={
 name=WaterFe,element=Fe,temperature=298.15,
 concentration=.01,
 species={Fe(s),Fe^2+,Fe^3+,Fe(OH)2(s),Fe(OH)3(s)},
 boundary convention=equal concentrations,
 constants={}},pourbaix/check=WaterFe,
 pourbaix/calculate=WaterFe,
 pourbaix/pH min=0,pourbaix/pH max=14,
 pourbaix/finalize=WaterFe}
\begin{tikzpicture}
\begin{axis}[pourbaix/axis,width=8cm,height=6cm,
 xmin=0,xmax=14,ymin=-1.2,ymax=1.5]
 \addpourbaix{WaterFe};
 \addpourbaixwater[
  prepare={diagram=WaterFe,oxygen pressure=.21},
  domain labels=false,
  boundary style={blue,dashed,thick}];
\end{axis}
\end{tikzpicture}
\end{codeexample}
The pressure 0.21 bar is an illustrative oxygen partial pressure; neither the
ambient total pressure nor humidity is inferred. Use the separate preparation
key before plotting when the generated equation keys must remain available
outside the plotting command. Multiple preparations can be isolated in TeX groups.

\section{Characteristic extra ticks}
\label{sec:extra-ticks}
A selected diagram can request ordinary PGFPlots \emph{extra ticks}, without
replacing the axis's normal ticks or the user's existing extra ticks. No ticks
are added by default. This is a presentation operation: it changes neither
thermodynamic data nor boundary identifiers, selection, or diagram handles.

\subsection{Selection and eligibility}
\begin{key}{/pgfplots/pourbaix/plot/extra ticks=\meta{all or none} (initially none)}
Set both \texttt{extra x ticks} and \texttt{extra y ticks}. The bare key means
\texttt{all}. Boolean \texttt{true} and \texttt{false} are accepted synonyms for
\texttt{all} and \texttt{none}. For mixed subsets, use the two separate keys.
For example, \texttt{extra ticks=all,extra y ticks=none} enables only x ticks.
\end{key}
\begin{key}{/pgfplots/pourbaix/plot/extra x ticks=\meta{selection} (initially none)}
Accepts \texttt{all}, \texttt{none}, or a braced comma-separated list of final
vertical-boundary identifiers, such as \texttt{\{a,c\}}. The bare key means
\texttt{all}; an empty list also disables the request. For every eligible boundary
add its pH to the native \texttt{extra x ticks}. A boundary must have a nonzero
visible segment after viewport clipping. Unknown identifiers or redox identifiers
in this list produce an error before submitting any plots or ticks.
\end{key}
\begin{key}{/pgfplots/pourbaix/plot/extra y ticks=\meta{selection} (initially none)}
Accepts \texttt{all}, \texttt{none}, or final redox identifiers, such as
\texttt{\{1,3\}}. For $E=a\,\mathrm{pH}+b$, add $b$ in volts only if the
retained boundary has support at pH zero, zero is in the viewport's pH interval,
and $b$ is within its potential interval. There must be a visible boundary
segment. An oblique line retained only from pH 7 onward does \emph{not} create an
intercept tick by extrapolation. A requested but ineligible valid identifier is
simply skipped; an unknown or vertical identifier is an error.
\end{key}
Selections use final chemical identifiers even when \texttt{boundary/1/label}
changes the displayed text. Removed or provisional supports contribute no ticks.
Tick requests are independent of the switches for drawing boundary lines, fills
and domain labels: they query the selected, clipped geometry.

Repeated coordinates are coalesced at a roundoff resolution of 64 machine
epsilons times the larger of 1 and the coordinate magnitudes. This comparison
uses unrounded pH or voltage values; label rounding does not merge distinct ticks.
Existing literal numeric extra ticks are recognized across successive calls,
including water overlays, and are not added twice. Existing PGF expressions,
TeX macros or range notation are preserved but not evaluated for duplicate
recognition. The original list and its ordering remain intact; newly requested
values are appended in increasing order. Normal major/minor ticks are unchanged.

\subsection{Styling and scope}
\begin{key}{/pgfplots/pourbaix/plot/extra x tick style=\marg{options} (initially empty)}
Append ordinary PGFPlots options to the native \texttt{extra x tick style} when
new x ticks are added. Use \texttt{tick label style}, \texttt{xticklabel},
\texttt{grid}, number formatting, and ordinary TikZ node options as appropriate.
This value is an option list, like the other Pourbaix appearance keys.
\end{key}
\begin{key}{/pgfplots/pourbaix/plot/extra y tick style=\marg{options} (initially empty)}
The corresponding native extra y tick options. The two styles are independent.
\end{key}
The example below defines the axis style
\texttt{pourbaix/characteristic ticks} to combine selection and appearance.
You can also configure native \texttt{extra x tick style} and
\texttt{extra y tick style} directly in the axis options.

Styles and tick additions survive the plotting command's temporary option group
but remain local to its caller's TeX group, normally the axis. They use PGFPlots'
normal deferred tick rendering. Native extra tick styles apply to the whole
extra-tick set of an axis, including pre-existing ticks; they are not independently
attached to each contributing plot. Later settings take normal PGF precedence.
Grouping an entire plotting call in an additional TeX group also scopes its tick
changes. Neither an unrelated axis nor later diagrams inherit those changes.
User-supplied \texttt{extra x tick labels} or \texttt{extra y tick labels} lists
must account for the appended positions; a numeric formatter is usually easier.
Coincidences with ordinary major ticks and label overlaps are not automatically
resolved: use native tick lists and styles when a different layout is desired.

Only fixed key paths and validated numbers are generated by Lua. Text and styles
remain inert TeX tokens and are captured before the local option group closes.
The adapter reads the macros used by PGFPlots to store existing extra ticks;
updates go through its public keys, never through private plot lists.

\subsection{Water intercepts}
\begin{key}{/pgfplots/pourbaix/water/extra y ticks=\meta{selection} (initially none)}
Select \texttt{all}, \texttt{none}, \texttt{H2}, \texttt{O2}, or
\texttt{\{H2,O2\}}. The bare key enables both. Add the pressure-dependent
intercepts only when pH zero and the corresponding potentials lie in the viewport.
No vertical ticks arise from the water model.
\end{key}
\begin{key}{/pgfplots/pourbaix/water/extra ticks=\meta{all or none} (initially none)}
Shorthand for the water \texttt{extra y ticks} key; bare means \texttt{all}.
\end{key}
\begin{key}{/pgfplots/pourbaix/water/extra y tick style=\marg{options} (initially empty)}
Append native extra y tick styling when this water call contributes new values.
The same axis-wide styling and scoping rules apply as for an element diagram.
\end{key}

\subsection{An executed selection and style example}
Nickel is useful here: its metal/aqueous boundary reaches pH zero, whereas the
metal/hydroxide boundary begins at the precipitation pH. Only the former creates
an intercept tick. The user's extra x tick at pH 3 is retained. The example also
selects the oxygen intercept from the water overlay, but not the hydrogen one.
\begin{codeexample}[width=11cm]
\pgfplotsset{
 pourbaix/characteristic ticks/.style={
  pourbaix/plot/extra x ticks=all,
  pourbaix/plot/extra y ticks={1},
  pourbaix/plot/extra x tick style={
   grid=major,tick label style={rotate=45,anchor=east},
   xticklabel={\pgfmathprintnumber[fixed,precision=2]{\tick}}},
  pourbaix/plot/extra y tick style={
   tick label style={text=blue},
   yticklabel={\pgfmathprintnumber[fixed,precision=3]{\tick}}}},
 pourbaix/declare={name=TickNi,element=Ni,
  species={Ni(s),Ni^2+,Ni(OH)2(s)},
  temperature=298.15,concentration=.01,
  boundary convention=equal concentrations,constants={}},
 pourbaix/check=TickNi,pourbaix/calculate=TickNi,
 pourbaix/pH min=0,pourbaix/pH max=14,pourbaix/finalize=TickNi}
\begin{tikzpicture}
\begin{axis}[pourbaix/axis,width=11cm,height=7cm,
 xmin=0,xmax=14,ymin=-1.2,ymax=1.5,
 extra x ticks={3},pourbaix/characteristic ticks]
 \addpourbaix{TickNi};
 \addpourbaixwater[prepare={diagram=TickNi},
  domain labels=false,extra y ticks={O2}];
\end{axis}
\end{tikzpicture}
\end{codeexample}
The rounded labels are presentation only. The precipitation pH is 7.397 for
this dataset; the metal/aqueous intercept is approximately $-0.31615935\,\mathrm V$.
For explicit vertical selection use \texttt{extra x ticks=\{a\}} in the
\cs{addpourbaix} options. To suppress both requests supplied by the style, pass
\texttt{extra ticks=none} to that call. The water request is independent.

\section{Classification pics outside PGFPlots}
The following examples use the five-species iron declaration, including its
explicit temperature. Both are ordinary TikZ pictures with no \texttt{axis}.
The oxidation view uses equally spaced group positions: spacing is schematic,
not proportional to differences in oxidation number. Labels show signed values
(nonintegers use six significant digits); stored values retain their precision.

% Repeat the earlier inline declaration in the surrounding manual scope:
% codeexample executions have their own local scope.
\pgfplotsset{pourbaix/declare={
  name=Fe1E-2,element=Fe,temperature=298.15,
  species={Fe(s),Fe^2+,Fe^3+,Fe(OH)2(s),Fe(OH)3(s)},
  concentration=1e-2,boundary convention=equal concentrations,
  constants={
    standard potential={oxidized=Fe^3+,reduced=Fe^2+,value=0.77},
    standard potential={oxidized=Fe^2+,reduced=Fe(s),value=-0.44},
    solubility product={solid=Fe(OH)2(s),value=1e-15},
    solubility product={solid=Fe(OH)3(s),value=1e-38}
  }
}}
\begin{codeexample}[width=0pt]
\begin{tikzpicture}
  \pic (states) {pourbaix oxidation=Fe1E-2};
\end{tikzpicture}
\end{codeexample}
The acidity view gives each oxidation group its own row, ordered from more acidic
to less acidic. Horizontal spacing does not encode an equilibrium constant.
\begin{codeexample}[width=0pt]
\begin{tikzpicture}[
  pourbaix/species/.append style={draw=blue!60!black}]
  \pic (acid) {pourbaix acidity=Fe1E-2};
\end{tikzpicture}
\end{codeexample}
Plain and ConTeXt documents use the same \verb|\pic| commands inside their normal
TikZ picture delimiters. A pic name provides TikZ's usual prefix for its nodes.
For example, \texttt{states-species-1} is the first species in the original
declaration list, and \texttt{states-state-1} is the lowest oxidation group.
The suffixes \texttt{-species-\meta{index}}, \texttt{-state-\meta{group}} and
\texttt{-caption} are available for anchoring annotations. Use distinct pic names
when several schematics share a picture.

\subsection{TikZ pic reference}
\begin{key}{/tikz/pics/pourbaix oxidation=\meta{diagram name}}
Pic type for the increasing-oxidation classification. Inside a TikZ picture, use:
\begin{verbatim}
\pic (name) {pourbaix oxidation=Fe1E-2};
\end{verbatim}
The name
argument is read literally; it must be visible in the current scope. It can be
omitted when an explicit handle is supplied with the following key.
\end{key}
\begin{key}{/tikz/pics/pourbaix acidity=\meta{diagram name}}
Pic type for the per-oxidation-group acidity order. The same scope, naming,
handle and styling rules apply. It does not create a PGFPlots axis.
\end{key}
\begin{key}{/tikz/pourbaix/handle=\meta{handle}}
Default empty: resolve the diagram name from the pic argument. A nonempty handle
takes precedence. Capture one using \verb|\pgfplotsgetpourbaix|.
For the last declared diagram, use \verb|\let\saved\pgfplotslastpourbaix|.
Then pass \verb|pourbaix/handle/.expanded=\saved| in the pic options. This also supports
anonymous diagrams and snapshots whose original name was later redefined.
\end{key}
\begin{key}{/tikz/pourbaix/column sep=\meta{dimension}}
Default \texttt{2.5cm}. Horizontal unit for group positions or acidity positions.
This is diagram layout, not a physical scale.
\end{key}
\begin{key}{/tikz/pourbaix/row sep=\meta{dimension}}
Default \texttt{0.9cm}. Vertical layout unit. Acidity group rows are separated
by 1.6 units to leave room for their arrows.
\end{key}
\begin{key}{/tikz/pourbaix/species}
Style for species nodes. Default: drawn outline, 2pt rounded corners, 4pt inner
separation. Customize with normal TikZ \texttt{.style} or \texttt{.append style}.
\end{key}
\begin{key}{/tikz/pourbaix/state}
Style for oxidation labels and tie markers. Default: 2pt inner separation.
\end{key}
\begin{key}{/tikz/pourbaix/oxidation number format=\meta{arabic or roman}}
Default \texttt{arabic}. Formats oxidation-number labels in the oxidation pic,
acidity pic and combined map. Roman numerals are upright and retain the sign:
for example, $-\mathrm{I}$ and $+\mathrm{IV}$. Zero remains $0$.
Nonintegers and values with absolute magnitude above 3999 retain the ordinary
numeric notation; they are never rounded to an integer for Roman conversion.
Set this required choice on an individual pic or on its enclosing TikZ picture.
It follows TeX grouping and leaves species charges, boundary identifiers,
classifications and exported numeric data unchanged. Unknown choices are errors.
Aya's classification example demonstrates the option on all three pics.
\end{key}
\begin{key}{/tikz/pourbaix/axis}
Style for schematic arrows and ticks. Default: \texttt{semithick}. This TikZ key
is distinct from the PGFPlots axis-label style in the plot reference.
\end{key}
\begin{key}{/tikz/pourbaix/caption}
Style for caption nodes. Default: 3pt inner separation.
\end{key}
\begin{key}{/tikz/pourbaix/oxidation caption=\meta{text}}
Default \texttt{increasing oxidation number}. This is ordinary TeX node text,
stored by TikZ and never passed through the Lua chemistry parser.
\end{key}
\begin{key}{/tikz/pourbaix/acidity direction=\meta{increasing or decreasing}}
Default \texttt{decreasing}. Controls the left-to-right species order within
an oxidation group in the acidity pic and combined map. It also sets the
corresponding acidity caption; set a custom caption afterwards if desired.
This required choice is TeX-scoped and has no effect on the oxidation-only pic,
scientific snapshots, species indices or boundary identifiers. Equal-acidity
markers remain ties in either direction. The vertical oxidation order of the
combined map remains increasing upwards.
\end{key}
\begin{key}{/tikz/pourbaix/acidity caption=\meta{text}}
Default \texttt{decreasing acidity}. The text and all styles follow TikZ grouping.
\end{key}

\section{A standalone map of species and final boundaries}
The map combines both classifications and the final pair graph. Oxidation numbers
increase from bottom to top; within each row, formal acidity decreases from left
to right. Positions are schematic, not pH or potential coordinates. Every declared
species appears. A species absent from all selected cells is dimmed and dashed;
retained pairs are linked with the same definitive identifiers as the E-pH plot.
A link may represent several separated retained pieces; consult the segment
readouts for their intervals.

\begin{key}{/tikz/pics/pourbaix map=\meta{name}}
Use \texttt{\string\pic \{pourbaix map=\meta{name}\};} in any TikZ picture,
including outside PGFPlots. Requires a selected diagram. The existing
\texttt{/tikz/pourbaix/handle} option can override the name with a captured
handle. The pic has ordinary local scope and respects \texttt{pic actions}.
\end{key}

\begin{codeexample}[width=7cm]
\pgfplotsset{pourbaix/declare={
 name=MapFe,element=Fe,temperature=298.15,
 concentration=.01,
 species={Fe(s),Fe^2+,Fe^3+,Fe(OH)2(s),Fe(OH)3(s)},
 boundary convention=equal concentrations,
 constants={}},
 pourbaix/check=MapFe,pourbaix/calculate=MapFe,
 pourbaix/pH min=0,pourbaix/pH max=14,
 pourbaix/finalize=MapFe}
\begin{tikzpicture}
 \pic (iron) {pourbaix map=MapFe};
\end{tikzpicture}
\end{codeexample}

\begin{key}{/tikz/pourbaix/map row sep=\meta{dimension} (initially 1.8cm)}
Vertical spacing of oxidation groups. The existing \texttt{column sep}, initially
\texttt{2.5cm}, controls horizontal spacing within a group. Neither distance
encodes numerical oxidation differences or thermodynamic constants.
\end{key}
The map reuses \texttt{pourbaix/species}, \texttt{pourbaix/state},
\texttt{pourbaix/axis}, \texttt{pourbaix/caption} and both classification-caption
text keys. Its additional TikZ styles are:
\begin{description}
\item[\texttt{pourbaix/map boundary}] Initially \texttt{semithick}.
\item[\texttt{pourbaix/map boundary label}] Initially
  \texttt{midway,fill=white,inner sep=1.5pt}.
\item[\texttt{pourbaix/map removed species}] Initially \texttt{dashed,text=gray}.
\item[\texttt{pourbaix/map long edge}] Initially \texttt{bend right=55}; applied
  to links spanning more than one oxidation row, to avoid an intervening species.
\item[\texttt{pourbaix/map long label}] Initially \texttt{pos=.3}; applied after
  the base label style on those long links.
\item[\texttt{pourbaix/map boundary }\meta{id}] Optional per-edge style, applied
  after the base edge style. No initial definition.
\end{description}
Dense pair graphs may need spacing or curve adjustments; automatic graph layout
and collision removal are not claimed. Species nodes use the familiar pic names
\texttt{(iron-species-1)}, \texttt{(iron-species-2)}, etc.; oxidation-row labels
use \texttt{(iron-state-1)}, etc. Ordinary TikZ paths can reference these anchors.
Prepared-only oxidation and acidity pics remain available independently and do
not require selection.

\section{Automatically composed diagram reports}
\label{sec:reports}
A report assembles existing views of one immutable selected snapshot. It includes
oxidation and acidity classifications, a labeled species map, the final diagram
with diagnostic overlays, thermodynamic inputs and recovered constants, and all
provisional boundary equations. It does not perform hidden checks or calculations.
Prepare, check, calculate and finalize the diagram first.

\begin{verbatim}
\pgfplotspourbaixreport[
  axis options={xmin=0,xmax=14,ymin=-1.5,ymax=1.5},
  plot options={removed boundary style={red,dashed}}
]{Fe1E-2}
\end{verbatim}
The command is used in ordinary document flow, outside an axis or picture. Its
syntax is identical in all three supported formats. It uses ordinary TikZ pictures
and PGFPlots axes; it does not require a LaTeX table package. Each card and picture
is an indivisible block, but page breaks are possible between blocks. Do not wrap
an entire report in a box or an unbreakable example environment. Reports do not
set page dimensions, fonts or global compatibility levels.

\subsection{Options and styles}
\begin{command}{\pgfplotspourbaixreport\oarg{options}\marg{name}}
Resolve \meta{name} once and compose the requested sections from that selected
snapshot. The named alias, last handle and stored data remain unchanged.
An unknown name, invalid handle or unselected diagram raises an error. Ordinary
PGF diagnostics apply to unknown option keys; unknown section names are errors.
An explicit handle takes precedence over the name, which may then be empty.
Options are scoped to the report. The section list may be empty.
\end{command}

Keys below belong to \texttt{/pgfplots/pourbaix/report/}. Values are required;
there are no implicit values for bare keys.
\begin{key}{/pgfplots/pourbaix/report/handle=\meta{handle} (initially empty)}
Use a captured selected snapshot. As elsewhere, use \texttt{/.expanded} when a
macro supplies its decimal handle. This does not rebind its original name.
\end{key}
\begin{key}{/pgfplots/pourbaix/report/sections=\meta{list}}
An ordered comma-separated list of section names. Default:
\texttt{header}, \texttt{oxidation}, \texttt{acidity}, \texttt{map},
\texttt{plot}, \texttt{events}, \texttt{inputs}, \texttt{constants},
\texttt{boundaries} (in that order).
Each section runs in its own \cs{foreach} group. Reordering and omission are
supported. The \texttt{events} section is omitted when there are no elimination
events. Other empty data lists produce only their heading.
\end{key}
\begin{key}{/pgfplots/pourbaix/report/axis options=\meta{key list} (initially empty)}
Ordinary PGFPlots options, applied after \texttt{pourbaix/axis},
\texttt{width=11cm,height=7cm}. Supply explicit finite decimal \texttt{xmin},
\texttt{xmax}, \texttt{ymin}, \texttt{ymax} if the \texttt{plot} section is requested.
pH is dimensionless and potential is in volts. No plotting bounds are inferred.
For macro-supplied limits use the corresponding key's \texttt{/.expanded} handler.
The viewport may crop the selected pH range; it does not change the selection.
\end{key}
\begin{key}{/pgfplots/pourbaix/report/plot options=\meta{key list} (initially empty)}
Options of \cs{addpourbaix}, applied after \texttt{show removed}.
For example, \texttt{show removed=false} suppresses both diagnostic overlays.
The report supplies its captured handle last so every block describes the same
snapshot. Options and styles remain TeX tokens and follow ordinary expansion rules.
\end{key}
\begin{key}{/pgfplots/pourbaix/report/text width=\meta{dimension} (initially 11cm)}
Text width of headings and data cards; TikZ adds inner padding. This does not set
axis dimensions or scale schematics. Choose a width suitable for the host page.
\end{key}
The public TikZ styles are:
\begin{itemize}
\item \texttt{pourbaix/report heading};
\item \texttt{pourbaix/report card};
\item \texttt{pourbaix/report classification}.
\end{itemize}
Headings use dark blue with white text; cards have a pale blue background and
rounded borders. Classifications default to \texttt{scale=.8,transform shape}.
Use \texttt{/.append style} or \texttt{/.style} as with any TikZ style. Existing
species, map, plot and node styles continue to apply. There is no automatic
layout or label-collision solver. Numerical display uses PGF's number printer
with precision~5; all calculations and exported coefficients retain full precision.

\begin{verbatim}
\tikzset{pourbaix/report card/.append style={fill=green!3}}
\pgfplotspourbaixreport[
  handle/.expanded=\saved,
  sections={header,inputs,constants,boundaries},
  text width=10cm
]{}
\end{verbatim}
This numerical-only report needs no axis options. The full report itself uses
\cs{foreach} to traverse sections and data cards. A document can also use
\cs{foreach} to declare and report several independent systems; see Section~\ref{sec:report-examples}.

\subsection{Reading the numerical cards}
Inputs are numbered independently of boundary labels. Each gives a balanced
reaction and its $\log_{10}K$, together with its user, default or database origin.
For imported potentials the source $E^\circ$ is also printed. Source identifiers
refer to the pinned catalogue references in the thermodynamic-library chapter;
standard-state conventions are displayed, not silently harmonized. Input log
constants are a common free-energy representation: they are not all solubility
products or acidity constants.

Each provisional pair card shows its species, selection status and equation.
An arrow from a provisional name to a final identifier links that pair to the
plot and species map. Retained pH intervals refer to the full selected window,
not the cropped viewport. A vertical locus has a singleton pH interval; its
potential extent is shown graphically and remains available in segment data.
The card also gives the balanced reaction, $\log_{10}K$, $E^\circ$ when relevant,
and the linear combination of numbered inputs used to derive the log constant.
Never interpret this combination as an unweighted sum of standard potentials.
Identity and empty loci are reported without inventing a finite line.

A separate section lists water data, selected solid phases, recovered hydroxide
solubility products and cumulative formation constants in logarithmic form.
Logarithms avoid overflow and clarify which reaction a constant belongs to.
Elimination cards give the affected species, retired limits, replacement pair,
pH interval and reaction log constant. They distinguish strict disproportionation
from zero-width events. Inactive supports need not all result from dismutation;
the general selection-status definitions still apply.

\clearpage\part{Scientific model and computation}\label{part:science}
\section{Scientific meaning of the input}
\subsection{Element balance and limiting boundaries}
Let $\nu_i$ be the number of atoms of the selected element in species $i$. For an
aqueous species let $c_i=[i]$; for a solid define its generalized concentration as
$c_i=n_i/V$. The diagram concentration is
\[
  c_{\mathrm t}=\sum_i \nu_i c_i.
\]
With $c_i$ expressed in mol/L and $V$ in litres, $c_{\mathrm t}$ counts moles of the
selected element per litre. All five iron species in the first example have
$\nu_i=1$, so their contributions enter with unit coefficients.

The intended construction is a \emph{limiting} diagram: a boundary between two
selected species neglects the contributions of all other species. This is not a
simultaneous full speciation calculation. Aqueous domains describe predominance;
solid domains describe existence.

\subsection{Boundary between an aqueous species and a solid}
The solid is at its existence limit. Its incipient amount has a generalized
concentration negligible compared with the concentration of the aqueous species.
Contributions from species outside the pair are neglected in the limiting
construction. This is a rule about amounts and concentrations, not an assignment
of a zero solid activity. The calculation stage uses unit solid activity and the aqueous element balance.

\subsection{Boundary between two solids}
At a boundary between two solids A and B, both solids are at their existence
limit. The boundary rule is simply
\[
  a_{\mathrm A}=a_{\mathrm B}=1.
\]
These dimensionless activities are the only input required for the two solids.
The algorithm does not determine their amounts or the concentrations of other
species at this boundary. In particular, it does not use the total element
balance to assign solid amounts or translate an existence limit into a zero
activity. This rule is independent of the diagram concentration, of the number
of element atoms in either solid formula, and of the aqueous boundary convention.
It supplies boundary input; locating the boundary remains a later computational
step.

\subsection{Two explicit aqueous conventions}
Consider two aqueous species A and B with element counts $a$ and $b$, so that the
limiting balance is $a c_{\mathrm A}+b c_{\mathrm B}=c_{\mathrm t}$.

With \texttt{boundary convention=equal concentrations}, the rule is
\[
 c_{\mathrm A}=c_{\mathrm B}
   =\frac{c_{\mathrm t}}{a+b}.
\]
With \texttt{boundary convention=equal element concentrations}, the rule is
\[
 a c_{\mathrm A}=b c_{\mathrm B}=\frac{c_{\mathrm t}}{2},\qquad
 c_{\mathrm A}=\frac{c_{\mathrm t}}{2a},\quad
 c_{\mathrm B}=\frac{c_{\mathrm t}}{2b}.
\]
These rules coincide when $a=b$; they differ otherwise. The choice is mandatory
even when all supplied species happen to contain one atom of the element.
Neither rule sets both species concentrations equal to $c_{\mathrm t}$.

For example, the formulas \texttt{Cr2O7\string^2-} and \texttt{Cr\string^3+}
contain respectively two and one chromium atoms. At $c_{\mathrm t}=0.012$ mol/L,
equal species concentrations give $0.004$ and $0.004$ mol/L. Equal element
concentrations give $0.003$ and $0.006$ mol/L. This illustrates concentration
bookkeeping only; it does not assert thermodynamic predominance of either species.

\begin{codeexample}[width=3.2cm]
\pgfplotsset{pourbaix/declare={
  element=Cr,temperature=298.15,
  species={Cr2O7^2-,Cr^3+},
  constants={},
  concentration=1.2e-2,
  boundary convention=equal element concentrations
}}
\let\chromiumhandle\pgfplotslastpourbaix
\texttt{Anonymous input saved.}
\end{codeexample}
An empty constant list is explicit input, not a certificate that a boundary can
be computed. Declaration alone inserts no missing thermodynamic datum; the
explicit checking step described below performs closure and water-default selection.

\section{Preparation before boundary calculations}
Preparation runs automatically after input validation and before a declaration
is committed. A failure leaves an existing diagram unchanged. Its products are
classifications and candidate pairs, not boundary equations or a stability map.

\subsection{Increasing oxidation number}
An elemental species has oxidation number zero; a monatomic ion has oxidation
number equal to its charge. For ordinary compounds containing only the selected
element X, hydrogen and oxygen, the automatic assignment uses
\[
  n_{\mathrm X}=\frac{q-h+2o}{\nu_{\mathrm X}},
\]
where $q$ is the species charge and $h,o$ are its H and O atom counts. This assumes
H(+I) and O($-$II), consistent with the conventional rules summarized in the
\href{https://old.goldbook.iupac.org/html/O/O04365_old.html}{IUPAC Gold Book}.
These are restricted rules, not a general inference of bonding from a formula.

An additional element in the formula, a charged polyatomic elemental species,
a nonintegral automatic result, or a compound with H or O as the selected
element requires an explicit assignment. Hydrogen and oxygen elemental species
and monatomic ions still follow the first two rules. The parser cannot detect
every peroxide, hydride or other exceptional bonding situation from composition
alone: supply an override when the conventional assumptions do not apply.

Use \texttt{species data=\{species=Fe3O4(s),oxidation number=8/3\}} inside a
declaration containing that solid to provide a chosen effective average for the
mixed-valence case. Signed decimals and integer fractions are accepted. An
explicit value takes precedence even when an automatic value would be available;
it is recorded as explicit rather than independently certified by the library.
Species are grouped by their assigned numerical value, in increasing order.
A supplied standard potential must name an oxidized member whose oxidation
number is greater than that of the reduced member.

\subsection{Decreasing stoichiometric acidity}
Within one oxidation group, define the formal index
\[
  \eta_i=\frac{h_i-2o_i}{\nu_i}.
\]
When the other elements have the same composition per atom of the selected
element, the species can be compared using H$^+$ and H$_2$O. Higher $\eta$ comes
first. This places Fe$^{2+}$ before Fe(OH)$_2$(s), and dichromate before chromate.
The rule is stoichiometric and uses no equilibrium constants; it is not an order
of acid strength, $pK_a$, precipitation pH or actual domain stability.

Different normalized spectator or ligand compositions invalidate this comparison.
For such a group, supply a distinct positive \texttt{acidity rank} for every
species in that group; smaller ranks come first. Partial explicit orders are
rejected. The same explicit-order requirement applies to multi-species groups
when H or O is the selected element, since the comparison reagents themselves
then contain the selected element. Explicit ranks may also override an otherwise
automatic order.

Equal automatic indices remain ties. Their display order is alphabetical by
normalized species identifier, solely for reproducibility. The acidity pic marks
a tie with an equals sign; it does not invent a strict acidity relation. For
example, FeO(s) and Fe(OH)$_2$(s) have the same formal index. Ties and ligand
dependence may make a nominal vertical candidate degenerate or independent of
pH; that question belongs to the subsequent reaction and boundary analysis.

\begin{codeexample}[width=0pt]
\pgfplotsset{pourbaix/declare={
  name=FeLigands,element=Fe,temperature=298.15,
  species={Fe^2+,Fe(CN)6^4-},constants={},
  concentration=1e-2,
  boundary convention=equal concentrations,
  species data={species=Fe^2+,acidity rank=1},
  species data={species=Fe(CN)6^4-,
    oxidation number=2,acidity rank=2}
}}
\begin{tikzpicture}
  \pic {pourbaix acidity=FeLigands};
\end{tikzpicture}
\end{codeexample}
The explicit ranks here demonstrate user control; they are not an inferred
chemical ordering for the cyanide system.

\subsection{All provisional boundaries}
Every unordered pair is included: $N$ species give $N(N-1)/2$ candidates,
including non-neighbouring species and pairs that may never bound stable domains.
Different oxidation numbers identify a Nernst candidate, nominally horizontal
or oblique. Equal oxidation numbers identify an equilibrium-quotient candidate,
nominally vertical (acid/base, precipitation or complexation).

Preparation does not choose between horizontal and oblique for a redox pair,
balance reactions or evaluate constants. The separate checking step balances
reactions and derives constants, but calculates no slopes, intercepts or stable
domains. A candidate's category is not proof that a finite boundary exists.
In particular, the pH dependence of complexation must eventually be checked
against the chosen ligand conditions.

The deterministic order is increasing oxidation number, then decreasing acidity
within each group. For each species in that ordered list, enumerate all later
species. Independent counters name vertical candidates
\texttt{a-temp}, \texttt{b-temp}, \ldots, \texttt{z-temp}, \texttt{aa-temp}, \ldots,
and redox candidates \texttt{1-temp}, \texttt{2-temp}, \ldots.
For the five-species iron example there are two vertical candidates and eight
redox candidates. The defaults \texttt{a-temp} and \texttt{b-temp} refer to the
Fe(II) and Fe(III) ion/hydroxide pairs, respectively.

Names are local to a diagram, so distinct diagrams may both have a boundary named
\texttt{a-temp}. Prefixes and the suffix can be customized, and a specific pair
can have an explicit name. Reversing the pair in an override has no effect.
Duplicate final names or repeated overrides of the same pair are errors. Adding
species can renumber default names; use a pair-specific override for a name that
must survive such edits. No inventory depends on the numerical values of the
constants or on temperature corrections.

\section{Checking thermodynamic sufficiency}
\label{sec:thermodynamic-check}
Checking is an explicit step after declaration. This separation allows incomplete
input to be inspected with the classification pics. Every boundary
calculation or potential-pH plotting command must require a checked snapshot.
The current check determines constants for all provisional pairs; it does not
construct boundary equations or decide whether a pair bounds stable domains.

\subsection{A complete check and a derived potential}
\begin{codeexample}[width=0pt]
\pgfplotsset{pourbaix/declare={
  name=FeChecked,element=Fe,temperature=298.15,
  species={Fe(s),Fe^2+,Fe^3+},
  concentration=1e-2,
  boundary convention=equal concentrations,
  constants={
    standard potential={
      oxidized=Fe^3+,reduced=Fe^2+,value=0.77},
    standard potential={
      oxidized=Fe^2+,reduced=Fe(s),value=-0.44}
  }
},pourbaix/check=FeChecked}
\pgfplotsgetpourbaix{FeChecked}\checkediron
\pgfkeys{
  /pgfplots/pourbaix/data/\checkediron/2-temp/.cd,
  standard potential/.get=\ironpotential
}
Derived potential:\space
$E^\circ=\pgfmathprintnumber[precision=5]{\ironpotential}$ V.
\end{codeexample}
Here \texttt{2-temp} denotes Fe(III)/Fe(s), because this example contains only
three species. Adding species can change default labels. The independent identity is
\[
 E^\circ_{\mathrm{Fe(III)/Fe(0)}}
 =\frac{E^\circ_{\mathrm{Fe(III)/Fe(II)}}
       +2E^\circ_{\mathrm{Fe(II)/Fe(0)}}}{3}
 =-0.036\overline{6}\ \mathrm V.
\]
Potentials are weighted by electron counts, not averaged without weights.

\subsection{Reaction normalization and thermodynamic closure}
For each supported pair, the library balances the selected element with the
smallest positive integer coefficients, then balances O with H$_2$O, H with
H$^+$, and charge with electrons. For a redox pair the direction is reduction
(higher to lower assigned oxidation number). For an equal-oxidation pair it is
\texttt{first} to \texttt{second} in the prepared inventory. A supplied
\texttt{equilibrium constant} instead uses the explicitly given input direction;
reversing that direction changes the sign of $\log_{10}K^\circ$.

The implementation uses additive reaction data,
\[
 \Delta_{\mathrm r}G^\circ=-RT\ln K^\circ,
 \qquad
 \log_{10}K^\circ=\frac{nFE^\circ}{RT\ln 10}.
\]
For half-reactions the latter is a bookkeeping representation on the common
input electrode reference. It is not a reference-free chemical equilibrium
constant involving an independently measurable electron activity. All potentials
must use that same reference, and all constants must use compatible standard
states at the declared temperature. See the IUPAC definitions of
\href{https://goldbook.iupac.org/terms/view/S05915}{standard equilibrium constant}
and \href{https://goldbook.iupac.org/terms/view/S05913}{standard electromotive force}.
The conversion evaluates $R/F=k_{\mathrm B}/e$ using the exact SI defining values
$k_{\mathrm B}=1.380649\times10^{-23}$ J/K and
$e=1.602176634\times10^{-19}$ C; see the
\href{https://physics.nist.gov/cuu/Constants/index.html}{NIST constants database}.

Each input supplies a reaction vector and a logarithmic constant. Elimination
constructs their linear span. A candidate is determined precisely when its
reaction vector belongs to that span within numerical precision. Combining
reactions then gives its logarithmic constant. Thus no particular chain of
neighbouring species is privileged, and alternative input paths can expose
inconsistent cycles. Each result retains coefficients referring to the original
input records, so its derivation can be inspected.

Calculations remain logarithmic: a constant too large or too small for a finite
positive floating-point value still has a usable \texttt{log K}. The direct
\texttt{equilibrium constant} key is then absent, never clipped to zero or infinity.
The stoichiometric elimination cleanup threshold is $10^{-12}$ in coefficient
units. Reconstructed reactions are independently checked with a $10^{-9}$
relative tolerance per coefficient (unit absolute scale for zeros).
The configurable cycle tolerance is separate: it applies to the absolute
residual in $\log_{10}K$ after elimination. This is numerical validation, not an
uncertainty propagation or a least-squares adjustment of inconsistent data.

Lua constant records accept the same optional \texttt{reaction} string.
Validated records preserve it alongside the signed coefficients in
\texttt{reaction\_coefficients}.

\subsection{Hydroxides and the water ion product}
Without an explicit \texttt{reaction}, the \texttt{solubility product} shorthand describes hydroxide dissolution
into a free monatomic ion and hydroxide ions. For example,
\[
 \mathrm{Fe(OH)_2(s)}=
 \mathrm{Fe^{2+}}+2\mathrm{OH^-}.
\]
Together with H$_2$O $=$ H$^+$ + OH$^-$, this determines the
acid-form dissolution constant $K^\circ_{\mathrm{acid}}=K_s^\circ/(K_w^\circ)^2$.
The opposite, ion-to-solid direction exported for \texttt{a-temp} in the
five-species iron example has $K^\circ=(K_w^\circ)^2/K_s^\circ$.
Free ions used in this algebra may be auxiliary reaction participants even if
not selected as diagram species. They never create additional domains.

An explicit water constant takes priority. If absent, the checker first tries
to derive it from supplied reactions. If it is still undetermined and the
working temperature is exactly 298.15 K, it uses $K_w^\circ=10^{-14}$, marked
\texttt{default} in the results. This is the usual undergraduate approximation,
not a precision water-property correlation. The
\href{https://www.iapws.org/relguide/Ionization.html}{IAPWS ionization release}
describes dependence on temperature and density; no such correlation or implicit
pressure is implemented here. At other temperatures there is no default.
Missing water data cause an error only if they prevent determination of a pair.

A user may instead give a logarithmic constant directly. The following example
also derives the omitted solubility product, $\log_{10}K_s^\circ=-15$:
\begin{codeexample}[width=0pt]
\pgfplotsset{pourbaix/declare={
  name=Hydroxide,element=Fe,temperature=310,
  species={Fe^2+,Fe(OH)2(s)},
  concentration=1e-2,
  boundary convention=equal concentrations,
  constants={
    equilibrium constant={first=Fe^2+,
      second=Fe(OH)2(s),log value=-12},
    water ion product={log value=-13.5}
  }
},pourbaix/check=Hydroxide}
\pgfplotsgetpourbaix{Hydroxide}\hydroxidehandle
\pgfkeys{
  /pgfplots/pourbaix/data/\hydroxidehandle/species 2/.cd,
  log solubility product/.get=\logsolubility
}
Derived $\log_{10}K_s^\circ=\logsolubility$.
\end{codeexample}
These are illustrative input data, not a temperature-dependent reference table.

\subsection{Missing data and current reaction limits}
The checker lists every unresolved candidate by provisional name and oriented
species pair. With \texttt{thermodynamic database=none}, an empty constant list
for Fe(s), Fe(II) and Fe(III) leaves all three redox candidates unresolved.
The default library can instead supply the missing independent values. A message asking for an independent
standard potential does not mean that every listed potential must be supplied:
two independent values may determine the third. Equivalent equilibrium data
may also complete the reaction span. Failed checks publish no partial results.

The current balancing model requires all elements other than the selected
element, H and O to have matching compositions per selected-element atom.
Pairs with different ligands require additional participants and reaction data;
explicit acidity ranks alone do not provide them. Explicit constant reactions
may contain auxiliary ligands, but do not extend this canonical boundary model
or supply ligand activities. Non-hydroxide dissolution inference and boundary
models targeting H or O remain unsupported and receive a reaction-model diagnostic.
They may still be declared and displayed in classification pics.

Even a determined equal-oxidation pair can have zero net proton coefficient.
For instance, a hydration reaction between two solids need not define a finite
vertical pH boundary. This stage certifies numerical reaction data only. It does
not supply missing ligand activities, choose an activity model, solve speciation,
assign solid amounts, or establish domain stability.

\subsection{Checking keys and result data}
\begin{key}{/pgfplots/pourbaix/check=\meta{name}}
Checks the named diagram in the current TeX scope. On success creates an
independent checked snapshot and updates that name and
\texttt{\string\pgfplotslastpourbaix} locally. Existing captured preparation handles
remain unchanged. On failure neither aliases nor the last handle change.
Checking an already checked snapshot reuses its handle and exports its values
again in the current scope. There is no default name.
\end{key}
\begin{key}{/pgfplots/pourbaix/check handle=\meta{handle}}
Checks a captured or anonymous handle; only the last handle is updated.
No named alias changes. Supply a decimal handle or use
\texttt{check handle/.expanded=\string\myhandle} explicitly.
\end{key}
\begin{key}{/pgfplots/pourbaix/diagram/thermodynamic tolerance=\meta{number} (initially 1e-8)}
Finite positive absolute tolerance for the logarithmic residual of a redundant
input reaction. A larger value may accommodate rounded input data; it does not
fit, average or modify them. The selected independent input basis determines
results for inconsistent data accepted within this tolerance.
\end{key}

After a successful check, results are stored as ordinary PGF value keys.
Their root is
\begin{quote}
\texttt{/pgfplots/pourbaix/data/}\meta{checked handle}\texttt{/}
\end{quote}
Use \texttt{.get}, \texttt{\string\pgfkeysgetvalue} or
\texttt{\string\pgfkeysvalueof} as for other PGF data. Each candidate contributes:
\begin{description}
  \item[\texttt{\meta{boundary}/log K}] Finite base-ten logarithm for the
    normalized oriented reaction (formal bookkeeping for a half-reaction).
  \item[\texttt{\meta{boundary}/standard potential}] Standard reduction potential
    in V, for redox candidates only.
  \item[\texttt{\meta{boundary}/equilibrium constant}] Positive dimensionless
    constant for non-redox candidates, present only when representable.
  \item[\texttt{\meta{boundary}/electrons}] Positive electron count for reduction;
    zero for a non-redox reaction.
  \item[\texttt{\meta{boundary}/status}] \texttt{provided} for a directly supplied
    oriented pair, otherwise \texttt{derived}. Conversion of a solubility product
    to the acid-form pair constant is a derivation.
  \item[\texttt{\meta{boundary}/source N}] Nonzero coefficient of input record N
    in the result's logarithmic constant. User inputs come first, followed by
    an inserted water default when applicable and then imported library rows.
    Absent keys mean zero contribution.
\end{description}
Water results use \texttt{water/log K}, \texttt{water/ion product} (when
representable), and \texttt{water/status}; the latter is
\texttt{provided}, \texttt{derived} or \texttt{default}. These keys are absent when
water is undetermined and unnecessary. Recoverable hydroxide products use
\texttt{species N/log solubility product}, \texttt{species N/solubility product}
(when representable) and \texttt{species N/status}; N is the original declaration
index, not the sorted position. All PGF result values obey TeX grouping.
Changing these value keys does not modify the authoritative Lua snapshot.

In Lua, every checked boundary has a \texttt{thermodynamics} record with
\texttt{log\_k}, \texttt{status}, optional \texttt{standard\_potential} or
\texttt{equilibrium\_constant}, a \texttt{reaction} and \texttt{sources}.
The reaction stores \texttt{first}, \texttt{second}, \texttt{electrons} and a
\texttt{coefficients} map: products have positive coefficients, reactants negative.
Each source has an input \texttt{index} and signed \texttt{coefficient}.
The diagram-level \texttt{thermodynamics} record contains \texttt{checked=true},
\texttt{water}, \texttt{solubilities} keyed by species identifier,
\texttt{inputs}, \texttt{input\_count} and \texttt{tolerance}.

\section{Calculating provisional boundary equations}
\label{sec:equations}
After a successful explicit thermodynamic check, the library can calculate the
locus of every provisional pair. This stage neither selects stable domains nor
clips lines at intersections. A calculated pair may never bound two stable
domains in the final diagram. No drawing is submitted by the calculation itself.

\subsection{Activities and the limiting balance}
The default model is an ideal dilute aqueous solution:
\[
 a_i=\frac{[i]}{c^\circ},\qquad c^\circ=1\,\mathrm{mol\,L^{-1}},
 \qquad a_{\rm solid}=a_{\rm water}=1,
 \qquad \mathrm{pH}=-\log_{10}a_{\mathrm{H}^+}.
\]
For two dissolved species containing $u$ and $v$ atoms of the selected element,
equal concentrations gives $[A]=[B]=c_t/(u+v)$; equal element concentrations gives
$[A]=c_t/(2u)$ and $[B]=c_t/(2v)$. At an aqueous--solid limit the dissolved
species alone accounts for the element balance: $[A]=c_t/u$. At a solid--solid
limit both activities are one; no amount balance or aqueous concentration is
computed. In particular, an incipient solid has unit activity, not zero activity.

The catalogue's original molality and molarity conventions are retained in its
metadata. Treating those numerical constants within this common ideal model is
a teaching approximation, not an activity-coefficient correction or a conversion
between standard states. No ionic-strength correction or temperature extrapolation
is performed. Potentials retain the reference of their input constants.

\begin{key}{/pgfplots/pourbaix/diagram/activity model=\meta{model} (initially ideal)}
Currently the only accepted model is \texttt{ideal}. It may be omitted or set
explicitly in a declaration; unknown values are errors. It fixes the activity
rules above, including the standard concentration and unit water activity.
\end{key}

\subsection{From a balanced reaction to a line}
Use the checked reaction orientation and signed stoichiometric coefficients
$\nu_i$, positive for products. Redox reactions are oriented as reductions and
consume $n>0$ electrons. Write $h=\nu_{\mathrm{H}^+}$ and
\[
 L=\nu_A\log_{10}a_A+\nu_B\log_{10}a_B,\qquad
 \log_{10}Q=L-h\,\mathrm{pH},\qquad
 f(T)=\frac{RT\ln 10}{F}.
\]
The electron is excluded from the chemical reaction quotient. Unit water
activity contributes zero to its logarithm. The reduction form of the
\href{https://goldbook.iupac.org/terms/view/09068}{IUPAC Nernst equation} gives
\[
 E=E^\circ-\frac{f(T)}n\log_{10}Q
   =\underbrace{\frac{f(T)h}{n}}_a\mathrm{pH}
    +\underbrace{\frac{f(T)}n(\log_{10}K-L)}_b.
\]
Here $E^\circ=f(T)\log_{10}K/n$ is already determined by thermodynamic closure.
With an explicit \texttt{nernst factor}, replace $f(T)$ everywhere by that
factor $s$, including during closure. The exact physical relation above then
becomes the selected numerical approximation.
For a nonredox reaction, $\log_{10}Q=\log_{10}K$, so that
\[
 \mathrm{pH}=\frac{L-\log_{10}K}{h}\qquad(h\ne0).
\]
The calculations use logarithms directly, without first forming tiny activities
or enormous equilibrium constants. The Nernst factor uses the exact SI ratio
$R/F=k_{\rm B}/e$ and the declared temperature; it is not rounded to
$0.06\,\mathrm V$. Finite output is validated before committing the result.

For illustrative iron data $E^\circ_{\rm III/II}=0.77\,\mathrm V$,
$E^\circ_{\rm II/0}=-0.44\,\mathrm V$, $K_{s,\rm II}=10^{-15}$,
$K_{s,\rm III}=10^{-38}$, $K_w=10^{-14}$ and $c_t=10^{-2}\,\mathrm{mol/L}$
at $298.15\,\mathrm K$, selected checks are:
\begin{center}
\begin{tabular}{ll}
Pair & Equation ($E$ in volts)\\\hline
Fe(II)/Fe(0) & $E=-0.44-f(T)$\\
Fe(III)/Fe(II) & $E=0.77$\\
Fe(II)/Fe(OH)$_2$(s) & $\mathrm{pH}=7.5$\\
Fe(III)/Fe(OH)$_3$(s) & $\mathrm{pH}=2$\\
Fe(OH)$_3$(s)/Fe(OH)$_2$(s) & $E=-f(T)\,\mathrm{pH}+0.77-9f(T)$
\end{tabular}
\end{center}
These values follow the stated illustrative inputs, not the fallback catalogue.
The last equation does not depend on $c_t$ or on the aqueous convention.

\subsection{Degenerate candidates}
Equal oxidation numbers identify a nonredox candidate, but do not guarantee a
unique vertical line. If $h=0$, the quotient is independent of pH and potential.
The residual $r=L-\log_{10}K$ then yields an \texttt{identity} if
$|r|$ does not exceed \texttt{thermodynamic tolerance}, and an \texttt{empty}
locus otherwise. The first means numerical equality throughout the plane under
the selected limiting assumptions; the second means no solution under those
assumptions. Neither means that both species are stable everywhere or nowhere.
For example, hydration between FeO(s) and Fe(OH)$_2$(s) has no proton term.
Such candidates retain their names, but receive no invented constant pH or
plottable expression. All candidates are processed, including these cases.

\subsection{Calculation and result keys}
\begin{key}{/pgfplots/pourbaix/calculate=\meta{name}}
Requires a checked diagram. On success creates a calculated snapshot and updates
its local name alias and \cs{pgfplotslastpourbaix}. An unchecked input produces a
clear error requesting \texttt{pourbaix/check}; checking is not implicit.
Repeated calculation of a calculated handle is idempotent. All thermodynamic
result keys are re-exported under the calculated handle. A failed calculation
changes neither aliases nor the last handle. Names are read literally.
\end{key}
\begin{key}{/pgfplots/pourbaix/calculate handle=\meta{handle}}
Same operation on a stable numeric handle; only the last handle is updated,
never a name alias. Use \texttt{.expanded} when the handle is stored in a macro.
Old prepared and checked snapshots remain unchanged. Both commands obey TeX
scoping and produce no typeset material.
\end{key}

All following paths are relative to
\texttt{/pgfplots/pourbaix/data/}\meta{handle}\texttt{/}.
The scalar \texttt{calculated} is \texttt{true} for calculated snapshots and
\texttt{false} for checked-only snapshots. The keys \texttt{activity model} and
\texttt{standard concentration} report \texttt{ideal} and $1$ (in mol/L);
the latter is exported only after calculation.

Under each \meta{boundary}\texttt{/}, the new readout keys are:
\begin{description}
\item[\texttt{equation type}] \texttt{potential}, \texttt{vertical},
  \texttt{identity} or \texttt{empty}. Inspect this before selecting a plotting form.
\item[\texttt{equation}] Plain ASCII equation text, suitable for \cs{texttt};
  it is not a TeX math formula or a PGF parser expression.
\item[\texttt{expression}] A PGF expression \texttt{(a)*x+(b)} for a potential
  line, or a decimal constant for a vertical line. Absent for degenerate cases.
\item[\texttt{a}, \texttt{b}] Potential-line slope (V per pH unit) and intercept
  (V), absent for other equation types.
\item[\texttt{pH}] Constant pH for a vertical line, absent for other types.
\item[\texttt{log Q offset}, \texttt{proton coefficient}] The dimensionless
  $L$ and signed $h$ used above, present for every candidate.
\item[\texttt{residual}] Dimensionless $L-\log_{10}K$, exported only for a
  degenerate candidate. Its tolerance is in decimal-logarithm units.
\end{description}
Read these stored values with PGF's \texttt{.get}, \cs{pgfkeysgetvalue}, or
\cs{pgfkeysvalueof}. They have no default before successful calculation.
Assignments to exported keys are local display copies and do not mutate Lua
snapshots. Decimal output retains 17 significant digits; use PGF number
formatting for presentation. Numerical Lua limits and the plotting engine's
own range limits still apply.

\subsection{An ordinary PGFPlots plot using calculated data}
The example deliberately plots just two complete provisional lines, not a
finished Pourbaix stability diagram. Equations enter through normal
\cs{addplot} commands, preserving PGFPlots' deferred visualization. The vertical
line uses explicitly chosen potential endpoints. No library-imposed axis range
or clipping domain is implied.
\begin{codeexample}[width=7cm]
\pgfplotsset{pourbaix/declare={
  name=Lines, element=Fe,
  temperature=298.15, concentration=1e-2,
  species={Fe(s),Fe^2+,Fe(OH)2(s)},
  boundary convention=equal concentrations,
  activity model=ideal,
  constants={
    standard potential={oxidized=Fe^2+,
      reduced=Fe(s),value=-0.44},
    solubility product={solid=Fe(OH)2(s),
      value=1e-15}
  }
},pourbaix/check=Lines,
  pourbaix/calculate=Lines}
\pgfplotsgetpourbaix{Lines}\lines
\pgfkeys{/pgfplots/pourbaix/data/\lines/.cd,
  1-temp/expression/.get=\potential,
  a-temp/pH/.get=\precipitation}
\begin{tikzpicture}
\begin{axis}[pourbaix/axis,width=7cm,
  height=5cm,xmin=0,xmax=14,
  ymin=-1,ymax=0,domain=0:14]
  \addplot[blue,samples=2] {\potential};
  \addplot[red] coordinates
    {(\precipitation,-1) (\precipitation,0)};
\end{axis}
\end{tikzpicture}
\end{codeexample}
Capture readout values before reusing a mutable diagram name. Calculating a
snapshot does not enqueue any plot and therefore works outside an axis too.
The pure Lua equivalent is \texttt{core.calculate(checked)}; the registry adds
\texttt{calculate(handle)}. Each boundary then has an \texttt{equation} record
with \texttt{kind}, \texttt{text}, optional \texttt{expression}, \texttt{a},
\texttt{b}, \texttt{ph}, and the quotient coefficients described above.
Its \texttt{log\_activities} table preserves the two endpoint logarithms;
\texttt{diagram.calculation} records the model and standard concentration.

\section{Finalizing the limiting diagram}
\label{sec:selection}
This stage reconstructs the sequence of existence or predominance domains at
successive pH values. It operates on calculated equations, preserving the
provisional inventory. Its output comprises retained boundary segments,
elimination events, and the portions of candidates not retained. It does not
enqueue drawings, solve full aqueous speciation, impose a potential window, or
restrict the result to the stability window of water.

\subsection{An empty common domain and disproportionation}
For three successive oxidation states Red, Int and Ox at fixed pH, let
$E_\ell$ be the Int/Red reduction potential and $E_u$ the Ox/Int reduction
potential. Int requires simultaneously
\[
 E\ge E_\ell,\qquad E\le E_u.
\]
If $E_\ell>E_u$, these half-domains have no common interval. The intermediate
is removed from the chain and Red and Ox become adjacent. This is the
\href{https://goldbook.iupac.org/terms/view/D01799}{disproportionation}
criterion used here: the same intermediate is oxidized and reduced. It does not
refer to two physically disconnected stable regions on opposite sides of the
pH--potential plane.

The replacement uses the checked equation of the \emph{product pair} Red/Ox.
It already exists in the all-pair inventory. In particular, do not average the
two old numerical potentials: the product pair can have different limiting
concentrations. Additive reaction free energies determine its thermodynamic
constant, and its own boundary convention determines its Nernst offset.

For example, take $E^\circ_{\mathrm{Cu^{II}/Cu^I}}=0.159\,\mathrm V$,
$E^\circ_{\mathrm{Cu^I/Cu}}=0.520\,\mathrm V$ and
$c_t=10^{-2}\,\mathrm{mol/L}$ at $298.15\,\mathrm K$. The upper bound of Cu(I)
is $0.159\,\mathrm V$ and its lower bound is approximately $0.402\,\mathrm V$.
It has no common interval and undergoes
\[
 2\mathrm{Cu}^+\longrightarrow\mathrm{Cu}^{2+}+\mathrm{Cu}(s).
\]
The replacement has $E^\circ=(0.159+0.520)/2$ and
$E=E^\circ+f(T)\log_{10}(c_t/c^\circ)/2\simeq0.280\,\mathrm V$.
Here the mean of the retired potentials happens to coincide with the replacement.
In general, use electron-weighted thermodynamic composition and evaluate the
product-pair quotient afresh; numerical line averaging is not the algorithm.

\subsection{Sweep algorithm and its meaning}
\begin{enumerate}
\item Form nodes from the requested pH limits and uniform step. Add every
  calculated vertical pH and every intersection of nonparallel potential lines
  inside the requested range. Classify each resulting open cell at its midpoint.
\item Within each oxidation group, retain the unique species favored by all its
  pairwise nonredox comparisons. For the checked reaction first $\to$ second,
  $r=\log_{10}Q-\log_{10}K<0$ favors the second; $r>0$ favors the first.
  Persistent equalities or incompatible cyclic comparisons are errors, with the
  pH and oxidation group identified. No name-based tie-breaking is used.
\item Traverse these representatives in increasing oxidation number. Whenever
  the last three have incompatible bounds, remove the intermediate, replace its
  adjacent edges with the direct product-pair edge, and repeat. This includes
  cascading eliminations. An earlier replacement can itself be retired later.
\item Record the surviving chain, edges and events in each cell. Merge touching
  pieces of the same source boundary. Merge identical event records across
  touching cells as well. A gap always starts a new record. Coincident lines of
  distinct species pairs are never identified as the same boundary.
\item At each vertical candidate, determine the species chains on either side.
  Retain the common potential interval of the two endpoint domains if it has
  positive extent. Otherwise that vertical candidate is inactive.
\end{enumerate}
The resulting chain gives at most one potential interval per species at a fixed
sample pH. Aqueous members have predominance domains; solid members have existence
domains. Lowest and highest members can have unbounded potential intervals.

This is a sequential construction in the declared \emph{limiting model}, not a
Gibbs-energy minimization with a single simultaneous activity vector for every
species. Different pairs use different limiting element balances. The algorithm
selects one acid/base representative per oxidation state before performing the
redox elimination. It does not infer additional species or claim a full
multiphase equilibrium solution. These choices matter near multiple boundaries;
provisional lines alone should not be interpreted as an exact speciation model.

A redox intermediate with numerically equal bounds has zero domain width.
It is removed with reason \texttt{zero width}, separately from a strictly
favorable \texttt{disproportionation}. Comparison resolution is
$32\epsilon\max(1,|E_\ell|,|E_u|)$ in volts, with
$\epsilon=2^{-52}$. Removing a zero-width domain is not a statement that its
species cannot occur at the common boundary.

\subsection{Mesh resolution, endpoints and reconstruction}
The default step is $0.1$ pH unit. The final uniform cell may be shorter.
Intersection refinement adds nodes rather than shifting events to the nearest
uniform node. For affine candidates the comparison signs cannot change inside
an event-free open cell in exact arithmetic. Thus, for example, an inversion at
pH $1.037$ is retained even with a step of $2$.

Nodes separated by at most
$32\epsilon\max(1,|\mathrm{pH}_{\min}|,|\mathrm{pH}_{\max}|)$ are treated as
numerically coincident. This resolution is exported; features below it are not
resolved. A step or range too small for this resolution is an error. The sweep
also rejects more than $100\,000$ cells, including after refinement, or more
than $2\,000\,000$ candidate line-pair intersection tests. These safeguards
prevent an accidental step from exhausting the TeX process. Increase the step
for excessive uniform sampling; reducing the species inventory may be necessary
for excessive intersection counts.

Reported pH intervals describe closures of classified open cells. Their shared
endpoints can host multiple coexisting species; cell chains are not assertions
of uniqueness at triple points. Zero-length potential segments are not generated.
The same equation can have several retained intervals separated by inactive
support; these are intentionally not merged.

A vertical segment uses its calculated pH exactly and the intersection of
one-sided potential extents. Probes are taken within the nearest event-free
neighborhoods. At a requested pH endpoint, one probe may lie just outside the
window solely to determine the boundary's one-sided domains; no segment outside
the window is exported. Potential endpoints can be infinite. A ray must be
clipped to the user's chosen axis bounds before supplying finite coordinates
for a plot.

\subsection{Configuration and commands}
\begin{key}{/pgfplots/pourbaix/pH min=\meta{number}}
Required finite lower pH limit for selection, with no initial value. This is a
scoped operation setting, independent of \texttt{xmin} and of the declaration.
Negative pH is allowed. Values are literal decimals; use PGF's expansion handlers
explicitly when supplying a macro.
\end{key}
\begin{key}{/pgfplots/pourbaix/pH max=\meta{number}}
Required finite upper limit, strictly greater than \texttt{pH min}. It does not
set an axis limit and has no initial value.
\end{key}
\begin{key}{/pgfplots/pourbaix/pH step=\meta{number} (default 0.1)}
Positive finite sampling step, initially $0.1$. Supplying the key without a value
restores $0.1$. Additional known event nodes are always inserted. All three
settings must precede \texttt{finalize} in a PGF key list. They do not alter a
previously selected snapshot.
\end{key}
The finalization action constructs the definitive domains and boundaries; it does
not choose a current diagram. Result paths are named \texttt{selection/} and
\texttt{selected}.

\begin{key}{/pgfplots/pourbaix/finalize=\meta{name}}
Requires calculated equations. Creates a selected snapshot, updates the local
name alias and \cs{pgfplotslastpourbaix}, and exports its results. No preceding
stage is run implicitly. Repeating selection with the same handle and settings
reuses the snapshot; changed settings create a new one. Failure preserves names,
the last handle and all older results. No default name or plot is implied.
\end{key}
\begin{key}{/pgfplots/pourbaix/finalize handle=\meta{handle}}
Same operation on a numeric handle, without changing any name alias. Use
\texttt{.expanded} with a captured handle macro. This is suitable for anonymous
or independently saved calculated snapshots.
\end{key}

\subsection{Result key reference}
Readouts retain their root path:\par\smallskip
\noindent\texttt{/pgfplots/pourbaix/data/}\meta{handle}\texttt{/}.\par\smallskip
The scalar \texttt{selected} is \texttt{true} after selection and \texttt{false}
for checked or calculated snapshots without selection. All paths below are
relative to its \texttt{selection/} subdirectory. They are stored PGF values,
readable with \texttt{.get}; editing them cannot change the Lua snapshot.

The top-level fields are \texttt{pH min}, \texttt{pH max}, \texttt{pH step},
\texttt{pH resolution}, \texttt{cell count}, \texttt{segment count},
\texttt{event count}, \texttt{candidate count} and \texttt{removed count}.
Indices start at one. Counts delimit the respective arrays; optional keys that
are inapplicable are absent, not zero.

\paragraph{\texttt{segment N/}.}
A reconstructed retained segment has a definitive \texttt{name} (a number or
lowercase letter sequence),
\texttt{source} (the provisional boundary name), \texttt{first species} and
\texttt{second species} (original declaration indices), \texttt{equation type},
\texttt{equation}, \texttt{expression}, \texttt{pH min}, and \texttt{pH max}.
Potential lines additionally have \texttt{a} and \texttt{b}; vertical segments
have \texttt{pH}, \texttt{potential min}, and \texttt{potential max} (volts).
Infinite bounds are the literal strings \texttt{-infinity} and \texttt{infinity},
not PGF numeric expressions. The source identifies the complete checked constants
and their provenance under the existing provisional-boundary path. Thus every
replacement retains access to its standard potential or equilibrium constant.

Segment numbers are deterministic within a snapshot, not permanent identifiers
across different declarations or ranges. Potential segments appear in discovery
order, followed by retained vertical segments in provisional-inventory order.
Use \texttt{source} to identify the chemical pair. Several segments may have the
same source when an inactive gap separates them.

\paragraph{\texttt{event N/}.}
An elimination record has \texttt{species}, \texttt{reduced species},
\texttt{oxidized species} (declaration indices), \texttt{lower source},
\texttt{upper source}, \texttt{replacement source}, \texttt{reason},
\texttt{pH min}, and \texttt{pH max}. Reasons are \texttt{disproportionation}
and \texttt{zero width}. An event's replacement may be transient; consult retained
segments to determine its final support. Zero events is a valid result.

The event also reports its balanced electron-free reaction: nonzero
\texttt{species I coefficient} values, \texttt{proton coefficient},
\texttt{water coefficient}, \texttt{log K}, \texttt{lower multiplier} and
\texttt{upper multiplier}. Coefficients are positive for products and negative
for reactants; absent species coefficients are zero. The primitive reaction is
formed by a positive multiple of the lower reduction and a negative multiple of
the upper reduction to cancel electrons, then divided by the coefficient gcd.
The same multipliers applied to their checked log constants give its log constant.
This is a thermodynamic reaction constant, not a calculated reaction extent or
an assignment of simultaneous concentrations to all three species.

\paragraph{\texttt{candidate N/}.}
One record per original candidate, in original order: \texttt{source},
\texttt{status}, and \texttt{segment count}. Status is \texttt{retained}
(whole candidate support in the pH window), \texttt{partial}, \texttt{inactive},
\texttt{outside range} (vertical pH outside the window), or \texttt{degenerate}
(no unique line). A clipped vertical ray or segment is \texttt{partial} because
its complete provisional line extended over all potentials.

\paragraph{\texttt{removed N/}.}
The complement of retained candidate support within the requested window:
\texttt{source}, \texttt{equation type}, \texttt{pH min}, \texttt{pH max},
and, for vertical pieces, \texttt{potential min} and \texttt{potential max}.
Endpoint closures can meet retained segments; this does not delete a coexistence
point. Degenerate equations have no line to remove. These records inventory all
inactive pieces, including nonadjacent pairs and acid/base exclusion. They must
not all be called disproportionation events: the separate event array carries
that causal information.

\subsection{Copper: one reconstructed replacement}
This inline example executes the complete sequence. The single retained line
replaces the two incompatible Cu(I) boundaries and spans the whole pH range.
\begin{codeexample}[width=7cm]
\pgfplotsset{pourbaix/declare={
 name=CuSweep,element=Cu,temperature=298.15,
 concentration=1e-2,species={Cu(s),Cu^+,Cu^2+},
 boundary convention=equal concentrations,
 thermodynamic database=none,
 constants={
  standard potential={oxidized=Cu^2+,
   reduced=Cu^+,value=.159},
  standard potential={oxidized=Cu^+,
   reduced=Cu(s),value=.520}}},
 pourbaix/check=CuSweep,
 pourbaix/calculate=CuSweep,
 pourbaix/pH min=0,pourbaix/pH max=14,
 pourbaix/finalize=CuSweep}
\pgfplotsgetpourbaix{CuSweep}\cu
\pgfkeys{/pgfplots/pourbaix/data/\cu/selection/.cd,
 segment 1/expression/.get=\cuLine,
 segment 1/pH min/.get=\cuLeft,
 segment 1/pH max/.get=\cuRight}
\begin{tikzpicture}
\begin{axis}[pourbaix/axis,width=7cm,height=5cm,
 xmin=0,xmax=14,ymin=0,ymax=.6]
 \addplot[blue,thick,samples=2,
  domain=\cuLeft:\cuRight] {\cuLine};
 \node at (axis cs:7,.45) {$\mathrm{Cu}^{2+}$};
 \node at (axis cs:7,.1) {$\mathrm{Cu}(s)$};
\end{axis}
\end{tikzpicture}
\end{codeexample}

\clearpage\part{Worked examples}\label{part:examples}
\section{Nine classical species inventories}
\label{sec:gallery}
These examples use explicit species lists at 298.15 K and
$c_t=10^{-2}\,\mathrm{mol/L}$, explicitly choosing equal \emph{element}
concentrations at aqueous--aqueous boundaries. This convention matters for
chromate/dichromate and molecular chlorine/chloride. All calculations and plots
execute inline. Each view shows retained domains and boundaries with the two
water equilibria superposed as blue dashed lines; the metal diagrams are not
clipped to the water stability window.

The zincate and plumbite labels use traditional formulae. In these examples, the explicitly adopted notation identities are
\[
 \mathrm{ZnO}_2^{2-}+2\mathrm{H}_2\mathrm O
   \equiv\mathrm{Zn(OH)}_4^{2-},\qquad
 \mathrm{HPbO}_2^-+\mathrm{H}_2\mathrm O
   \equiv\mathrm{Pb(OH)}_3^-.
\]
With unit water activity their numerical aqueous activities are identified.
These are teaching notation conventions, not measured hydration equilibria or
additional species domains. Catalogue records identify them as such. Do not
select both notations as independent chemical species in the same diagram.

The zinc and lead windows extend respectively to pH 16 and 20 to display the
amphoteric branches with the retained source constants. In particular, the lead
extension is a formal extrapolation of the ideal activity model to illustrate
its algebra: it is not a prediction for a physically dilute solution at pH 20.
The selected sets exclude many real species (including higher lead oxides,
chromium hydroxo complexes, manganate and chlorate); their omission is deliberate.
A region is stable only within the stated limiting construction and species set.
A missing region need not be a plotting error: selection can eliminate it.

The catalogue records and their provenance are documented in Part~\ref{part:data}.
Source standard states are not fully harmonized; the ideal concentration model
is an approximation when using data tabulated on another activity scale.

\subsection{Copper}
\begin{codeexample}[width=11cm]
\pgfplotsset{pourbaix/declare={
 name=GalleryCu,element=Cu,
 species={Cu(s),Cu^2+,Cu2O(s),Cu(OH)2(s)},
 temperature=298.15,concentration=.01,
 boundary convention=equal element concentrations,
 constants={}},
 pourbaix/check=GalleryCu,
 pourbaix/calculate=GalleryCu,
 pourbaix/pH min=0,pourbaix/pH max=14,
 pourbaix/finalize=GalleryCu}
\begin{tikzpicture}
\begin{axis}[pourbaix/axis,width=11cm,height=7cm,
 xmin=0,xmax=14,ymin=-0.8,ymax=1.5]
 \addpourbaix[final diagram]{GalleryCu};
 \addpourbaixwater[
  prepare={diagram=GalleryCu},domain labels=false];
\end{axis}
\end{tikzpicture}
\end{codeexample}
\subsection{Zinc}
\begin{codeexample}[width=11cm]
\pgfplotsset{pourbaix/declare={
 name=GalleryZn,element=Zn,
 species={Zn(s),Zn^2+,Zn(OH)2(s),ZnO2^2-},
 temperature=298.15,concentration=.01,
 boundary convention=equal element concentrations,
 constants={}},
 pourbaix/check=GalleryZn,
 pourbaix/calculate=GalleryZn,
 pourbaix/pH min=0,pourbaix/pH max=16,
 pourbaix/finalize=GalleryZn}
\begin{tikzpicture}
\begin{axis}[pourbaix/axis,width=11cm,height=7cm,
 xmin=0,xmax=16,ymin=-2,ymax=0.7]
 \addpourbaix[final diagram,
  boundary label style={pos=.75},
  domain/ZnO2^2-={label style={rotate=90,font=\small}}
 ]{GalleryZn};
 \addpourbaixwater[
  prepare={diagram=GalleryZn},domain labels=false];
\end{axis}
\end{tikzpicture}
\end{codeexample}
\subsection{Aluminium}
\begin{codeexample}[width=11cm]
\pgfplotsset{pourbaix/declare={
 name=GalleryAl,element=Al,
 species={Al(s),Al^3+,Al(OH)3(s),Al(OH)4^-},
 temperature=298.15,concentration=.01,
 boundary convention=equal element concentrations,
 constants={}},
 pourbaix/check=GalleryAl,
 pourbaix/calculate=GalleryAl,
 pourbaix/pH min=0,pourbaix/pH max=14,
 pourbaix/finalize=GalleryAl}
\begin{tikzpicture}
\begin{axis}[pourbaix/axis,width=11cm,height=7cm,
 xmin=0,xmax=14,ymin=-2.4,ymax=0.8]
 \addpourbaix[final diagram,
  boundary label style={pos=.75},
  domain/Al(OH)4^-={label style={rotate=90,font=\small}}
 ]{GalleryAl};
 \addpourbaixwater[
  prepare={diagram=GalleryAl},domain labels=false];
\end{axis}
\end{tikzpicture}
\end{codeexample}
\subsection{Silver}
\begin{codeexample}[width=11cm]
\pgfplotsset{pourbaix/declare={
 name=GalleryAg,element=Ag,
 species={Ag(s),Ag^+,Ag2O(s)},
 temperature=298.15,concentration=.01,
 boundary convention=equal element concentrations,
 constants={}},
 pourbaix/check=GalleryAg,
 pourbaix/calculate=GalleryAg,
 pourbaix/pH min=0,pourbaix/pH max=14,
 pourbaix/finalize=GalleryAg}
\begin{tikzpicture}
\begin{axis}[pourbaix/axis,width=11cm,height=7cm,
 xmin=0,xmax=14,ymin=-0.5,ymax=1.5]
 \addpourbaix[final diagram]{GalleryAg};
 \addpourbaixwater[
  prepare={diagram=GalleryAg},domain labels=false];
\end{axis}
\end{tikzpicture}
\end{codeexample}
\subsection{Lead}
\begin{codeexample}[width=11cm]
\pgfplotsset{pourbaix/declare={
 name=GalleryPb,element=Pb,
 species={Pb(s),Pb^2+,Pb(OH)2(s),HPbO2^-},
 temperature=298.15,concentration=.01,
 boundary convention=equal element concentrations,
 constants={}},
 pourbaix/check=GalleryPb,
 pourbaix/calculate=GalleryPb,
 pourbaix/pH min=0,pourbaix/pH max=20,
 pourbaix/finalize=GalleryPb}
\begin{tikzpicture}
\begin{axis}[pourbaix/axis,width=11cm,height=7cm,
 xmin=0,xmax=20,ymin=-1.5,ymax=1]
 \addpourbaix[final diagram,
  boundary label style={pos=.75},
  domain/HPbO2^-={label style={rotate=90,font=\small}}
 ]{GalleryPb};
 \addpourbaixwater[
  prepare={diagram=GalleryPb},domain labels=false];
\end{axis}
\end{tikzpicture}
\end{codeexample}
\subsection{Nickel}
\begin{codeexample}[width=11cm]
\pgfplotsset{pourbaix/declare={
 name=GalleryNi,element=Ni,
 species={Ni(s),Ni^2+,Ni(OH)2(s)},
 temperature=298.15,concentration=.01,
 boundary convention=equal element concentrations,
 constants={}},
 pourbaix/check=GalleryNi,
 pourbaix/calculate=GalleryNi,
 pourbaix/pH min=0,pourbaix/pH max=14,
 pourbaix/finalize=GalleryNi}
\begin{tikzpicture}
\begin{axis}[pourbaix/axis,width=11cm,height=7cm,
 xmin=0,xmax=14,ymin=-1.3,ymax=1]
 \addpourbaix[final diagram]{GalleryNi};
 \addpourbaixwater[
  prepare={diagram=GalleryNi},domain labels=false];
\end{axis}
\end{tikzpicture}
\end{codeexample}
\subsection{Chromium}
\begin{codeexample}[width=11cm]
\pgfplotsset{pourbaix/declare={
 name=GalleryCr,element=Cr,
 species={Cr(s),Cr^3+,Cr(OH)3(s),Cr2O7^2-,CrO4^2-},
 temperature=298.15,concentration=.01,
 boundary convention=equal element concentrations,
 constants={}},
 pourbaix/check=GalleryCr,
 pourbaix/calculate=GalleryCr,
 pourbaix/pH min=0,pourbaix/pH max=14,
 pourbaix/finalize=GalleryCr}
\begin{tikzpicture}
\begin{axis}[pourbaix/axis,width=11cm,height=7cm,
 xmin=0,xmax=14,ymin=-1.5,ymax=1.7]
 \addpourbaix[final diagram]{GalleryCr};
 \addpourbaixwater[
  prepare={diagram=GalleryCr},domain labels=false];
\end{axis}
\end{tikzpicture}
\end{codeexample}
\subsection{Manganese}
\begin{codeexample}[width=11cm]
\pgfplotsset{pourbaix/declare={
 name=GalleryMn,element=Mn,
 species={Mn(s),Mn^2+,Mn(OH)2(s),MnO2(s),MnO4^-},
 temperature=298.15,concentration=.01,
 boundary convention=equal element concentrations,
 constants={}},
 pourbaix/check=GalleryMn,
 pourbaix/calculate=GalleryMn,
 pourbaix/pH min=0,pourbaix/pH max=14,
 pourbaix/finalize=GalleryMn}
\begin{tikzpicture}
\begin{axis}[pourbaix/axis,width=11cm,height=7cm,
 xmin=0,xmax=14,ymin=-1.8,ymax=1.8]
 \addpourbaix[final diagram]{GalleryMn};
 \addpourbaixwater[
  prepare={diagram=GalleryMn},domain labels=false];
\end{axis}
\end{tikzpicture}
\end{codeexample}
\subsection{Chlorine}
\begin{codeexample}[width=11cm]
\pgfplotsset{pourbaix/declare={
 name=GalleryCl,element=Cl,
 species={Cl^-,Cl2(aq),HClO(aq),ClO^-},
 temperature=298.15,concentration=.01,
 boundary convention=equal element concentrations,
 constants={}},
 pourbaix/check=GalleryCl,
 pourbaix/calculate=GalleryCl,
 pourbaix/pH min=0,pourbaix/pH max=14,
 pourbaix/finalize=GalleryCl}
\begin{tikzpicture}
\begin{axis}[pourbaix/axis,width=11cm,height=7cm,
 xmin=0,xmax=14,ymin=0,ymax=2]
 \addpourbaix[final diagram,
  boundary/1={label style={anchor=north east}},
  boundary/4={label style={anchor=south west}},
  domain/Cl2(aq)={label={},label style={
   pin={[pin distance=10mm,pin edge={black,thin}]
     60:$\mathrm{Cl}_2(\mathrm{aq})$}}}
 ]{GalleryCl};
 \addpourbaixwater[
  prepare={diagram=GalleryCl},domain labels=false];
\end{axis}
\end{tikzpicture}
\end{codeexample}
\section{Composing reports with foreach}\label{sec:report-examples}
The following complete example executes here, without externalized figures.
All four systems use $298.15\,\mathrm{K}$, $c_t=10^{-2}\,\mathrm{mol/L}$ and equal
\emph{element} concentrations on aqueous--aqueous boundaries. The choice matters
for dichromate, which contains two chromium atoms. Changing a species set changes
the limiting construction; these reports are not full speciation calculations.

Iron uses two hydroxide solids. Copper deliberately retains the three oxidation
states to expose the incompatible Cu(I) band. Zinc includes the amphoteric
hydroxide and tetrahydroxozincate. Chromium uses a reduced set supported by the
current catalogue: metal, Cr(II), Cr(III) and dichromate. It omits chromium
hydroxide, chromate and hydroxo complexes and must not be described as a complete
classical chromium diagram. A five-species chromium set is included in
Section~\ref{sec:gallery}.
The catalogue combines sources with different standard-state conventions;
a versioned source alone does not establish thermodynamic consistency.

\begin{verbatim}
\tikzset{
  pourbaix/report Zn/.style={
    pourbaix/species 4 label/.style={xshift=-8pt}},
  pourbaix/report Cu/.style={
    pourbaix/species 2 label/.style={anchor=east}}}
\foreach \Element/\Species/\pHmax/\Emin/\Emax in {
 Fe/{Fe(s),Fe^2+,Fe^3+,Fe(OH)2(s),Fe(OH)3(s)}/14/-1.2/1.2,
 Cu/{Cu(s),Cu^+,Cu^2+}/14/-.2/.8,
 Zn/{Zn(s),Zn^2+,Zn(OH)2(s),Zn(OH)4^2-}/16/-1.8/.5,
 Cr/{Cr(s),Cr^2+,Cr^3+,Cr2O7^2-}/14/-1.5/1.5}{
  \tikzset{pourbaix/report \Element/.try}
  \pgfplotsset{pourbaix/declare={name=TeachingReport,
    element/.expanded=\Element,species/.expanded=\Species,
    temperature=298.15,concentration=.01,
    boundary convention=equal element concentrations,
    constants={}},
    pourbaix/check=TeachingReport,
    pourbaix/calculate=TeachingReport,
    pourbaix/pH min=0,pourbaix/pH max/.expanded=\pHmax,
    pourbaix/finalize=TeachingReport}
  \pgfplotspourbaixreport[
    axis options={xmin=0,xmax/.expanded=\pHmax,
      ymin/.expanded=\Emin,ymax/.expanded=\Emax},
    plot options={boundary label style={pos=.25},
      final domain label style={font=\small},
      removed domain label style={anchor=south east,font=\small}}
  ]{TeachingReport}
}
\end{verbatim}
The executions below add a page break and a heading before each element; the
scientific data and report options are identical to the displayed source.
The zinc window extends to pH~16 to show the amphoteric branch; the dilute
activity model remains a formal teaching approximation at such high alkalinity.
Element-specific TikZ styles adjust two labels without altering their data;
\texttt{/.try} makes the absent styles for other elements harmless.

\tikzset{
  pourbaix/report Zn/.style={
    pourbaix/species 4 label/.style={xshift=-8pt}},
  pourbaix/report Cu/.style={
    pourbaix/species 2 label/.style={anchor=east}}}
\foreach \Element/\Species/\pHmax/\Emin/\Emax in {
 Fe/{Fe(s),Fe^2+,Fe^3+,Fe(OH)2(s),Fe(OH)3(s)}/14/-1.2/1.2,
 Cu/{Cu(s),Cu^+,Cu^2+}/14/-.2/.8,
 Zn/{Zn(s),Zn^2+,Zn(OH)2(s),Zn(OH)4^2-}/16/-1.8/.5,
 Cr/{Cr(s),Cr^2+,Cr^3+,Cr2O7^2-}/14/-1.5/1.5}{
  \tikzset{pourbaix/report \Element/.try}
  \clearpage
  \subsubsection*{Teaching report: \Element}
  \pgfplotsset{pourbaix/declare={name=TeachingReport,
    element/.expanded=\Element,species/.expanded=\Species,
    temperature=298.15,concentration=.01,
    boundary convention=equal element concentrations,constants={}},
    pourbaix/check=TeachingReport,pourbaix/calculate=TeachingReport,
    pourbaix/pH min=0,pourbaix/pH max/.expanded=\pHmax,pourbaix/finalize=TeachingReport}
  \pgfplotspourbaixreport[
    axis options={xmin=0,xmax/.expanded=\pHmax,
      ymin/.expanded=\Emin,ymax/.expanded=\Emax},
    plot options={boundary label style={pos=.25},
      final domain label style={font=\small},
      removed domain label style={anchor=south east,font=\small}}
  ]{TeachingReport}
}
\clearpage

\clearpage\part{Thermodynamic data and provenance}\label{part:data}
\section{The aqueous thermodynamic library}
\label{sec:database}
When supplied data do not determine every provisional pair, the checker can draw
on an offline library before reporting an error. The catalogue,
\texttt{aqueous-298}, contains 87 entries: 28 standard reduction potentials,
12 acidity constants, 32 cumulative hydroxide formation constants, 7
hydroxide solubility products and 8 further reaction records (including two
explicit notation identities). It is a curated starting collection, not a complete
periodic-table database. Its edition is \texttt{2026-10-08.1}.

\subsection{Fallback, precedence and scope}
Declaration and classification remain independent of the library. During an
explicit check the order is:
\begin{enumerate}
  \item Validate and combine all supplied reaction data, including redundant
    cycles. A contradictory user cycle remains an error.
  \item Determine water data from those inputs, or apply the authorized
    298.15 K approximation.
  \item If a provisional pair remains undetermined, consider library entries
    for the element, temperature and chosen solid phases. Entries with both
    endpoints selected are considered first, then those with one selected
    endpoint, then auxiliary aqueous connections. Within each tier, use the
    stable catalogue order shown below.
  \item Add only reaction vectors outside the current span; a library number
    cannot replace an already determined reaction constant. Stop when the
    supported provisional pairs are determined. If some remain unresolved,
    report their names and the unsuccessful library lookup.
\end{enumerate}

This is a fallback policy, not a fit of user data to a reference database.
A literature value may be skipped because its reaction is already determined,
even if its number differs. For example, the three copper potentials in the
source table are not exactly cycle-consistent; the checker uses independent
information according to the stated ordering. It does not average such entries
or retroactively change a user's potential. To control a particular datum,
supply it explicitly in \texttt{constants}.

For copper, the one-electron values 0.159 and 0.520 V imply a two-electron
potential of 0.3395 V, whereas the separate direct entry is 0.337 V.
The 2.5 mV difference exceeds numerical roundoff. Endpoint priority can select
the direct datum for a two-species diagram and the one-electron data when
Cu$^+$ is included. Supply the same consistent independent constants when
comparing diagrams whose numerical inputs must remain identical.

All entries refer to water as solvent at 298.15 K. The current library does not
extrapolate in temperature, pressure or ionic strength. Gases, amalgams,
nonaqueous solvents and complexes with ligands other than hydroxide are excluded
from this catalogue. Aqueous intermediates may enter the reaction
algebra without becoming diagram species. An unselected solid is never added
as an intermediate. The original species list, concentration balance and
classification pics are unchanged by lookup.

\subsection{Sources and conventions}
The original 24 potentials were selected from the
\href{https://en.wikipedia.org/w/index.php?title=Standard_electrode_potential_(data_page)&oldid=1377006472}{Wikipedia data page, revision 1377006472}.
They are reduction potentials in V versus SHE. If potential fallback is used,
supplied potentials must share that reference; no electrode-reference conversion
is performed. The source describes unit aqueous
activities using a 1 mol/dm$^3$ reference and gives its temperature and pressure
conventions. The elementary-metal entries do not add an amalgam or gas phase.

The original equilibrium data come from the USGS
\href{https://github.com/usgs-coupled/phreeqc3/blob/b0b3be767158ccc3322d2c816625cf470045e67e/database/phreeqc.dat}{PHREEQC database at revision b0b3be7}.
The catalogue records full revision identifiers, retrieval date, source line
ranges, original logarithmic values and conversion rules. The downloaded file's
SHA-256 fingerprint is retained in its source metadata. Only the reference
\texttt{log\_k} fields are used; analytical temperature expressions are not
evaluated. Thus this catalogue does not claim to reproduce every calculation
made by a full PHREEQC installation. See the USGS descriptions of
\href{https://water.usgs.gov/water-resources/software/PHREEQC/documentation/phreeqc3-html/phreeqc3-5.htm}{database files}
and \href{https://water.usgs.gov/water-resources/software/PHREEQC/documentation/phreeqc3-html/phreeqc3-50.htm}{aqueous species data}.

PHREEQC's aqueous activity convention is based on molality. The source
conventions are recorded and have \emph{not} been converted to a common
molarity standard state in this catalogue. Combining these sources is
therefore a documented teaching-data approximation, not a certified precision
thermodynamic dataset. Users requiring a homogeneous standard-state convention
must supply a compatible dataset and disable fallback. No concentration-to-activity
conversion is inferred by checking numerical sufficiency.

The scientific audit is documented in \texttt{doc/scientific-audit.md}.
Independent reconstruction confirms the reactions and logarithmic conversions
of all 59 selected USGS database records against the pinned files. The catalogue
cycle check identifies the copper discrepancy described above, treating distinct
solid phases separately. These checks establish source fidelity and algebraic
consistency of the conversions; they do not establish experimental uncertainties
or primary-measurement provenance for every compiled potential.

When a source supplies $E^\circ$, that potential remains fixed if the user chooses
a rounded Nernst factor. When the source supplies $\log K$, that logarithm remains
fixed and the equivalent potential is $s\log K/n$. This distinction applies to
the aqueous chlorine potential and to the logarithmic manganese and hypochlorite
data, respectively.

Acidity values distinguish $pK_a=-\log_{10}K_a$ from $\log K_a$. Stepwise phosphate
constants are differences of cumulative protonation constants. The carbonate
entry uses dissolved CO$_2$, not an independently defined molecular H$_2$CO$_3$
species. Borate notation follows the source: H$_3$BO$_3$ and H$_2$BO$_3^-$.
Formula spelling is significant in lookup: use the identifiers in the tables;
no general structural-isomer or hydration alias inference is performed.

For a source hydrolysis reaction
\[
 m\mathrm M^{z+}+n\mathrm{H_2O}
 =\mathrm M_m(\mathrm{OH})_n^{(mz-n)+}+n\mathrm H^+,
\]
the library stores the cumulative hydroxide formation constant for
$m\mathrm M^{z+}+n\mathrm{OH^-}=\text{complex}$, using
\[
 \log_{10}\beta_{m,n}
 =\log_{10}K_{\mathrm{hydrolysis}}-n\log_{10}K_{w,\mathrm{source}}.
\]
The source water field is $\log_{10}K_{w,\mathrm{source}}=-14$. For hydroxide
solids, the acid-form dissolution value is converted with
$\log_{10}K_s=\log_{10}K_{\mathrm{acid}}+n\log_{10}K_{w,\mathrm{source}}$.
The transformed constants are fixed catalogue values. Supplying a different
water constant changes subsequent conversions between reaction forms, not the
stored catalogue number. Source reaction coefficients and logarithms remain
available in the catalogue records for these conversions.

\subsection{Declared defaults for solid phases}
The library supplies documented defaults that can be overridden per species.
The default phases are:
\begin{center}
\begin{tabular}{ll}
Formula & Default phase identifier\\\hline
Al(OH)$_3$(s) & \texttt{gibbsite}\\
Fe(OH)$_3$(s) & \texttt{amorphous}\\
Zn(OH)$_2$(s) & \texttt{epsilon}\\
Mn(OH)$_2$(s) & \texttt{pyrochroite}\\
Fe(OH)$_2$(s), Cd(OH)$_2$(s), Pb(OH)$_2$(s) & \texttt{unspecified}
\end{tabular}
\end{center}
The elementary-metal records also use \texttt{unspecified}.
\texttt{unspecified} means the selected source does not identify a crystal form;
it is not a newly defined polymorph. For Al(OH)$_3$, an additional
\texttt{amorphous} entry is available. Unknown or unavailable phase identifiers
never fall back to a different phase. They cause an insufficiency error if the
missing data cannot otherwise be supplied or derived. Phase metadata select a
lookup record; they do not allow simultaneous polymorph domains with identical
species identifiers in the current parser.

\subsection{An iron declaration completed by the library}
\begin{codeexample}[width=0pt]
\pgfplotsset{pourbaix/declare={
  name=LibraryIron,element=Fe,temperature=298.15,
  species={Fe(s),Fe^2+,Fe^3+,Fe(OH)2(s),Fe(OH)3(s)},
  concentration=1e-2,
  boundary convention=equal concentrations,
  thermodynamic database=aqueous-298,
  constants={}
},pourbaix/check=LibraryIron}
\pgfplotsgetpourbaix{LibraryIron}\libraryiron
\pgfkeys{
  /pgfplots/pourbaix/data/\libraryiron/species 5/.cd,
  log solubility product/.get=\librarylogks,
  thermodynamic phase/.get=\libraryphase
}
Phase: \libraryphase.\quad
$\log_{10}K_s^\circ=\pgfmathprintnumber{\librarylogks}$.
\end{codeexample}
The Fe(II) hydroxide solubility product is derived from two source potentials;
it is not an additional tabulated solubility record. Source provenance distinguishes
that derivation from the imported Fe(III) amorphous-hydroxide datum.

\subsection{A phase override and a hydroxide complex}
\begin{codeexample}[width=0pt]
\pgfplotsset{pourbaix/declare={
  name=AmorphousAl,element=Al,temperature=298.15,
  species={Al^3+,Al(OH)3(s)},constants={},
  concentration=1e-2,
  boundary convention=equal concentrations,
  species data={species=Al(OH)3(s),
    thermodynamic phase=amorphous}
},pourbaix/check=AmorphousAl}
\pgfplotsgetpourbaix{AmorphousAl}\amorphousal
\pgfkeys{
  /pgfplots/pourbaix/data/\amorphousal/species 2/.cd,
  log solubility product/.get=\allogks
}
$\log_{10}K_s^\circ=\pgfmathprintnumber{\allogks}$.
\end{codeexample}

\begin{codeexample}[width=0pt]
\pgfplotsset{pourbaix/declare={
  name=HydroxoZn,element=Zn,temperature=298.15,
  species={Zn^2+,Zn(OH)4^2-},constants={},
  concentration=1e-2,
  boundary convention=equal concentrations
},pourbaix/check=HydroxoZn}
\pgfplotsgetpourbaix{HydroxoZn}\hydroxozinc
\pgfkeys{
  /pgfplots/pourbaix/data/\hydroxozinc/species 2/.cd,
  log beta/.get=\zinclogbeta
}
$\log_{10}\beta_4=\pgfmathprintnumber{\zinclogbeta}$.
\end{codeexample}
Here $\beta_4$ refers to Zn$^{2+}+4$OH$^-=$Zn(OH)$_4^{2-}$.
Its logarithm is 14.8, while the corresponding cumulative hydrolysis logarithm
is $-41.2$. They are different constants for different reactions.

\subsection{Library keys and provenance}
\begin{key}{/pgfplots/pourbaix/diagram/thermodynamic database=\meta{selection} (initially aqueous-298)}
The choices are \texttt{aqueous-298} and \texttt{none}. The latter disables
library fallback but retains user-data closure and the authorized water default.
There is no network request at declaration, checking, plotting or typesetting.
The selection belongs to the immutable diagram snapshot.
\end{key}
\begin{key}{/pgfplots/pourbaix/species data/thermodynamic phase=\meta{identifier}}
Optional solid-phase lookup override. Identifiers start with a lowercase letter
and contain lowercase letters, digits and hyphens. Aqueous species cannot have
this field. Omission selects the documented default when one exists. Supplied
numerical constants remain authoritative even if the library lacks that phase.
\end{key}

Additional PGF values under the checked handle's data root are:
\begin{description}
  \item[\texttt{database}, \texttt{database version}, \texttt{database consulted}]
    Selection, edition, and whether missing supported pairs triggered lookup.
  \item[\texttt{input N/origin}] \texttt{user}, \texttt{default} or \texttt{database}.
  \item[\texttt{input N/log K}] The logarithmic value actually used in closure.
  \item[\texttt{input N/reaction}] Original explicit equation, when supplied,
    as detokenized text suitable for printing; aliases retain their input spelling.
  \item[\texttt{input N/coefficient/}\meta{species}] Signed net coefficient
    in the input reaction (negative for reactants). Only nonzero participants
    have keys. Species identifiers are canonical, for example \verb|OH^-|.
  \item[\texttt{input N/database id}, \texttt{input N/source id}]
    Stable catalogue and source identifiers for imported rows.
  \item[\texttt{input N/kind}, \texttt{input N/temperature}]
    Imported constant category and reference temperature in K.
  \item[\texttt{species N/thermodynamic phase}] Selected solid lookup phase,
    whether defaulted or overridden; this alone does not prove a record was used.
  \item[\texttt{species N/log beta}, \texttt{species N/beta}]
    Recoverable cumulative formation constant for a selected hydroxide complex
    present in the catalogue, and its finite positive value when representable.
    \texttt{species N/formation status} records its origin.
  \item[\texttt{\meta{boundary}/pKa}] Available for an aqueous non-redox reaction
    with unit coefficients of both endpoints and one proton released. A
    cumulative multi-proton constant is not labelled as a single-step $pK_a$.
\end{description}
A result's \texttt{status} can also be \texttt{database} when its reaction
is exactly one imported row. Combinations remain \texttt{derived}; their
\texttt{source N} coefficients lead back to input origins and catalogue IDs.
Failed checks publish none of these values. All exported keys obey TeX grouping.

In Lua, \texttt{core.database.list(element)} returns a defensive copy of the
matching entries; omit the element to list all 87. \texttt{get(id)} returns one
entry or reports an unknown identifier. \texttt{sources()} and
\texttt{phase\_defaults()} return copied metadata. Each entry contains its
\texttt{id}, \texttt{kind}, \texttt{element}, \texttt{first}, \texttt{second},
\texttt{coefficients}, numerical value, temperature, solvent, standard-state
label, source ID and locator. Converted records retain source logarithms and
conversion information. The catalogue's \texttt{version} is distinct from the
library implementation version.

A checked snapshot has a \texttt{thermodynamics} record. Its library fields are
\texttt{database\_version}, \texttt{database\_sources},
\texttt{database\_consulted}, selected \texttt{phases}, and hydroxide
\texttt{formations}. An imported input's \texttt{metadata.record}
is a complete copied catalogue entry. Original \texttt{constants} are unchanged.
The catalogue is maintained in \texttt{luapourbaix-database.lua}; adding
entries requires a source, reaction, conditions and regression validation.
No public runtime registration protocol is promised in this version.

\subsection{Catalogue of the catalogue}
The 13 supplementary records and their provenance are documented in
Section~\ref{sec:gallery}.
The following tables are part of this manual's source. Endpoints omit balancing
water, protons, hydroxide and electrons; the complete vectors are in the
catalogue records. All values refer to 298.15 K. The tables identify the actual
formula spelling used by lookup and the quantity being reported.

\subsubsection{Standard reduction potentials (Wikipedia)}
\begin{longtable}{@{}>{\raggedright\arraybackslash}p{42mm}p{87mm}r@{}}
Catalogue ID & Endpoints & $E^\circ$/V\\\hline
\endhead
\texttt{fe-iii-ii} & $\mathrm{Fe}^{3+}$ $\to$ $\mathrm{Fe}^{2+}$ & 0.771\\
\texttt{fe-ii-metal} & $\mathrm{Fe}^{2+}$ $\to$ $\mathrm{Fe}\mathrm{(s)}$ & -0.44\\
\texttt{fe-hydroxide-ii-metal} & $\mathrm{Fe}(\mathrm{O}\mathrm{H})_{2}\mathrm{(s)}$ $\to$ $\mathrm{Fe}\mathrm{(s)}$ & -0.89\\
\texttt{zn-ii-metal} & $\mathrm{Zn}^{2+}$ $\to$ $\mathrm{Zn}\mathrm{(s)}$ & -0.7618\\
\texttt{al-iii-metal} & $\mathrm{Al}^{3+}$ $\to$ $\mathrm{Al}\mathrm{(s)}$ & -1.662\\
\texttt{cu-ii-i} & $\mathrm{Cu}^{2+}$ $\to$ $\mathrm{Cu}^{+}$ & 0.159\\
\texttt{cu-i-metal} & $\mathrm{Cu}^{+}$ $\to$ $\mathrm{Cu}\mathrm{(s)}$ & 0.52\\
\texttt{cu-ii-metal} & $\mathrm{Cu}^{2+}$ $\to$ $\mathrm{Cu}\mathrm{(s)}$ & 0.337\\
\texttt{ag-i-metal} & $\mathrm{Ag}^{+}$ $\to$ $\mathrm{Ag}\mathrm{(s)}$ & 0.7996\\
\texttt{cd-ii-metal} & $\mathrm{Cd}^{2+}$ $\to$ $\mathrm{Cd}\mathrm{(s)}$ & -0.4\\
\texttt{pb-ii-metal} & $\mathrm{Pb}^{2+}$ $\to$ $\mathrm{Pb}\mathrm{(s)}$ & -0.126\\
\texttt{ni-ii-metal} & $\mathrm{Ni}^{2+}$ $\to$ $\mathrm{Ni}\mathrm{(s)}$ & -0.257\\
\texttt{co-ii-metal} & $\mathrm{Co}^{2+}$ $\to$ $\mathrm{Co}\mathrm{(s)}$ & -0.28\\
\texttt{sn-iv-ii} & $\mathrm{Sn}^{4+}$ $\to$ $\mathrm{Sn}^{2+}$ & 0.151\\
\texttt{sn-ii-metal} & $\mathrm{Sn}^{2+}$ $\to$ $\mathrm{Sn}\mathrm{(s)}$ & -0.13\\
\texttt{cr-iii-ii} & $\mathrm{Cr}^{3+}$ $\to$ $\mathrm{Cr}^{2+}$ & -0.407\\
\texttt{cr-ii-metal} & $\mathrm{Cr}^{2+}$ $\to$ $\mathrm{Cr}\mathrm{(s)}$ & -0.9\\
\texttt{cr-dichromate-iii} & $\mathrm{Cr}_{2}\mathrm{O}_{7}^{2-}$ $\to$ $\mathrm{Cr}^{3+}$ & 1.33\\
\texttt{mn-ii-metal} & $\mathrm{Mn}^{2+}$ $\to$ $\mathrm{Mn}\mathrm{(s)}$ & -1.185\\
\texttt{mg-ii-metal} & $\mathrm{Mg}^{2+}$ $\to$ $\mathrm{Mg}\mathrm{(s)}$ & -2.372\\
\texttt{ca-ii-metal} & $\mathrm{Ca}^{2+}$ $\to$ $\mathrm{Ca}\mathrm{(s)}$ & -2.868\\
\texttt{v-iii-ii} & $\mathrm{V}^{3+}$ $\to$ $\mathrm{V}^{2+}$ & -0.26\\
\texttt{v-ii-metal} & $\mathrm{V}^{2+}$ $\to$ $\mathrm{V}\mathrm{(s)}$ & -1.13\\
\texttt{v-v-iv} & $\mathrm{V}\mathrm{O}_{2}^{+}$ $\to$ $\mathrm{V}\mathrm{O}^{2+}$ & 1\\
\end{longtable}

\subsubsection{Acidity constants (USGS)}
\begin{longtable}{@{}>{\raggedright\arraybackslash}p{42mm}p{87mm}r@{}}
Catalogue ID & Endpoints & $pK_a$\\\hline
\endhead
\texttt{phosphate-1} & $\mathrm{H}_{3}\mathrm{P}\mathrm{O}_{4}\mathrm{(aq)}$ $\to$ $\mathrm{H}_{2}\mathrm{P}\mathrm{O}_{4}^{-}$ & 2.168\\
\texttt{phosphate-2} & $\mathrm{H}_{2}\mathrm{P}\mathrm{O}_{4}^{-}$ $\to$ $\mathrm{H}\mathrm{P}\mathrm{O}_{4}^{2-}$ & 7.207\\
\texttt{phosphate-3} & $\mathrm{H}\mathrm{P}\mathrm{O}_{4}^{2-}$ $\to$ $\mathrm{P}\mathrm{O}_{4}^{3-}$ & 12.346\\
\texttt{carbonate-1} & $\mathrm{C}\mathrm{O}_{2}\mathrm{(aq)}$ $\to$ $\mathrm{H}\mathrm{C}\mathrm{O}_{3}^{-}$ & 6.352\\
\texttt{carbonate-2} & $\mathrm{H}\mathrm{C}\mathrm{O}_{3}^{-}$ $\to$ $\mathrm{C}\mathrm{O}_{3}^{2-}$ & 10.329\\
\texttt{ammonium} & $\mathrm{N}\mathrm{H}_{4}^{+}$ $\to$ $\mathrm{N}\mathrm{H}_{3}\mathrm{(aq)}$ & 9.252\\
\texttt{bisulfate} & $\mathrm{H}\mathrm{S}\mathrm{O}_{4}^{-}$ $\to$ $\mathrm{S}\mathrm{O}_{4}^{2-}$ & 1.988\\
\texttt{sulfide-1} & $\mathrm{H}_{2}\mathrm{S}\mathrm{(aq)}$ $\to$ $\mathrm{H}\mathrm{S}^{-}$ & 6.994\\
\texttt{sulfide-2} & $\mathrm{H}\mathrm{S}^{-}$ $\to$ $\mathrm{S}^{2-}$ & 12.918\\
\texttt{hydrogen-fluoride} & $\mathrm{H}\mathrm{F}\mathrm{(aq)}$ $\to$ $\mathrm{F}^{-}$ & 3.18\\
\texttt{borate-source-notation} & $\mathrm{H}_{3}\mathrm{B}\mathrm{O}_{3}\mathrm{(aq)}$ $\to$ $\mathrm{H}_{2}\mathrm{B}\mathrm{O}_{3}^{-}$ & 9.24\\
\end{longtable}

\subsubsection{Cumulative hydroxide formation (USGS)}
\begin{longtable}{@{}>{\raggedright\arraybackslash}p{42mm}p{87mm}r@{}}
Catalogue ID & Endpoints & $\log_{10}\beta$\\\hline
\endhead
\texttt{fe-ii-oh-1} & $\mathrm{Fe}^{2+}$ $\to$ $\mathrm{Fe}\mathrm{O}\mathrm{H}^{+}$ & 4.5\\
\texttt{fe-ii-oh-2} & $\mathrm{Fe}^{2+}$ $\to$ $\mathrm{Fe}(\mathrm{O}\mathrm{H})_{2}\mathrm{(aq)}$ & 7.43\\
\texttt{fe-ii-oh-3} & $\mathrm{Fe}^{2+}$ $\to$ $\mathrm{Fe}(\mathrm{O}\mathrm{H})_{3}^{-}$ & 11\\
\texttt{fe-iii-oh-1} & $\mathrm{Fe}^{3+}$ $\to$ $\mathrm{Fe}\mathrm{O}\mathrm{H}^{2+}$ & 11.81\\
\texttt{fe-iii-oh-2} & $\mathrm{Fe}^{3+}$ $\to$ $\mathrm{Fe}(\mathrm{O}\mathrm{H})_{2}^{+}$ & 22.33\\
\texttt{fe-iii-oh-3} & $\mathrm{Fe}^{3+}$ $\to$ $\mathrm{Fe}(\mathrm{O}\mathrm{H})_{3}\mathrm{(aq)}$ & 29.44\\
\texttt{fe-iii-oh-4} & $\mathrm{Fe}^{3+}$ $\to$ $\mathrm{Fe}(\mathrm{O}\mathrm{H})_{4}^{-}$ & 34.4\\
\texttt{fe-iii-dimer} & $\mathrm{Fe}^{3+}$ $\to$ $\mathrm{Fe}_{2}(\mathrm{O}\mathrm{H})_{2}^{4+}$ & 25.05\\
\texttt{fe-iii-trimer} & $\mathrm{Fe}^{3+}$ $\to$ $\mathrm{Fe}_{3}(\mathrm{O}\mathrm{H})_{4}^{5+}$ & 49.7\\
\texttt{al-oh-1} & $\mathrm{Al}^{3+}$ $\to$ $\mathrm{Al}\mathrm{O}\mathrm{H}^{2+}$ & 9\\
\texttt{al-oh-2} & $\mathrm{Al}^{3+}$ $\to$ $\mathrm{Al}(\mathrm{O}\mathrm{H})_{2}^{+}$ & 17.9\\
\texttt{al-oh-3} & $\mathrm{Al}^{3+}$ $\to$ $\mathrm{Al}(\mathrm{O}\mathrm{H})_{3}\mathrm{(aq)}$ & 25.1\\
\texttt{al-oh-4} & $\mathrm{Al}^{3+}$ $\to$ $\mathrm{Al}(\mathrm{O}\mathrm{H})_{4}^{-}$ & 33.3\\
\texttt{cu-oh-1} & $\mathrm{Cu}^{2+}$ $\to$ $\mathrm{Cu}\mathrm{O}\mathrm{H}^{+}$ & 6\\
\texttt{cu-oh-2} & $\mathrm{Cu}^{2+}$ $\to$ $\mathrm{Cu}(\mathrm{O}\mathrm{H})_{2}\mathrm{(aq)}$ & 14.32\\
\texttt{cu-oh-3} & $\mathrm{Cu}^{2+}$ $\to$ $\mathrm{Cu}(\mathrm{O}\mathrm{H})_{3}^{-}$ & 15.1\\
\texttt{cu-oh-4} & $\mathrm{Cu}^{2+}$ $\to$ $\mathrm{Cu}(\mathrm{O}\mathrm{H})_{4}^{2-}$ & 16.4\\
\texttt{zn-oh-1} & $\mathrm{Zn}^{2+}$ $\to$ $\mathrm{Zn}\mathrm{O}\mathrm{H}^{+}$ & 5.04\\
\texttt{zn-oh-2} & $\mathrm{Zn}^{2+}$ $\to$ $\mathrm{Zn}(\mathrm{O}\mathrm{H})_{2}\mathrm{(aq)}$ & 11.1\\
\texttt{zn-oh-3} & $\mathrm{Zn}^{2+}$ $\to$ $\mathrm{Zn}(\mathrm{O}\mathrm{H})_{3}^{-}$ & 13.6\\
\texttt{zn-oh-4} & $\mathrm{Zn}^{2+}$ $\to$ $\mathrm{Zn}(\mathrm{O}\mathrm{H})_{4}^{2-}$ & 14.8\\
\texttt{cd-oh-1} & $\mathrm{Cd}^{2+}$ $\to$ $\mathrm{Cd}\mathrm{O}\mathrm{H}^{+}$ & 3.92\\
\texttt{cd-oh-2} & $\mathrm{Cd}^{2+}$ $\to$ $\mathrm{Cd}(\mathrm{O}\mathrm{H})_{2}\mathrm{(aq)}$ & 7.65\\
\texttt{cd-oh-3} & $\mathrm{Cd}^{2+}$ $\to$ $\mathrm{Cd}(\mathrm{O}\mathrm{H})_{3}^{-}$ & 8.7\\
\texttt{cd-oh-4} & $\mathrm{Cd}^{2+}$ $\to$ $\mathrm{Cd}(\mathrm{O}\mathrm{H})_{4}^{2-}$ & 8.65\\
\texttt{pb-oh-1} & $\mathrm{Pb}^{2+}$ $\to$ $\mathrm{Pb}\mathrm{O}\mathrm{H}^{+}$ & 6.29\\
\texttt{pb-oh-2} & $\mathrm{Pb}^{2+}$ $\to$ $\mathrm{Pb}(\mathrm{O}\mathrm{H})_{2}\mathrm{(aq)}$ & 10.88\\
\texttt{pb-oh-3} & $\mathrm{Pb}^{2+}$ $\to$ $\mathrm{Pb}(\mathrm{O}\mathrm{H})_{3}^{-}$ & 13.94\\
\texttt{pb-oh-4} & $\mathrm{Pb}^{2+}$ $\to$ $\mathrm{Pb}(\mathrm{O}\mathrm{H})_{4}^{2-}$ & 16.3\\
\texttt{mn-oh-1} & $\mathrm{Mn}^{2+}$ $\to$ $\mathrm{Mn}\mathrm{O}\mathrm{H}^{+}$ & 3.41\\
\texttt{mn-oh-3} & $\mathrm{Mn}^{2+}$ $\to$ $\mathrm{Mn}(\mathrm{O}\mathrm{H})_{3}^{-}$ & 7.2\\
\texttt{mg-oh-1} & $\mathrm{Mg}^{2+}$ $\to$ $\mathrm{Mg}\mathrm{O}\mathrm{H}^{+}$ & 2.56\\
\end{longtable}

\subsubsection{Hydroxide solubility (USGS)}
\begin{longtable}{@{}>{\raggedright\arraybackslash}p{42mm}p{87mm}r@{}}
Catalogue ID & Endpoints & $\log_{10}K_s$\\\hline
\endhead
\texttt{al-gibbsite} & $\mathrm{Al}(\mathrm{O}\mathrm{H})_{3}\mathrm{(s)}$ $\to$ $\mathrm{Al}^{3+}$ & -33.89\\
\texttt{al-amorphous} & $\mathrm{Al}(\mathrm{O}\mathrm{H})_{3}\mathrm{(s)}$ $\to$ $\mathrm{Al}^{3+}$ & -31.2\\
\texttt{fe-iii-amorphous} & $\mathrm{Fe}(\mathrm{O}\mathrm{H})_{3}\mathrm{(s)}$ $\to$ $\mathrm{Fe}^{3+}$ & -37.109\\
\texttt{mn-pyrochroite} & $\mathrm{Mn}(\mathrm{O}\mathrm{H})_{2}\mathrm{(s)}$ $\to$ $\mathrm{Mn}^{2+}$ & -12.8\\
\texttt{zn-epsilon} & $\mathrm{Zn}(\mathrm{O}\mathrm{H})_{2}\mathrm{(s)}$ $\to$ $\mathrm{Zn}^{2+}$ & -16.5\\
\texttt{cd-hydroxide} & $\mathrm{Cd}(\mathrm{O}\mathrm{H})_{2}\mathrm{(s)}$ $\to$ $\mathrm{Cd}^{2+}$ & -14.35\\
\texttt{pb-hydroxide} & $\mathrm{Pb}(\mathrm{O}\mathrm{H})_{2}\mathrm{(s)}$ $\to$ $\mathrm{Pb}^{2+}$ & -19.85\\
\end{longtable}

\section{Supplementary thermodynamic records}
The pinned USGS \href{https://github.com/usgs-coupled/phreeqc3/blob/b0b3be767158ccc3322d2c816625cf470045e67e/database/minteq.v4.dat}{MINTEQ v4 file}
provides the following reference logarithms at 25 degrees Celsius. Acid
\emph{dissolution} constants must not be confused with solubility products.
\begin{center}\small
\begin{tabular}{lll}
Solid & Acid dissolution reaction & $\log_{10}K$\\\hline
Cuprite & $\mathrm{Cu_2O}+2\mathrm H^+\to2\mathrm{Cu}^++\mathrm{H_2O}$ & $-1.406$\\
Copper hydroxide & $\mathrm{Cu(OH)_2}+2\mathrm H^+\to\mathrm{Cu}^{2+}+2\mathrm{H_2O}$ & $8.674$\\
Silver oxide & $\mathrm{Ag_2O}+2\mathrm H^+\to2\mathrm{Ag}^++\mathrm{H_2O}$ & $12.574$\\
Nickel hydroxide & $\mathrm{Ni(OH)_2}+2\mathrm H^+\to\mathrm{Ni}^{2+}+2\mathrm{H_2O}$ & $12.794$\\
Pyrolusite & $\mathrm{MnO_2}+4\mathrm H^++2\mathrm e^-\to\mathrm{Mn}^{2+}+2\mathrm{H_2O}$ & $41.38$
\end{tabular}\end{center}
{\emergencystretch=2em
The source locators are lines 10949--10952 (cuprite),
10953--10956 (copper hydroxide), 10961--10964 (silver oxide),
10965--10968 (nickel hydroxide) and 11005--11008 (pyrolusite).\par}
The three unnamed crystal forms are recorded as
\texttt{unspecified}; cuprite and pyrolusite are explicit phase defaults.
The same file supplies $\log K=127.794$ for permanganate reduction to Mn(II)
(reverse of lines 470--477), and $\log K=-14.56$ for
\[\mathrm{Cr_2O_7}^{2-}+\mathrm{H_2O}\to2\mathrm{CrO_4}^{2-}+2\mathrm H^+.\]
This reverses lines 5000--5007. Electron-weighted conversion is $E^\circ=f\log K/n$.

For amorphous chromium hydroxide, the
\href{https://boris-portal.unibe.ch/server/api/core/bitstreams/f01d2d61-ad0e-4baa-b1ac-ccaa870866df/content}{chromium benchmark, Computational Geosciences 19 (2015), 497--521},
Table 4, gives $\log K=9.35$ at 25 degrees Celsius for
$\mathrm{Cr(OH)_3}+3\mathrm H^+\to\mathrm{Cr}^{3+}+3\mathrm{H_2O}$.
The default phase is \texttt{amorphous}. MINTEQ's Cr(III) hydrolysis entry at
20 degrees Celsius and nonzero ionic strength is not silently substituted.

Chlorine uses $E^\circ(\mathrm{Cl_2(aq)/Cl^-})=1.396\,\mathrm V$ versus SHE
from Table 3 (page 43) of
\href{https://www.collectionscanada.gc.ca/obj/thesescanada/vol2/002/MR60643.PDF}{Chlorine Dioxide as a Potable Water Disinfectant},
which credits Lide (1999). This is a compiled value, not a newly verified primary
measurement; crucially, it refers to \emph{aqueous} chlorine, not chlorine gas.
The pinned USGS \href{https://github.com/usgs-coupled/phreeqc3/blob/b0b3be767158ccc3322d2c816625cf470045e67e/database/llnl.dat}{LLNL file}
provides $pK_a(\mathrm{HClO})=7.5692$ (reverse of lines 4493--4497).
Its water oxidation to \emph{dissolved} oxygen has $\log K=-85.9951$
(lines 712--714); $\mathrm{Cl^-}+\tfrac12\mathrm{O_2(aq)}\to\mathrm{ClO^-}$
has $\log K=-15.1014$ (lines 898--900). Eliminating that dissolved oxygen gives
\[
 \mathrm{ClO^-}+2\mathrm H^++2\mathrm e^-\to\mathrm{Cl^-}+\mathrm{H_2O},
 \qquad \log K=85.9951/2+15.1014.
\]
No gaseous-oxygen standard potential is mixed into this conversion.
Full source revisions, SHA-256 hashes for the pinned database files, reactions
and transformation notes are stored in the offline catalogue.

\clearpage\part{Technical reference}\label{part:technical}
\section{Expansion and the TeX/Lua boundary}
Data keys store their arguments without expansion. The bridge reads their macro
replacement text with \texttt{token.get\_macro}; it never inserts that text into a
Lua expression or calls \texttt{load} on it. Quotes and backslashes therefore remain
input data and are rejected where they violate the relevant grammar.
Use the standard PGF \texttt{.expanded} handler when expansion is intended:
\begin{codeexample}[width=3.2cm]
\def\chosenconcentration{2e-2}
\pgfplotsset{pourbaix/declare={
  element=Fe,temperature=298.15,species={Fe(s),Fe^2+},
  constants={standard potential={
    oxidized=Fe^2+,reduced=Fe(s),value=-0.44}},
  concentration/.expanded=\chosenconcentration,
  boundary convention=equal concentrations
}}
\texttt{Expanded input saved.}
\end{codeexample}
Expansion follows TeX and PGF semantics; it is not a chemical expression parser.
Ordinary PGF styles and key handlers remain executable TeX configuration.
The no-evaluation guarantee concerns chemical data crossing into Lua, not an
attempt to sandbox arbitrary TeX supplied as a style.

Validated ASCII names and decimal handles are returned through local macro
assignments. For classification pics, a dedicated renderer converts validated
formulas to a fixed set of generic TeX and TikZ commands. Explicit reaction strings are validated and exported as detokenized text;
their participants are exported in sorted order. Raw input is never
emitted as executable TeX. Styling remains in TeX keys, outside the Lua data path.
The library preserves the catcodes in force around its loader. No LaTeX-only or
ConTeXt-only command is used by the generic interface.

\section{Deferred plotting and implementation architecture}
PGFPlots separates its survey and visualization phases. Plot commands contribute
data while an axis is read; drawing is completed at axis end, after information
needed for limits and transformations has been collected. The installed
\texttt{pgfplots.code.tex} stores plot records and visualizes them in its axis-end
routine. Its source is maintained in the
\href{https://github.com/pgf-tikz/pgfplots}{PGFPlots repository}.

The library does not patch those internal lists. Classification pics are ordinary
TikZ schematics. The plotting interface captures the selected snapshot when
requested and submits ordinary addplot commands through the survey machinery
for visualization alongside other plots. Looking up a mutable name only
at axis end would incorrectly allow a later declaration to alter an earlier plot.
Deferring all data submission until visualization would also bypass normal limit
collection; that is not the intended integration.

Lua declarations accept \texttt{nernst\_factor="auto"} (or omission) and
positive numeric values. Normalized snapshots store the effective numeric
\texttt{nernst\_factor} and \texttt{nernst\_factor\_source}.

The pure core is \texttt{luapourbaix.lua}. Engine access and the session registry
belong to \texttt{luapourbaix-tex.lua}. Pic rendering lives in
\texttt{luapourbaix-tikz.lua}; E-pH clipping and plot-helper commands live in
\texttt{luapourbaix-plot.lua}. The sole bridge global is \texttt{luapourbaix}.
The generic entry point is:
\begin{quote}
\texttt{pgflibrarypgfplots.pourbaix.code.tex}
\end{quote}

\subsection{Lua data API}
\begin{description}
  \item[\texttt{core.declare(raw)}] Returns a validated diagram table or raises a
    diagnostic. The input has \texttt{name} (optional), \texttt{element},
    \texttt{species} (a string list or Lua array of strings),
    \texttt{temperature} (K), \texttt{concentration}, \texttt{boundary\_convention} and
    \texttt{constants} (an array of records). Constant records use
    \texttt{kind="standard potential"} with \texttt{oxidized}, \texttt{reduced},
    \texttt{value}, or \texttt{kind="solubility product"} with \texttt{solid} and
    \texttt{value} or \texttt{log\_value}. Also supported are
    \texttt{kind="equilibrium constant"} with \texttt{first}, \texttt{second}
    and one value field, and \texttt{kind="water ion product"} with one value
    field. The caller's table is not mutated.
  \item[\texttt{core.parse\_species(text, element)}] Returns the normalized
    \texttt{id}, \texttt{formula}, \texttt{phase}, \texttt{charge}, elemental
    \texttt{composition} and \texttt{element\_count}.
  \item[\texttt{core.check(diagram)}] Returns a checked copy or raises an error;
    the preparation snapshot is never mutated. The pure core loads its pure
    thermodynamic dependency through Lua's package path; the engine bridge
    supplies that dependency using the format's file resolver.
  \item[\texttt{core.new\_registry()}] Returns \texttt{add(raw)},
    \texttt{check(handle)} and \texttt{get(handle)} functions. A successful check
    returns a new handle and checked copy; checking an already checked handle
    returns that handle and a copy. Adding validates before allocating a handle;
    retrieval returns a defensive copy. The additional \texttt{build(raw, options)}
    executes all four stages and allocates one final handle only on complete
    success; options contain \texttt{ph\_min}, \texttt{ph\_max} and optional
    \texttt{ph\_step}. Calls use the dot syntax.
  \item[\texttt{luapourbaix.get(handle)}] Retrieves a defensive copy from the
    engine bridge's registry. The diagram contains \texttt{schema\_version=7},
    \texttt{model="limiting"}, the declared fields, normalized \texttt{species},
    \texttt{species\_by\_id} and constants with explicit units.
    Temperature metadata and the preparation arrays described below are also
    included. An anonymous diagram has no \texttt{name} field.
  \item[\texttt{core.aqueous\_boundary\_concentrations(diagram, A, B)}]
    Returns two concentrations from the declared convention and element balance.
    Requires distinct declared aqueous species. It neither locates a boundary
    nor evaluates a potential or equilibrium constant. Solid arguments are errors.
  \item[\texttt{core.solid\_boundary\_activities(diagram, A, B)}]
    Returns the two dimensionless activities $1,1$ for distinct declared solids
    at their mutual existence limit. It does not compute amounts, concentrations
    of other species or boundary coordinates. Aqueous arguments are errors.
\end{description}
The Lua data schema is versioned independently of the TeX interface. Consumers
must check its schema version before interpreting records.

\subsection{Prepared Lua records}
The raw selector \texttt{thermodynamic\_database} defaults to
\texttt{aqueous-298}. The optional raw Lua fields corresponding to declaration
customizations are:
\begin{description}
  \item[\texttt{species\_data}] Records containing \texttt{species}, with optional
    \texttt{oxidation\_number}, \texttt{acidity\_rank} and
    \texttt{thermodynamic\_phase} fields.
  \item[\texttt{boundary\_labels}] A record with \texttt{vertical\_prefix},
    \texttt{redox\_prefix} and \texttt{suffix}.
  \item[\texttt{boundary\_names}] Records with \texttt{first}, \texttt{second}
    and \texttt{name}.
\end{description}

The validated diagram retains the original species-list order. Each species
additionally has its declaration \texttt{index}, assigned \texttt{oxidation\_number},
\texttt{oxidation\_basis}, \texttt{acidity\_basis}, and either an automatic
\texttt{acidity\_index} or an explicit \texttt{acidity\_rank}.
\texttt{oxidation\_groups} is an increasing array of records containing
\texttt{oxidation\_number}, \texttt{acidity\_basis} and an acidity-ordered
\texttt{species} array. \texttt{ordered\_species} flattens these groups.

\texttt{boundaries} is the deterministic all-pair array, and
\texttt{boundaries\_by\_name} provides a name lookup. Each candidate has
\texttt{name}, normalized species identifiers \texttt{first} and \texttt{second},
\texttt{kind} (\texttt{vertical} or \texttt{redox}), \texttt{law}
(\texttt{equilibrium quotient} or \texttt{Nernst}), a descriptive candidate
\texttt{geometry}, and \texttt{provisional=true}. Preparation adds no equation or stability fields; calculation later adds a
separate equation record. Reaction coefficients are added only to a
successfully checked snapshot, inside each boundary's \texttt{thermodynamics}.
Registry retrieval returns defensive copies of this complete prepared snapshot.

\section{Lua rendering and geometry API}
The selected record adds \texttt{final\_boundaries}, \texttt{final\_by\_name}
and \texttt{domains} inside \texttt{selection}. Each final boundary has
\texttt{name}, \texttt{source}, \texttt{kind}, declaration indices
\texttt{first}/\texttt{second}, segment indices, and its equation and
thermodynamic records. Domain records carry species, kind, pH bounds and optional
lower/upper source names.

The pure \texttt{luapourbaix-plot.lua} module provides
\texttt{geometry(diagram, bounds)} with finite \texttt{xmin}, \texttt{xmax},
\texttt{ymin}, \texttt{ymax}. It returns clipped domains, diagnostic bands, lines
and interior label anchors. Its \texttt{render(diagram, options)} produces generic
plot-helper commands. The bridge supplies these to TeX only after successful
completion. Style tokens remain entirely on the TeX side. The pure TikZ renderer
accepts \texttt{render(diagram, "map")} for the standalone map.

Tests check definitive numbering through \texttt{aa}, clipped-line equations,
coverage of the iron viewport by its domain polygons, diagnostic copper-band
limits, zoom clipping, missing prerequisites and options, and map orientation.
All three formats verify actual plot visualization occurs at axis end, scoped
style-token transport, handle capture before name rebinding, overlays, and
coexistence with plots before and after the library command. Visual examples
remain inline and execute during manual compilation.

\section{Lua selection records}
\texttt{core.select(calculated, options)} returns a new selected copy;
\texttt{options} requires \texttt{ph\_min}, \texttt{ph\_max}, and optionally
\texttt{ph\_step}. The registry's \texttt{select(handle, options)} returns a handle
and defensive copy. The diagram's \texttt{selection} contains \texttt{segments},
\texttt{events}, \texttt{inventory}, \texttt{removed}, \texttt{nodes}, and
\texttt{cells}. Each cell has \texttt{ph\_min}, \texttt{ph\_max}, its midpoint
\texttt{sample}, surviving \texttt{species} indices in oxidation order, and
adjacent \texttt{boundaries} source names. Each segment has an \texttt{equation}
record; each event has a \texttt{reaction} coefficient table. Infinite Lua bounds
are signed \texttt{math.huge}; only the TeX readouts turn these into strings.

Analytic tests cover copper disproportionation, successive eliminations,
zero-width intermediates, an off-grid inversion, coarse/fine reconstruction
agreement, iron precipitation limits, acid-only unbounded verticals, boundaries
on the requested range endpoints, tied solids, parameter validation and snapshot
rollback. All three TeX formats test scalar transport, scoped aliases and normal
PGFPlots plotting of the selected expression.

\section{Diagnostics and verification}
Plot option errors suppress the affected diagram request, even if TeX continues
in nonstop mode. A failed water preparation inside a plotting command cannot
fall back to an earlier water record. Existing records remain available for a
later valid request. This guard covers option processing before submission;
arbitrary executable TikZ styles can still fail during PGFPlots visualization.

An invalid declaration does not replace any valid object or consume a successful
handle. Diagnostics cover missing fields, unknown keys, malformed formulas,
invalid numbers, duplicate species or constants, foreign-element species and
references to undeclared species. A missing boundary convention is always an
error. Unknown names and engines without Lua also produce explicit diagnostics.
When TeX continues after a reported error, the invalid declaration remains
uncommitted. A later valid declaration may still succeed.

Declaration validation certifies input structure and the concentration
convention. The explicit check additionally verifies supported reaction balances,
thermodynamic sufficiency and cycle consistency. Neither certifies the physical
plausibility of a species set or the stability of a chemical boundary. Missing
reaction models are reported separately from missing numerical constants.

\texttt{make test} runs the pure Lua core without engine tables, exercises both
concentration conventions and the solid-solid unit-activity rule, checks formula
counts and validation failures, and
compiles shared declaration and classification-pic tests under all three formats. Those tests check
name scoping, independent snapshots, explicit expansion, inert data transport
and transaction failure. Negative compilations check the required convention and
the rejection of pdfTeX. Preparation tests check oxidation overrides, complete
explicit acidity orders, formal ties, all-pair counts, custom labels, alphabetic
numbering beyond \texttt{z}, temperature validation and rejection of reversed
oxidized/reduced roles. Ordinary plots and standalone classification pics are
exercised together without replacing PGFPlots' survey mechanism. Thermodynamic
tests independently check electron-weighted potential composition, hydroxide
conversion, inferred water constants, contradictory cycles, reaction coefficients
for dichromate reduction, logarithmic overflow avoidance, provenance, result-key
scoping and atomic failure under all three formats. Catalogue tests balance
every imported reaction, exercise every entry against its endpoint pair, verify
source metadata and defensive copies, and check default phases, overrides,
user-data precedence, aqueous intermediates and temperature rejection.
An optional \texttt{tests/lua/source-audit.lua} script independently reconstructs
59 USGS reactions and logarithmic constants from the three pinned upstream files
without network access. With only \texttt{phreeqc.dat}, it checks the original
50 entries. The ordinary suite also checks phase-aware catalogue cycles and
the distinct meanings of source potentials and source $\log K$ values.

\texttt{make doc} compiles this manual, generates its index and recompiles it.
\texttt{make check} runs both suites. All artifacts remain in \texttt{build/}.
\texttt{make dist} runs these checks and builds a local CTAN source archive,
checksums and draft submission metadata in \texttt{build/dist/}; it does not
publish anything. \texttt{make manifest} refreshes the distribution file list
after source files are added or removed. \texttt{make clean} removes only the
generated build tree. Release metadata and the local preparation procedure are
documented in \texttt{release.json} and \texttt{doc/releasing.md}.
The examples above are part of this source and execute within the manual;
no example is externalized and no precompiled example PDF is included.

\clearpage\printindex
\end{document}
